\documentclass[11pt,twoside]{article}
\usepackage{amssymb}
\usepackage{amsmath}
\usepackage{mathrsfs}
\usepackage{amsthm}
\usepackage{amsfonts}
\usepackage{txfonts}
% \usepackage{newtxtext}   % 正文文本字体
% \usepackage{newtxmath}   % 如果需要数学字体
\usepackage{enumerate}
%\usepackage[retainorgcmds]{IEEEtrantools}
\usepackage{hyperref}
%\usepackage{esint}
\usepackage{latexsym,amssymb}
\usepackage{color}
\usepackage{slashed}
\usepackage[normalem]{ulem}
\newcommand{\rev}[1]{{\color{red}#1}}
\newcommand{\revdel}[1]{{\color{red}\sout{#1}}}
\newcommand{\revnote}[1]{{\color{red}\footnotesize\textit{#1}}}

\usepackage{indentfirst}
\usepackage{graphicx}
\usepackage{pgf,tikz}
\usepackage{mathrsfs}
\usetikzlibrary{arrows}

\pagestyle{empty}

%\usepackage{refcheck}


%\usepackage[notref,notcite]{showkeys}

%\let\mathbb=\varmathbb
%\DeclareSymbolFont{letters}{OML}{ztmcm}{m}{it}

\allowdisplaybreaks

\pagestyle{myheadings}
\markboth{\footnotesize\rmfamily\scshape Jun Cao, Haonan Lu, Guangzheng Wang and Chenfeng Zhu}
{\footnotesize\rmfamily\scshape 
Revised Density Property for  Triebel--Lizorkin Spaces on Bounded  Domains 
in Full Smoothness Range}

\textwidth=15truecm
\textheight=21.6truecm
\oddsidemargin 0.35cm
\evensidemargin 0.35cm



\makeatletter
\def\@cite#1#2{#1}
\makeatother
\def\rr{{\mathbb R}}
\def\rn{{\mathbb{R}^n}}
\def\urn{\mathbb{R}_+^{n+1}}
\def\zz{{\mathbb Z}}
\def\cc{{\mathbb C}}
\def\aa{{\mathbb A}}
\def\nn{{\mathbb N}}
\def\nno{\nn_0}
\def\hh{{\mathbb H}}
\def\ff{{\mathbb F}}
\def\bbg{{\mathbb G}}
\def\ca{{\mathcal A}}
\def\cb{{\mathcal B}}
\def\ccc{{\mathcal C}}
\def\cd{{\mathcal D}}
\def\ce{{\mathcal E}}
\def\cf{{\mathcal F}}
\def\hf{{\mathscr F }}
\def\cg{{\mathcal G}}
\def\ch{{\mathcal H}}
\def\hhh{{\mathscr H}}
\def\cj{{\mathcal J}}
\def\cci{{\mathcal I}}
\def\cl{{\mathcal L}}
\def\cm {{\mathcal M}}
\def\cn{{\mathcal N}}
\def\cp{{\mathcal P}}
\def\cq{{\mathcal Q}}
\def\mr{{\mathcal R}}
\def\cs{{\mathcal S}}
\def\ct{{\mathcal T}}
\def\cx{{\mathcal X}}
\def\cy{{\mathcal Y}}
\def\mi{{\mathrm I}}
\def\mj{{\mathrm J}}
\def\mh{{\mathrm H}}
\def\mrm{\mathrm}
\def\cw{\mathcal W}
\def\sxing{S^{*}}
\def\qxing{Q^{*}}
\def\pxing{P^{*}}
\def\Qi{\frac{11}{10}Q_i}
\def\Qkk{\frac{6}{5}Q_k}
\def\Qk{\frac{11}{10}Q_k}
\def\Qii{\frac{6}{5}Q_i}
\def\Qj{\frac{11}{10}Q_j}

\def\fz{\infty }
\def\az{\alpha}
\def\bz{\beta}
\def\dz{\delta}
\def\bdz{\Delta}
\def\ez{\epsilon}
\def\gz{{\gamma}}
\def\bgz{{\Gamma}}
\def\kz{\kappa}
\def\lz{\lambda}
\def\blz{\Lambda}
\def\oz{{\omega}}
\def\boz{{\Omega}}
\def\tz{\theta}
\def\sz{\sigma}
\def\vz{\varphi}
\def\vez{\varepsilon}

\def\lf{\left}
\def\r{\right}   
\def\lfz{{\lfloor}}              %中括号
\def\rfz{{\rfloor}}
\def\la{\langle}                 %尖括号
\def\ra{\rangle}
\def\hs{\hspace{0.25cm}}
\def\ls{\lesssim}                %小于约等于
\def\gs{\gtrsim}                 %大于约等于
\def\tr{\Delta}
\def\pa{\partial}                %偏导
\def\ov{\overline}               %上划线（闭包）
\def\noz{\nonumber}
\def\wz{\widetilde}              %波浪线
\def\wh{\widehat}                %尖帽子
\def\ev{\equiv}                  %恒等于
\def\com{\complement}            %补集符号
\def\bh{\backslash}              %差集符号
\def\gfz{\genfrac{}{}{0pt}{}}
\def\wee{\wedge}                 %尖型交
\def\st{\subset}                 %包含于

\def\dis{\displaystyle}
\def\dlimsup{\displaystyle\limsup}
\def\dist{\mathop\mathrm{\,dist\,}}
\def\loc{{\mathop\mathrm{\,loc\,}}}
\def\supp{\mathop\mathrm{\,supp\,}}
\def\atom{\mathop\mathrm{\,atom\,}}

\def\cappqk{\text{cap}_{\cf_{p,q}^{\sz}(\boz)}(K)}
\def\capft{\text{cap}_{\cf_{p,q}^{\sz}(\boz)}(\left\{|f|>t \right\})}
\def\hcondition{(\text{H})_{\sz,p,q}-condition}
\def\qiuhetwo{\sum_{j \in \Lambda_2} \sum_{\substack{\bz\in\nn^n \\ \bz \leq \az}} \binom{\az}{\bz}}
\def\fpjzuhel{\sum_{\bz \leq \az} \binom{\az}{\bz} D^{\bz}f(x) D^{\az - \bz} \psi_j(x)}
\def\fpjzuhely{\sum_{\bz \leq \az} \binom{\az}{\bz} D^{\bz}f(y) D^{\az - \bz} \psi_j(y)}
\def\fpjzuhe{D^{\bz}f(x) \sum_{j \in \Lambda_2}D^{\az - \bz} \psi_j(x)}
\def\fpjzuhey{D^{\bz}f(y) \sum_{j \in \Lambda_2}D^{\az - \bz} \psi_j(y)}
\def\we{w_{\varepsilon}}
\def\axdk{\sum_{\lvert \alpha \rvert \leq k}}
\def\adk{\sum_{\substack{\az\in\nno^n \\ \lvert \alpha \rvert = k}}}
\def\huae{\varepsilon}
\def\ve{v_{\varepsilon}}
\def\bozthreee{\boz_{3\huae}}
\def\bozfour{\Omega_{14\huae}}
\def\fhuaf{{F}_{p,q}^s(\boz)}
\def\fq{\mathring{F}_{p,q}^s(\boz)}
\def\fw{\widetilde{F}_{p,q}^s(\boz)}
\def\huaf{{\cf}_{p,q}^s(\boz)}
\def\huafwe{{\cf}_{p,q}^s(\boz,d^{-aq})}
\def\huafweab{{\hf}_{p,q}^{s}(\boz)}
\def\huafsz{{\cf}_{p,q}^\sz(\boz)}
\def\huafq{\mathring{\cf}_{p,q}^s(\boz)}
\def\huafqw{\mathring{\cf}_{p,q}^s(\boz,d^{-aq})}
\def\huafqwab{\mathring{\hf}_{p,q}^{s}(\boz)}
\def\huafw{\widetilde{\cf}_{p,q}^s(\boz)}
\def\capfk{\sum_{k = -\infty}^{\infty} 2^{kp} \text{cap}_{\huafsz} (\{|f| > 2^k\})}
\def\cro{\ccc_\rho}
\def\mz{{\frac n2(\frac 1{p_\Phi^-}-\frac12)}}
\def\mzx{{\frac n2(\frac 1{\wz p_\Phi}-\frac12)}}
\def\hl{{\mathrm H}_l}
\def\jl{{\mathrm J}_l}
\def\rb{{\rho(\mu(B))[\mu(B)]^{1/2}}}
\def\rkb{{\rho(\mu(2^kB))[\mu(2^kB)]^{1/2}}}
\def\rjb{{\rho(\mu(B_j))[\mu(B_j)]^{1/2}}}
\def\dxt{\,\frac{d\mu(x)\,dt}{t}}
\def\dyt{\,\frac{d\mu(y)\,dt}{t}}
\def\dt{\,\frac{dt}{t}}
\def\ds{\,\frac{ds}{s}}
\def\cub{\chi_{U_k(B)}}
\def\lp{{L^p(\cx)}}
\def\lx{{L^\Phi(\cx)}}
\def\xt{{\cx\times(0,\fz)}}
\def\ml{{M ^{M,\ez}_\Phi(L)}}
\def\mlx{{M ^{M,\ez}_\Phi(L^*)}}
\def\mly{{M ^{M_1}_{\Phi,L}(\cx)}}
\def\tx{{T_\Phi(\cx)}}
\def\txx{{(T_\Phi(\cx))^*}}
\def\txz{{T_\Phi^\fz(\cx)}}
\def\txv{{T_{\Phi,{\rm v}}^\fz(\cx)}}
\def\txl{{T_{\Phi,0}^\fz(\cx)}}
\def\txy{{T_{\Phi,1}^\fz(\cx)}}
\def\txb{{T_{2,b}^2(\cx)}}
\def\txe{{T_{2}^2(\cx)}}
\def\ttx{{\wz T_\Phi(\cx)}}
\def\dfb{\dz_k(f,B)}
\def\tml{(t^2L)^Me^{-t^2L}}
\def\tmlx{(t^2L^*)^Me^{-t^2L^*}}
\def\tmy{(t^2L)^{M_1}e^{-t^2L}}
\def\iem{(I-e^{-r_B^2L})^M}
\def\ilm{(I-[I+r_B^2L]^{-1})^M}
\def\Iem{\lf(I-e^{-r_B^2L}\r)^M}
\def\Ilm{\lf(I-[I+r_B^2L]^{-1}\r)^M}
\def\iemx{(I-e^{-r_B^2L^{\ast}})^M}
\def\Iemx{\lf(I-e^{-r_B^2L^{\ast}}\r)^M}
\def\tl{t^2Le^{-t^2L}}
\def\tlx{t^2L^*e^{-t^2L^*}}
\def\vy{{V_{k,1}(B)}}
\def\pme{(\Phi,\,M,\,\ez)_L}
\def\pml{(\Phi,\,\wz M,\,\ez)_{L^*}}
\def\pmx{(\Phi,\,M,\,\ez)_{L^*}}
\def\bmol{{\mathop\mathrm{\,BMO}^M_{\rho,L}(\cx)}}
\def\bmolx{{\mathop\mathrm{\,BMO}^M_{\rho,L^*}(\cx)}}
\def\vmol{{\mathop\mathrm{\,VMO}^M_{\rho,L}(\cx)}}
\def\tvmol{{\mathop\mathrm{\wz{\,VMO}}^M_{\rho,L}(\cx)}}
\def\vmo{{\mathop\mathrm{\,VMO}_{\rho,L}(\cx)}}
\def\bmo{{\mathop\mathrm{\,BMO}_{\rho,L}(\cx)}}
\def\bmox{{\mathop\mathrm{\,BMO}_{\rho,L^*}(\cx)}}
\def\bmoxx{{(\bmox)^*}}
\def\vmox{{(\vmo)^*}}
\def\hx{{H_{\Phi,L}(\cx)}}
\def\hxx{{H_{\Phi,L^*}(\cx)}}
\def\hmfl{{H_{\Phi,\mathrm {fin},L}^{\mathrm {mol},\ez, M}(\cx)}}
\def\hmflx{{H_{\Phi,\mathrm {fin},L^*}^{\mathrm {mol},\ez, M}(\cx)}}
\def\hmx{{H_{\Phi,L}^{M,\ez}(\cx)}}
\def\bmx{{B_{\Phi,L}^{M,\ez}(\cx)}}
\def\byx{{B_{\Phi,L}^{M_1,\ez_1}(\cx)}}
\def\bx{{B_{\Phi,L}(\cx)}}
\def\bxx{{B_{\Phi,L^*}(\cx)}}
\def\plm{\pi_{L,M}}
\def\ply{\pi_{L,1}}
\def\vlp{{L^{p(\cdot)}(\rn)}}
\def\vhs{H^{p(\cdot)}(\rn)}
\def\cps{\cl_{q,p(\cdot),s}(\rn)}
\def\cpss{\cl_{1,p(\cdot),s}(\rn)}
\def\tlve{f_{p(\cdot),q(\cdot)}^{w(\cdot),\phi}(\rn)}
\def\btlve{F_{p(\cdot),q(\cdot)}^{w(\cdot),\phi}(\rn)}
\def\beve{b_{p(\cdot),q(\cdot)}^{w(\cdot),\phi}(\rn)}
\def\bbeve{\dot B_{p,q}^{\az}(\rn)}
\def\mbeve{{\mathcal{N}}_{p,q,\lz}^s(\rn)}
\def\dmbeve{{FH}_{(p,\lz),q}^s(\rn)}
\def\lpqbl{{\mathcal{H}}^{(p,\lz),q}_{\bz}(\rr^{2n})}
\def\Lz{\Lambda}
\def\d{\mathcal{D}}
\def\s{\mathcal{S}}
\newcommand{\Fdot}{\dot{\mathfrak F}}

\newcounter{assum}
\renewcommand{\theassum}{\Roman{assum}}
\newtheorem{assumption}[assum]{Assumption}
\newtheorem{theorem}{Theorem}[section]
\usepackage{amsthm}

\newtheoremstyle{mylem}
{}{}{\itshape}{}{\bfseries}{.}{ }
{\thmname{#1}\thmnumber{ #2}\thmnote{\normalfont([#3])}}

\theoremstyle{mylem}
\newtheorem{lemma}[theorem]{Lemma}

\newtheorem{corollary}[theorem]{Corollary}
\newtheorem{proposition}[theorem]{Proposition}
\newtheorem{example}[theorem]{Example}
\theoremstyle{definition}
\newtheorem{remark}[theorem]{Remark}
\newtheorem{definition}[theorem]{Definition}
\renewcommand{\appendix}{\par
\setcounter{section}{0}%
\setcounter{subsection}{0}%
\setcounter{subsubsection}{0}%
\gdef\thesection{\@Alph\c@section}%
\gdef\thesubsection{\@Alph\c@section.\@arabic\c@subsection}%
\gdef\theHsection{\@Alph\c@section.}%
\gdef\theHsubsection{\@Alph\c@section.\@arabic\c@subsection}%
\csname appendixmore\endcsname
}

\newtheorem{notation}[theorem]{Notation}

\numberwithin{equation}{section}

\definecolor{qqqqff}{rgb}{0,0,1}
\definecolor{ffqqqq}{rgb}{1,0,0}



\begin{document}
\arraycolsep=1pt
\title{\bf Revised Density Property for  Triebel--Lizorkin Spaces on Bounded  Domains 
in Full Smoothness Range
\footnotetext{\hspace{-0.35cm} 2020 {\it Mathematics Subject Classification}.
Primary 46E35; Secondary 42B35, 26A33.
\endgraf {\it Key words and phrases.}  Triebel--Lizorkin space, 
Density property, Distance weight, Uniform domain, Hardy inequality.}}
\author{Jun Cao, Haonan Lu, Guangzheng Wang and Chenfeng Zhu\footnote{Corresponding author,
E-mail: \texttt{chenfengzhu@zjut.edu.cn}/{\color{red}\today}/Final version.}}
\date{}
\maketitle

\vspace{-0.8cm}

\begin{center}
\begin{minipage}{13cm}\small
{\noindent{\bf Abstract.}
Let $\Omega\subset\mathbb{R}^n$ be  a bounded domain and
$d(x)=d(x,\partial\Omega)$ be the boundary distance function.
Suppose that $s\in(0,\infty)$, $p,q\in[1,\infty)$ and $a\in(0,\infty)$.
This paper investigates the problem of compactly supported smooth functions approximation in 
weighted Triebel--Lizorkin space $\mathcal{F}_{p,q}^{s}(\Omega,d^{-aq})$, whose intrinsic norm is  
augmented by suitable Hardy-type terms. Since weighted 
Meyers--Serrin approximation is not available in general,
we formulate such problem to a revised version by comparing the closures of 
$C_{\mathrm{c}}^\infty(\Omega)\cap\mathcal{F}_{p,q}^{s}(\Omega,d^{-aq})$
and $C^\infty(\Omega)\cap\mathcal{F}_{p,q}^{s}(\Omega,d^{-aq})$ within $\mathcal{F}_{p,q}^{s}(\Omega,d^{-aq})$.
This formulation reduces to the classical density property in the unweighted low-order case (i.e. $a=0$ and $s\in(0,1)$). 
The main results prove that the revised density property holds in the full smoothness range $s\in(0,\infty)$.
A new feature in the higher-order smoothness
case $s\in [1,\infty)$ is that the distance weight structure arises naturally from Leibniz's rule under Whitney cutoffs. 
Finally, in the range $s\in(0,1)$, 
we show that the density property is stable when the intrinsic norm is replaced by a restriction norm,
provided the domain supports a suitable Hardy inequality.
}
\end{minipage}
\end{center}

%%%%%%%%%%%%%%%%%%%%%%%%%%%%%%%%%%%%%%%%%%%%%%%%%%%%%%%%%%%%%%%%%%%%%%%%%
%%%%%%%%%%%%%%%%%%%%%%%%%%%  section 1  %%%%%%%%%%%%%%%%%%%%%%%%%
%%%%%%%%%%%%%%%%%%%%%%%%%%%%%%%%%%%%%%%%%%%%%%%%%%%%%%%%%%%%%%%%%%%%%%%%%

\vspace{0.2cm}

\tableofcontents

\section{Introduction}

Let $\Omega$ be a domain in the Euclidean space $\rn$
with $\pa\boz$ being its boundary.  Denote by $d=d(\cdot,\partial\Omega):=\inf_{y\in\pa\boz}|\cdot-y|$
the \emph{boundary distance function}.
For any $s\in(0,\infty)$, 
$p,q\in[1,\fz)$ and $a\in\rr$,  define the \emph{weighted Triebel--Lizorkin space}
$\mathcal{F}_{p, q}^{s}(\boz,d^{-aq})$ by setting  
\begin{align}\label{eqn-TLs1}
\mathcal{F}_{p, q}^{s}(\boz,d^{-aq}):=\left\{f\in L^p(\boz): \ \|f\|_{\mathcal{F}_{p, q}^{s}(\boz,d^{-aq})}<\fz\right\},
\end{align}
where, in the case $s=k+\sigma\in (0,\infty)\setminus \nn$ with $k\in\nn_0$ and $\sigma\in (0,1)$, the
\emph{fractional-order Triebel--Lizorkin norm} (see, e.g.,  [\cite{p2019}] for the case $a=0$) is defined by 
\begin{align}\label{eqn-1.3}
&\|f\|_{\mathcal F_{p,q}^s(\Omega,d^{-aq})}\nonumber\\
&\quad:=\sum_{\substack{\az\in {\nno^n} \\ |\alpha| \leq k}}\left\| D^{\alpha}f \right\|_{L^p(\Omega)}
+\sum_{\substack{\beta\in {\nno^n} \\ |\beta|=k }}\left\{\int_{\Omega} \left[\int_{\Omega} 
\frac{|D^\beta f(x)-D^\beta f(y)|^q}
{|x - y|^{n+\sigma q}d(y,\pa\boz)^{aq}} \,dy \right]^{\frac{p}{q}}\, dx \right\}^{\frac{1}{p}} \nonumber\\
&\quad\,=:\| f \|_{W^{k,p}(\Omega)} 
+\sum_{\substack{\beta\in {\nno^n} \\ |\beta|=k }}\left|D^\beta f\right|_{\mathcal{F}_{p,q}^{\sz}(\Omega,d^{-aq})}.
\end{align}
In the integer-case $s=k\in\nn$, the \emph{integer-order Triebel--Lizorkin norm} (see, e.g.,  [\cite{t2010}, p.\,101, (2)] for the case $\boz=\rn$ and $a=0$) is defined by 
\begin{align}\label{1614}
&\|f\|_{\mathcal F_{p,q}^s(\boz,d^{-aq})}\nonumber\\
&\quad:=\left\| f \right\|_{L^p(\Omega)}+\left\{
\int_{\boz}
\left[
\int_{|h|<\frac{d(x,\pa\boz)}{2(k+1)} }
\frac{\left|\Delta_h^{k+1}f(x)\right|^q}
{|h|^{n+kq}d(x,\partial \Omega)^{aq}}
\,d h
\right]^\frac{p}{q}
\,d x
\right\}^\frac{1}{p}.
\end{align}
Here, for any $m\in\nn$,  $x\in \Omega$ and $h\in \rn$  satisfying $m|h|<d(x,\pa\boz)$, 
the \emph{$m$-order difference} of $f$ at $x$ is defined by
\begin{align*}
\Delta_h^m f(x)
:=\sum_{\nu=0}^m (-1)^{m-\nu}\binom{m}{\nu}f(x+\nu h).
\end{align*}
Note that the norm 
$\|f\|_{\mathcal F_{p,q}^s(\boz,d^{-aq})}$ is equivalent if we replace
the weight $d(x,\partial \Omega)$ by $d(x+h,\partial \Omega)$ in \eqref{1614}, due to the fact 
$|h|<\frac{d(x,\pa\boz)}{2(k+1)}$.
The range of \(h\) ensures that all points $x+\nu h$, $0\leq\nu\leq k+1$, remain in $\Omega$, it is permissible to suitably enlarge the upper bound of \(h\), provided that all the above points remain in \(\Omega\).
Moreover, in the case  $a=0$, we simply write $\mathcal F_{p,q}^s(\Omega)$ 
instead of $\mathcal F_{p,q}^s(\Omega,1)$. In the case
$a>0$, we have the inclusion
$$\mathcal{F}_{p, q}^{s}(\boz,d^{-aq})\subset\mathcal F_{p,q}^s(\Omega)
\subset \mathcal{F}_{p, q}^{s}(\boz,d^{aq}),$$ whenever $\boz$ is bounded.


In this paper, we investigate the density property of the Triebel--Lizorkin
spaces $\mathcal{F}_{p,q}^{s}(\boz,d^{-aq})$. To be precise, 
denote by $C_{\mathrm{c}}^{\fz}(\boz)$ the set of all smooth functions with compact support in 
$\boz$.
Let 
\begin{align}\label{eqn-1.4}
\mathcal F_{p,q}^{s,1}(\boz,d^{-aq}):=\overline{C_{\mathrm{c}}^{\fz}(\boz)\cap \mathcal F_{p, q}^{s}(\boz,d^{-aq})}^{\|\cdot\|_{\mathcal F_{p,q}^s(\boz,d^{-aq})}}
\end{align}
endowed with the norm
$\|\cdot\|_{\mathcal F_{p,q}^s(\Omega,d^{-aq})}$ as in \eqref{eqn-TLs1}. 
If $a=0$,  we simply write ${\mathcal F}_{p,q}^{s,1}(\boz)$. Motivated by the research in [\cite{c2000}, \cite{cjlz2022}, \cite{dk2021}, \cite{dk2022}, \cite{k2023}] on the density problem in function spaces, we naturally consider the following question.

\medskip 

{\noindent{\bf Question I}. When does the identity $\mathcal F_{p,q}^{s,1}(\boz,d^{-aq})=\mathcal F_{p,q}^{s}(\boz,d^{-aq})$ hold?}

\medskip
This question plays a fundamental role in the analysis of $\mathcal F_{p, q}^{s}(\boz,d^{-aq})$, since 
it determines whether the class of  ``simple" compactly supported smooth functions can approximate
every function in this space.  This problem is highly complicated and relies heavily on the weight exponent $a$, 
smoothness parameter $s$ and domain geometry. 
Below we provide a partial overview of existing research on this subject.
If $a=0$ (unweighted case),  Question {\bf I} was answered in the affirmative  by Caetano [\cite{c2000}], 
for $s\in (0,{(n-\mathrm{dim}_M(\partial\Omega))}/{p})$ in terms of the upper Minkowski dimension $\mathrm{dim}_M(\partial\Omega)$
of $\partial \Omega$ (see [\cite{klv2013}] for the definition of $\mathrm{dim}_M(\partial\Omega)$).


\medskip
{\noindent{\bf Theorem A} ([\cite{c2000}, Proposition 2.2]).   Let \(\Omega\subset\mathbb{R}^n\) be a bounded domain 
with  its boundary upper Minkowski dimension $\mathrm{dim}_M(\partial\Omega)<n$. Then for any $ p,q\in[1,\infty)$ and $s\in (0,
{(n-\mathrm{dim}_M(\partial\Omega))}/{p})$,
\begin{align*}
\mathcal F_{p,q}^{s,1}(\boz)={\mathcal F}_{p,q}^{s}(\boz).
\end{align*}
}

Note that the exponent ${(n-\mathrm{dim}_M(\partial\Omega))}/{p}$ in Theorem {\bf A}
is sharp since there are counterexamples showing that 
the conclusion of Theorem {\bf A} may fail 
for any $s>{(n-\mathrm{dim}_M(\partial\Omega))}/{p}$. This indicates that to deal with the 
case $s>{(n-\mathrm{dim}_M(\partial\Omega))}/{p}$, we
need more conditions.
Among them, one of the most commonly used conditions is the following Hardy condition: 
Let $s\in(0,\infty)$, $p,q,r\in[1,\infty)$ and $a,b\in\rr$.
We say a function $f$ satisfies the 
\emph{Hardy condition} $(H)_{s,p,q,a,r,b}$
if $f\in \mathcal F_{p, q}^{s}(\boz,d^{-aq})$ and
\begin{align*}
\int_{\boz}  \frac{|f(x)|^r}{d (x, \pa\boz)^{b}}\,dx < \infty.
\end{align*}
Note that the Hardy-type condition $(H)_{s,p,q,a,r,b}$ imposes a specific decay bound for functions 
$f\in \mathcal{F}_{p,q}^{s}(\Omega,d^{-aq})$ as $x$ approaches the boundary $\partial\Omega$.
If $\Omega$ is bounded, then it is easy to see $(H)_{s,p,q,a,r,b_2}\Rightarrow (H)_{s,p,q,a,r,b_1}$
for any $0<b_1<b_2<\infty$.
Recall that such condition was already used in [\cite{dk2022}] to study the density property of weighted 
fractional Sobolev space $\mathcal W^{s,p}(\Omega,d^{-ap})$ in the 
case $s\in (0,1)$ and $p\in [1,\infty)$. 
Recall that $\mathcal W^{s,p}(\Omega,d^{-ap})=\mathcal F_{p, p}^{s}(\boz,d^{-ap})$ and for 
any $f\in \mathcal W^{s,p}(\Omega,d^{-ap})$, its 
\emph{weighted Gagliardo norm} is defined by setting
\begin{align*}
\|f\|_{\mathcal W^{s,p}(\Omega,d^{-ap})}:=\| f \|_{L^{p}(\Omega)} 
+\left[ \int_{\Omega} \int_{\Omega} 
\frac{\left| f(x)- f(y)\right|^p}
{|x - y|^{n+s p}d(y,\pa\boz)^{a p}}\,dy\,dx \right]^{\frac{1}{p}}.
\end{align*}
The space $\mathring{\mathcal{W}}^{s,p}(\boz,d^{-ap})$ is
defined similarly to $\mathcal F_{p,q}^{s,1}(\boz,d^{-aq})$ by replacing $\mathcal F_{p, q}^{s}
(\boz,d^{-aq})$ with $\mathcal{W}^{s,p}(\boz,d^{-ap})$ in \eqref{eqn-1.4}. To study the relation between
$\mathring{\mathcal{W}}^{s,p}(\boz,d^{-ap})$
and $\mathcal{W}^{s,p}(\boz,d^{-ap})$,
we first recall the definition of uniform domains (see, e.g., [\cite{go1979}, \cite{v1988}]).
A domain \(\boz\subset\rn\) is called a \emph{uniform domain} if there exist constants 
\(c_1,c_2>0\) such that any two points \(x,y\in\boz\) can be connected by a rectifiable curve \(\bgz\subset\boz\) satisfying that
\begin{align*}
\begin{cases}
L(\bgz)\leq c_1 |x-y|, \ \text{where} \ L(\bgz)\ \text{denotes the Euclidean arc length of}\ \bgz;  \\
\min\{|x-z|,|y-z|\}\leq c_2 d(z,\pa\boz)\text{ for any } z\in\bgz.
\end{cases}
\end{align*}
The uniform domain is a global condition, whereas the \((\vez,\delta)\)-domain is a local condition (see \cite{j1981} for the definition of \((\vez,\delta)\)-domains). A uniform domain is necessarily an \((\vez,\delta)\)-domain for some \(\delta\), while a general \((\vez,\delta)\)-domain need not be a uniform domain. When the domain is bounded, a uniform domain is equivalent to some \((\vez,\delta)\)-domain.
We refer to \cite{j1981} for more details about the relation of
above domains.

Kijaczko [\cite{k2023}] proved the following result
(see [\cite{klv2013}] for the definitions of upper and lower Assouad co-dimensions $\overline{\mathrm{codim}}_A(\partial\Omega)$
and $\underline{\mathrm{codim}}_A(\partial\Omega)$.
In particular, it is known that
$\underline{\mathrm{codim}}_A(\partial\Omega)=n-\mathrm{dim}_M(\partial\Omega)$ as in 
Theorem {\bf A}, when $\partial \Omega$ is $n-\underline{\mathrm{codim}}_A(\partial\Omega)$-homogeneous, where a set is called $s$-homogeneous when its lower and upper Assouad dimensions both equal $s$).

\medskip
{\noindent{\bf Theorem B} ([\cite{k2023}, Theorem 4]).  Let $\Omega\subset\mathbb{R}^n$ be a 
bounded uniform domain with $\overline{\mathrm{codim}}_A(\partial\Omega)$ and $\underline{\mathrm{codim}}_A(\partial\Omega)$
being its boundary upper and lower Assouad co-dimensions.
Let $s\in(0,1)$,  $p\in[1,\infty)$ and $a\in[0, \underline{\mathrm{codim}}_A(\partial\Omega)/p)$. 
If $s> \overline{\mathrm{codim}}_A(\partial\Omega)/p-a$,
then
\begin{align*}
\mathring{\mathcal{W}}^{s,p}(\boz,d^{-ap}) = \left\{ f\in \mathcal{W}^{s,p}(\boz,d^{-ap}):
\int_\Omega \frac{|f(x)|^p}{d(x,\pa\boz)^{(s+a)p}}\,dx < \infty \right\}.
\end{align*}
}

For the case where $p\ne q\in [1,\infty)$, $s\in (0,1)$ and $a=0$, the following result was proven in [\cite{cjlz2022}].

\medskip
{\noindent{\bf Theorem C} ([\cite{cjlz2022}, Proposition 1]).   Let \(\Omega\subset\mathbb{R}^n\) be a bounded uniform domain, 
$s\in(0, 1)$ and $p,q\in[1,\infty)$. If $f\in\mathcal F_{p,q}^s(\Omega)$ satisfies the Hardy condition $(H)_{s,p,q,0,p,sp}$, then
\begin{align*}
\|f\|_{\mathcal F_{p,q}^{s,1}(\boz)}\sim\|f\|_{{\mathcal F}_{p,q}^{s}(\boz)},
\end{align*}
where the positive equivalence constants are independent of $f$.
}

Motivated by the aforementioned results, we aim to investigate Question {\bf I} in both the weighted case $a>0$ 
and the higher-order setting $s\ge 1$. As it turns out, both the weight function $d^{-ap}$ and higher-order regularity 
introduce additional difficulties to this problem. Let us first come to the weighted case 
$a>0$. In this case, 
one \emph{difficulty} is that the weighted Meyers--Serrin theorem (i.e. $C^\infty(\Omega)$
is dense in ${\mathcal F}_{p,q}^{s}(\boz,d^{-aq})$) is known to hold only under very strict condition on the parameter $a$,
even when $p=q\in [1,\infty)$ and $s\in (0,1)$ (see, e.g., [\cite{dk2021}, Theorem 12]). To circumvent it, 
we introduce another space 
\begin{align}\label{eqn-1.xx4}
\mathcal F_{p,q}^{s,2}(\boz,d^{-aq}):=\overline{C^\infty(\boz)\cap \mathcal F_{p, q}^{s}(\boz,d^{-aq})}^{\|\cdot\|_{\mathcal F_{p,q}^s(\boz,d^{-aq})}}
\end{align}
as the completion of $C^\infty(\boz)\cap \mathcal F_{p, q}^{s}(\boz,d^{-aq})$
with respect to the norm $\|\cdot\|_{\mathcal F_{p,q}^s(\boz,d^{-aq})}$ define as in 
\eqref{eqn-1.3}.
Note that if the weighted Meyers--Serrin theorem holds, then we immediately have the identity
$${\mathcal F}_{p,q}^{s}(\boz,d^{-aq})=\mathcal F_{p,q}^{s,2}(\boz,d^{-aq}).$$
Also, we know that such identity always holds in the unweighted case $a=0$.
Moreover, based on the Hardy condition in Theorem {\bf B}
and Theorem {\bf C}, we naturally introduce the following
two weighted Triebel--Lizorkin norms with Hardy term
\begin{align}\label{159}
\|f\|_{\widetilde{\mathcal{F}}_{p,q}^s(\Omega,d^{-aq})}
:=\|f\|_{\mathcal F_{p,q}^s(\Omega,d^{-aq})}
+\left\|\frac{f}{d^{s+a}}\right\|_{L^p(\Omega)}
+\left\|\frac{f}{d^{s+a+\frac{n}{q}}}\right\|_{L^q(\Omega)}
\end{align}
and
\begin{align*}
\|f\|_{\widehat{\mathcal{F}}_{p,q}^s(\Omega,d^{-aq})}
:=\|f\|_{\mathcal F_{p,q}^s(\Omega,d^{-aq})}
+\left\|\frac{f}{d^{s+a}}\right\|_{L^p(\Omega)}.
\end{align*}
We then define
\begin{align*}
\widetilde{\mathcal{F}}_{p,q}^s(\Omega,d^{-aq})
:=\left\{f\in L^p(\boz):\|f\|_{\widetilde{\mathcal{F}}_{p,q}^s(\Omega,d^{-aq})}<\infty\right\}
\end{align*}
and
\begin{align*}
\widehat{\mathcal{F}}_{p,q}^s(\Omega,d^{-aq})
:=\left\{f\in L^p(\boz):\|f\|_{\widehat{\mathcal{F}}_{p,q}^s(\Omega,d^{-aq})}<\infty\right\}.
\end{align*}
The spaces $\widetilde{\mathcal{F}}_{p,q}^{s,1}(\boz,d^{-aq})$ and $\widehat{\mathcal{F}}_{p,q}^{s,1}(\boz,d^{-aq})$
are defined similarly to that in \eqref{eqn-1.4},
and the spaces
$\widetilde{\mathcal{F}}_{p,q}^{s,2}(\boz,d^{-aq})$ and $\widehat{\mathcal{F}}_{p,q}^{s,2}(\boz,d^{-aq})$ are defined similarly to that in \eqref{eqn-1.xx4}.

\begin{remark}\label{remark}
Let $\Omega\subset\mathbb{R}^n$, $s\in(0,1)$,
$p,q\in[1,\infty)$ and $a\in(0,\infty)$.
\begin{enumerate}
\item[\rm(i)] 
If $\Omega$ supports the following Hardy
inequality: for any $f\in\mathcal F_{p,q}^s(\Omega,d^{-aq})$,
\begin{align*}
\left\|\frac{f}{d^{s+a}}\right\|_{L^p(\Omega)}
+\left\|\frac{f}{d^{s+a+\frac{n}{q}}}\right\|_{L^q(\Omega)}
\lesssim\|f\|_{\mathcal F_{p,q}^s(\Omega,d^{-aq})},
\end{align*}
then $\widetilde{\mathcal{F}}_{p,q}^s(\Omega,d^{-aq})=\widehat{\mathcal F}_{p,q}^{s}(\Omega,d^{-aq})
=\mathcal{F}_{p,q}^s(\Omega,d^{-aq})$
with equivalent norms.
\item[\rm(ii)]
If $\Omega$ supports the following Hardy
inequality: for any $f\in\mathcal F_{p,q}^s(\Omega,d^{-aq})$,
\begin{align*}
\left\|\frac{f}{d^{s+a}}\right\|_{L^p(\Omega)}
\lesssim\|f\|_{\mathcal F_{p,q}^s(\Omega,d^{-aq})},
\end{align*}
then $\widehat{\mathcal{F}}_{p,q}^s(\Omega,d^{-aq})
=\mathcal{F}_{p,q}^s(\Omega,d^{-aq})$ with equivalent norms.
\end{enumerate}
\end{remark}

Based on the aforementioned argument, if a suitable weighted Meyers--Serrin theorem
were available, we could address Question {\bf I}. However, in the absence
of such a theorem, we must modify Question {\bf I}. 
We shall reformulate
Question {\bf I} in three different smoothness regimes (see
Question {\bf I-1} for \(s\in(0,1)\), Question {\bf I-2} for
\(s\in(1,\infty)\setminus\mathbb N\), and Question {\bf I-3} for
\(s\in\mathbb N\)). 
These three modifications are essential for the approach adopted in this
paper. Moreover, we believe that without such modifications, 
the parameter range to which the original Question {\bf I} applies would be
quite limited.
We first reformulate Question {\bf I} for \(s\in(0,1)\) as
the following question.

\medskip 

{\noindent{\bf Question I-1}. 
When do the identities
\[
\widetilde{\mathcal{F}}_{p,q}^{s,1}(\boz,d^{-aq}) =\widetilde{\mathcal{F}}_{p,q}^{s,2}(\boz,d^{-aq})
\]
or
\[
\widehat{\mathcal{F}}_{p,q}^{s,1}(\boz,d^{-aq})= \widehat{\mathcal{F}}_{p,q}^{s,2}(\boz,d^{-aq})
\]
hold?}	


\medskip

The following Theorem \ref{thm-main} is our first main result, which 
provides an affirmative answer to Question {\bf I-1} in the case $s\in (0,1)$ (see [\cite{a2015}, \cite{pt2015}, 
\cite{s1970}] for the definition of Lipschitz domains).

\begin{theorem}\label{thm-main} 
Let \(\Omega\subset\mathbb{R}^n\) be a bounded domain, 
$s\in (0, 1)$, $p,q\in[1,\infty)$ and $a\in(0,\infty)$.
\begin{enumerate}
\item[\rm(i)]
Then $\widetilde{\mathcal{F}}_{p,q}^{s,1}(\boz,d^{-aq}) =\widetilde{\mathcal{F}}_{p,q}^{s,2}(\boz,d^{-aq})$.
\item[\rm(ii)]
Assume further that \(\Omega\subset\mathbb{R}^n\) is a Lipschitz domain and
$aq<1$. Then $$\widehat{\mathcal{F}}_{p,q}^{s,1}(\boz,d^{-aq})= \widehat{\mathcal{F}}_{p,q}^{s,2}(\boz,d^{-aq}).$$
\end{enumerate}
\end{theorem}


The Hardy term $\|{f}/{d^{s+a+\frac{n}{q}}}\|_{L^q(\Omega)}$
in \eqref{159} appears rather restrictive yet unavoidable within our approach. 
Indeed, for any $\delta\in (0,\infty)$, let
\begin{align}\label{eqn-OBD}
\boz_\delta:=\left\{x\in\boz: d(x,\partial \boz)<\dz\right\}.
\end{align}
Given $\varepsilon >0$ sufficiently small, in many scenarios of the proof of Theorem \ref{thm-main}, 
we need the Hardy term $\|{f}/{d^{s+a+\frac{n}{q}}}\|_{L^q(\Omega)}$ to bound the following term
\begin{align*}
\mathrm{I}:=\int_{\Omega \setminus \Omega_{14\varepsilon}} \left[  \int_{\Omega_{3\varepsilon}} 
\frac{| f(y)|^q}{|x - y|^{s q + n}d(y,\pa\boz)^{aq}}\,dy \right]^{\frac{p}{q}}\,dx.
\end{align*}
Note that in the unweighted case $a=0$, the term $\mathrm{I}$ can be bounded via the Hardy term $\|{f}/{d^{s}}\|_{L^q(\Omega)}$ as in Theorem {\bf C} via polar coordinates:
\begin{align*}
\mathrm{I}&\lesssim \int_{\Omega \setminus \Omega_{14\varepsilon}} \left[ \int_{\Omega_{3\varepsilon}} 
\frac{| f(y)-f(x)|^q}{|x - y|^{s q + n} }\,dy \right]^{\frac{p}{q}}\,dx\\
&\quad+\int_{\Omega \setminus \Omega_{14\varepsilon}} \left[  \int_{\Omega_{3\varepsilon}} 
\frac{| f(x)|^q}{|x - y|^{s q + n} }\,dy \right]^{\frac{p}{q}}\,dx \\
&\lesssim \|f\|_{\mathcal{F}_{p,q}^s(\boz)}^p+\int_{\Omega \setminus \Omega_{14\varepsilon}} 
\frac{|f(x)|^p}{d(x,\partial \boz)^{sp}}\,dx.
\end{align*}
Moreover, if $aq<1$, we can get rid of the Hardy term $\|{f}/{d^{s+a+\frac{n}{q}}}\|_{L^q(\Omega)}$
by virtue of the summation property of Whitney cubes (see the proof of Theorem \ref{thm-main} for more details). 
In particular, it is proven in [\cite{dilt2019}, Corollary 3.8] that if  $\partial\Omega\subset\rn$ is a closed set with 
upper Assouad dimension 
$\overline{\mathrm{dim}}_A(\Omega)=n-1$,  then 
the boundary  distance $d^{-aq}$ is an $A_1$ weight if and only if $aq\in(0,1)$.

Next, we come to the higher-order smoothness setting $s\in (1,\infty)\setminus\nn$. 
Here, Leibniz's rule for higher-order derivatives creates extra 
difficulties. Indeed, in the density problem, we always need to approximate 
a given function $f$ by the product of $f$ and a suitably chosen bump function $\psi$ over a Whitney cube
(see, for instance, \eqref{eqn-4.23} in the proof of 
Theorem \ref{thm-main}). Since $s>1$, Leibniz's rule 
gives rise to intermediate derivative terms of the form $D^\beta fD^{\alpha-\beta}\psi$ with $\beta<\alpha$. 
By virtue of the scaling properties of bump functions, the  term
$D^{\alpha-\beta}\psi$ induce a distance weight $d^{-|\alpha-\beta|}$,
which only disappears in the case $s\in (0,1)$.

The above argument points to the following three definitions of modified weighted intrinsic Triebel--Lizorkin spaces,
the first of which when $k=0$ reduces back to the space $\mathcal F_{p,q}^s(\Omega)$.

\begin{definition}\label{def-secmain}
Let \( k \in \nn_0 \), \(\sz\in(0,1) \) and \( p, q\in[1, \infty) \). Let \( \boz \subset \rn\) be a domain and
$s := k + \sz\in (0,\infty)\setminus \nn$. The \emph{modified weighted intrinsic Triebel--Lizorkin spaces}
$\hf_{p,q}^{s}(\boz)$,
$\widetilde{\mathscr{F}}_{p,q}^{s}(\boz)$ 
and $\widehat{\mathscr{F}}_{p,q}^{s}(\boz)$
are defined as
\begin{align*}
\hf_{p,q}^{s}(\boz):=\left\{ f \in L^p(\boz) : \| f \|_{\hf_{p,q}^{s}(\boz)} < \infty \right\},
\end{align*}
\begin{align*}
\widetilde{\mathscr{F}}_{p,q}^{s}(\boz) := \left\{ f \in L^p(\boz) : \| f \|_{\widetilde{\mathscr{F}}_{p,q}^{s}(\boz)} < \infty \right\}
\end{align*}
and
\begin{align*}
\widehat{\mathscr{F}}_{p,q}^{s}(\boz):= \left\{ f \in L^p(\boz) : \| f \|_{\widehat{\mathscr{F}}_{p,q}^{s}(\boz)} < \infty \right\},
\end{align*}
where
\begin{align}\label{eqn-GWN}
\| f \|_{\hf_{p,q}^{s}(\boz)}
:=\| f \|_{\cf_{p,q}^{s}(\boz)} + \sum_{\substack{\bz\in\nno^n \\ | \bz |\leq k}} 
\left|D^{\bz}f\right|_{\cf_{p,q}^{\sz}(\boz,d^{-(k-|\bz|)q})},
\end{align}
\begin{align*}
\| f \|_{\widetilde{\mathscr{F}}_{p,q}^{s}(\boz)}
:=\| f \|_{\hf_{p,q}^{s}(\boz)}+ \sum_{\substack{\bz\in\nno^n \\ |\bz| \leq k}}
\left(\left\|\frac{D^\beta f}{d^{s-|\beta|}}\right\|_{L^p(\Omega)}+\left\|\frac{D^\beta f}{d^{s-|\beta|+\frac{n}{q}}}\right\|_{L^q(\Omega)}\right)
\end{align*}
and
\begin{align*}
\| f \|_{\widehat{\mathscr{F}}_{p,q}^{s}(\boz)}
:=\| f \|_{\hf_{p,q}^{s}(\boz)}+\sum_{\substack{\bz\in\nno^n \\ |\bz| \leq k}}
\left\|\frac{D^\beta f}{d^{s-|\beta|}}\right\|_{L^p(\Omega)}+\sum_{\substack{\bz\in\nno^n \\|\bz| < k}} \left\|\frac{D^\beta f}{d^{s-|\beta|+\frac{n}{q}}}\right\|_{L^q(\Omega)}.
\end{align*}
\end{definition}
Similarly, the {\it completion spaces} $\widetilde{\hf}_{p,q}^{s,1}(\boz)$, $\widetilde{\hf}_{p,q}^{s,2}(\boz)$,
$\widehat{\hf}_{p,q}^{s,1}(\boz)$ and
$\widehat{\hf}_{p,q}^{s,2}(\boz)$
are 
defined, respectively, as 
\begin{align*}
\widetilde{\hf}_{p,q}^{s,1}(\boz):=\overline{C_{\mathrm{c}}^{\fz}(\boz)\cap {\widetilde{\hf}}_{p, q}^{s}(\boz)}^{\|\cdot\|_{{\widetilde{\hf}}_{p,q}^s(\boz)}}, \ \ \widetilde{\hf}_{p,q}^{s,2}(\boz):=\overline{C^\infty(\boz)\cap {\widetilde{\hf}}_{p, q}^{s}(\boz)}^{\|\cdot\|_{{\widetilde{\hf}}_{p,q}^s(\boz)}},
\end{align*}
\begin{align*}
\widehat{\hf}_{p,q}^{s,1}(\boz):=\overline{C_{\mathrm{c}}^{\fz}(\boz)\cap {\widehat{\hf}}_{p, q}^{s}(\boz)}^{\|\cdot\|_{{\widehat{\hf}}_{p,q}^s(\boz)}}
\ \ \text{and}\ \ 
\widehat{\hf}_{p,q}^{s,2}(\boz):=\overline{C^\infty(\boz)\cap {\widehat{\hf}}_{p, q}^{s}(\boz)}^{\|\cdot\|_{{\widehat{\hf}}_{p,q}^s(\boz)}}.
\end{align*}

\begin{remark}
Let \( \boz \subset \rn\) be a domain, \( k \in \nn_0 \), \(\sz\in(0,1) \), \( p, q\in[1, \infty) \)  and	$s := k + \sz$.
\begin{enumerate}
\item[\rm(i)] 
Recall that, if $s\in (0,1)$ and $\Omega$
is a bounded uniform domain, then, by Definition \ref{def-secmain} and  the  extension property of  ${\cf}_{p,q}^{s,2}(\boz)$ (see, e.g., [\cite{ps2017}]), we know 
\begin{align*}
\hf_{p,q}^{s}(\boz)=\mathcal{F}_{p,q}^{s}(\boz)={\hf}_{p,q}^{s,2}(\boz).
\end{align*}

\item[\rm(ii)] 
If $\Omega$ supports the following Hardy
inequality: for any $f\in\hf_{p,q}^{s}(\boz)$,
\begin{align*}
\sum_{\substack{\bz\in\nno^n \\ |\bz| \leq k}}
\left(\left\|\frac{D^\beta f}{d^{s-|\beta|}}\right\|_{L^p(\Omega)}+\left\|\frac{D^\beta f}{d^{s-|\beta|+\frac{n}{q}}}\right\|_{L^q(\Omega)}\right)
\lesssim\|f\|_{\hf_{p,q}^{s}(\boz)},
\end{align*}
then $\widetilde{\mathscr{F}}_{p,q}^{s}(\boz)
=\hf_{p,q}^{s}(\boz)$
with equivalent norms.

\item[\rm(iii)]
If $\Omega$ supports the following Hardy
inequality: for any $f\in\hf_{p,q}^{s}(\boz)$,
\begin{align*}
\sum_{\substack{\bz\in\nno^n \\ |\bz| \leq k}}
\left\|\frac{D^\beta f}{d^{s-|\beta|}}\right\|_{L^p(\Omega)}+\sum_{\substack{\bz\in\nno^n \\|\bz| < k}} \left\|\frac{D^\beta f}{d^{s-|\beta|+\frac{n}{q}}}\right\|_{L^q(\Omega)}
\lesssim\|f\|_{\hf_{p,q}^{s}(\boz)},
\end{align*}
then $\widehat{\mathscr{F}}_{p,q}^{s}(\boz)
=\hf_{p,q}^{s}(\boz)$ with equivalent norms.
\end{enumerate}
\end{remark}

We now refine  Question {\bf I} in the case $s\in(1,\infty)\setminus\mathbb{N}$ to the following question.
\medskip 

{\noindent{\bf Question I-2}.
When do the identities $\widetilde{\hf}_{p,q}^{s,1}(\boz) =\widetilde{\hf}_{p,q}^{s,2}(\boz)$
or $\widehat{\hf}_{p,q}^{s,1}(\boz)
=\widehat{\hf}_{p,q}^{s,2}(\boz)$ hold?}	

\medskip 

Now, we state the second main result of this paper as follows, which provides an answer to 
Question {\bf I-2} in an affirmative way.

\begin{theorem}\label{thm-secmain}
Let $k \in \nn$, $p, q \in [1, \infty)$ and $\sigma\in(0,1)$. Suppose that $\boz$ is a bounded domain and let $s := k + \sz$.
\begin{enumerate}
\item[\rm(i)]
Then $\widetilde{\hf}_{p,q}^{s,1}(\boz) =\widetilde{\hf}_{p,q}^{s,2}(\boz)$.

\item[\rm(ii)]
Assume further that \(\Omega\subset\mathbb{R}^n\) is a Lipschitz domain.
Then $\widehat{\hf}_{p,q}^{s,1}(\boz)=
\widehat{\hf}_{p,q}^{s,2}(\boz)$.
\end{enumerate} 
\end{theorem}

Now, we come to the integer-order smoothness setting. We define the weighted local difference seminorms 
and introduce the modified weighted  Triebel--Lizorkin space based on differences as follows.

\begin{definition}
Let $\Omega\subset\mathbb{R}^n$ be a domain,
$k\in\mathbb{N}$ and $p,q\in[1,\infty)$.
\begin{enumerate}
\item[\rm(i)] 
For any given measurable function $f$ on $\Omega$ and for any $r\in\{0,1,\dots,k\}$,
define its \emph{weighted local difference seminorm} by
\begin{align*}
[f]_{r,k}
:=
\left\{
\int_{\boz}
\left[
\int_{|h|<\frac{d(x,\pa\boz)}{2(k+1)}}
\frac{\left|\Delta_h^{k+1-r}f(x)\right|^q}
{|h|^{n+(k-r)q}d(x,\pa\boz)^{rq}}
\,d h
\right]^\frac{p}{q}
\,d x
\right\}^\frac{1}{p}.
\end{align*}
\item[\rm(ii)]
The \emph{modified weighted Triebel--Lizorkin space based on differences} $\mathfrak F^{k}_{p,q}(\boz)$ is defined to be the set of all measurable functions $f$ on $\boz$ such that
\begin{align}\label{diffnorm}
\|f\|_{\mathfrak F^{k}_{p,q}(\boz)}
:=
\|f\|_{L^p(\boz)}
+
\sum_{r=0}^k [f]_{r,k}
+
\left\|\frac{f}{d^k}\right\|_{L^p(\boz)}
\end{align}
is finite.
\item[\rm(iii)]
Define the completion spaces 
$\mathfrak F^{k,1}_{p,q}(\boz)$
and	$\mathfrak F^{k,2}_{p,q}(\boz)$	as
\begin{align*}
\mathfrak F^{k,1}_{p,q}(\boz)
:=
\overline{
C_{\mathrm{c}}^{\fz}(\boz)\cap \mathfrak F^{k}_{p,q}(\boz)
}^{\|\cdot\|_{\mathfrak F^{k}_{p,q}(\boz)}}
\end{align*}
and
\begin{align*}
\mathfrak F^{k,2}_{p,q}(\boz)
:=
\overline{
C^\infty(\boz)\cap \mathfrak F^{k}_{p,q}(\boz)
}^{\|\cdot\|_{\mathfrak F^{k}_{p,q}(\boz)}}.
\end{align*}
\end{enumerate}
\end{definition}

We now refine Question {\bf I} in the case $s\in\mathbb{N}$ to the following question.
\medskip 

{\noindent{\bf Question I-3}.
When does the identity $\mathfrak F^{k,1}_{p,q}(\boz)=\mathfrak F^{k,2}_{p,q}(\boz)$ hold?}

\medskip 

Now, we state the third main result of this paper as follows, which provides an answer to 
Question {\bf I-3} in an affirmative way.

\begin{theorem}\label{thm-third}
Let $\boz\subset\rn$ be a bounded domain, $k\in\nn$ and $p,q\in[1,\infty)$.  Then
\begin{align*}
\mathfrak F^{k,1}_{p,q}(\boz)=\mathfrak F^{k,2}_{p,q}(\boz).
\end{align*}
\end{theorem}

\begin{remark}
The above three main results, namely
Theorems~\ref{thm-main},~\ref{thm-secmain} and~\ref{thm-third}, concern
only distance weight $d(y,\pa\boz)$. If one also allows
distance weight $d(x,\pa\boz)$, the situation is as
follows. For \(s\in(0,1)\), a result analogous to
Theorem~\ref{thm-main} can be obtained simply by modifying the indices
in the Hardy term. 
For \(s\in(1,\infty)\setminus\mathbb N\), the distance weight arises from the Leibniz rule, and adding distance weight $d(x,\pa\boz)$ would affect this derivation, so the weights $d(y,\pa\boz)$ and $d(x,\pa\boz)$ cannot be considered simultaneously.
Hence such weights are not considered in Theorem~\ref{thm-secmain} here. 
For \(s\in\mathbb N\), the
restriction on the inner integral domain implies that the two distance
weights \(d(x,\pa\boz)\) and \(d(x+h,\pa\boz)\) are comparable, and hence
the corresponding norms are equivalent; therefore only one weight needs
to be considered in Theorem~\ref{thm-third}.
\end{remark}

In the remainder of the introduction, we study a question closely related to 
Question {\bf I}: whether the density property still holds when 
the intrinsic norm in \eqref{eqn-1.3} is replaced by another  restriction norm. 
%let  $s\in\rr$, $p\in(0,\infty)$, $q\in(0,\infty]$ and 
%$\varphi=\{\varphi_j\}_{j=0}^{\infty}$ be a family of unity of resolution.
%Recall that the \emph{Triebel--Lizorkin space} $F_{p,q}^s(\rn)$ on the whole Euclidean space
%$\rn$ is defined by setting
%\begin{align*}
%F_{p,q}^s(\rn):=\left\{f\in\mathcal{S}'(\rn): \|f\|_{F_{p,q}^s(\rn)}
%:=\left\|\left[\sum_{j=0}^{\infty}\left|2^{sj}\left(\varphi_j 
%\hat{f}\right)^\vee\right|^q\right]^{\frac{1}{q}}\right\|_{L^p(\rn)}<\infty \right\},
%\end{align*} 
%where $\hat{f}$ and $\check{f}$ denotes respectively the 
%Fourier and inverse Fourier transform of $f$. 
We study this question only in the case $s\in (0,1)$. To be precise, 
let $\Omega\subset \rn$ be a domain. We introduce the following definitions of 
weighted restriction and completion Triebel--Lizorkin spaces $F_{p, q}^{s}(\boz,d^{-aq})$ and 
${F}_{p,q}^{s,1}(\boz,d^{-aq})$, which reduce back 
to the corresponding Triebel--Lizorkin spaces  $F_{p, q}^{s}(\boz)$ and ${F}_{p,q}^{s,1}(\boz)$ 
from [\cite{cjlz2022,t2001,t2008}] when  $a=0$.


\begin{definition}
Let $\Omega\subset\rn$ be a domain, $s\in(0,1)$,  $p,q\in[1,\fz)$ and $a\in\rr$.
\begin{enumerate}
\item[\rm(i)]The \emph{weighted restriction Triebel--Lizorkin space} $F_{p, q}^{s}(\boz,d^{-aq})$ is defined by setting
$$
F_{p, q}^{s}(\boz,d^{-aq}):=\left\{f\in L^p(\boz): \exists\ g\in F_{p,q}^s(\rn,d^{-aq}) \text{ such that }g=f\text{ on }\Omega\right\}
$$ endowed with the norm
\begin{align}\label{eqn-wrctls}
\|f\|_{F_{p, q}^{s}(\boz,d^{-aq})}:=\inf\|g\|_{F_{p,q}^s(\rn,d^{-aq})},
\end{align}
where the infimum is taken over all $g\in F_{p,q}^s(\rn,d^{-aq})$ as above and
\begin{align*}
\|g\|_{F_{p,q}^s(\rn,d^{-aq})}
&:=\|g\|_{L^p(\rn)}+ \left\{ \int_{\rn} \left[\int_{\rn} \frac{|g(x) -g(y)|^q}
{|x - y|^{s q + n}}\,dy \right]^{\frac{p}{q}}\,dx \right\}^{\frac{1}{p}} \nonumber\\ 
&\hs\hs+\left\{ \int_{\rn} \left[\int_{\rn} \frac{|g(x)-g(y)|^q}{|x - y|^{s q + n}d(y,\pa\boz)^{aq}}\,dy \right]^{\frac{p}{q}}\,dx \right\}^{\frac{1}{p}}.
\end{align*}

\item[\rm(ii)]
The \emph{weighted restriction Triebel--Lizorkin space with Hardy term} $\widetilde{F}_{p, q}^{s}(\boz,d^{-aq})$ is defined by setting
\begin{align*}
\widetilde{F}_{p, q}^{s}(\boz,d^{-aq}):=\left\{
f\in L^p(\boz):
\|f\|_{\widetilde{F}_{p, q}^{s}(\boz,d^{-aq})}:=\|f\|_{F_{p, q}^{s}(\boz,d^{-aq})}
+\left\|\frac{f}{d^{s+a+\frac{n}{q}}}\right\|_{L^q(\Omega)}<\infty\right\}.
\end{align*}

\item[\rm(iii)]
The \emph{weighted completion Triebel--Lizorkin spaces} 
$F_{p,q}^{s,1}(\boz,d^{-aq})$, $F_{p,q}^{s,2}(\boz,d^{-aq})$,
$\widetilde{F}_{p,q}^{s,1}(\boz,d^{-aq})$ and
$\widetilde{F}_{p,q}^{s,2}(\boz,d^{-aq})$ are defined, respectively, by setting $$F_{p,q}^{s,1}
(\boz,d^{-aq}):=\overline{C_{\mathrm{c}}^{\fz}(\boz)\cap F_{p, q}^{s}(\boz,d^{-aq})}^{\|\cdot\|_{F_{p,q}^s(\boz,d^{-aq})}},$$ 
$$F_{p,q}^{s,2}(\boz,d^{-aq}):=\overline{C^\infty(\boz)\cap F_{p, q}^{s}(\boz,d^{-aq})}^{\|\cdot\|_{F_{p,q}^s(\boz,d^{-aq})}},$$
$$\widetilde{F}_{p,q}^{s,1}
(\boz,d^{-aq}):=\overline{C_{\mathrm{c}}^{\fz}(\boz)\cap\widetilde{F}_{p, q}^{s}(\boz,d^{-aq})}^{\|\cdot\|_{\widetilde{F}_{p,q}^s(\boz,d^{-aq})}}$$
and 
$$\widetilde{F}_{p,q}^{s,2}(\boz,d^{-aq}):=\overline{C^\infty(\boz)\cap\widetilde{F}_{p, q}^{s}(\boz,d^{-aq})}^{\|\cdot\|_{\widetilde{F}_{p,q}^s(\boz,d^{-aq})}}.$$
\end{enumerate}
\end{definition}

\begin{remark}
Let $a=0$, $p,q\in[1,\infty)$ and $s\in(\max\{0,\frac{n}{p} - \frac{n}{q}\},1)$. It is proved in [\cite{ps2017}, Theorem 1.2] that $\|\cdot\|_{F_{p,q}^s(\rn,1)}$ is equivalent to the classical Triebel--Lizorkin norm (see [\cite{t2010}, Definition 2.3.1] for the definition of classical Triebel--Lizorkin norms).
\end{remark}
We also need the following version of Hardy's inequality
(see, e.g.,  [\cite{dk2024}, \cite{k2023},\cite{iv2013}, \cite{iv2015}] for some related interesting results on this topic).

\begin{assumption}\label{1617}
Let $a\in(0,\infty)$, $ p, q\in[1,\infty)$, $s\in(0,1)$ and $\boz \subsetneqq \rn$ be a domain.
We assume that $\boz$ supports {Hardy's inequality}
$(HI)_{s,p,q,a}$ in the following sense: for any $f \in \mathcal F_{p, q}^{s}(\boz,d^{-aq})$ defined as in \eqref{eqn-TLs1},
\begin{equation}\label{eqn-Hardy}
\int_{\boz}  \frac{|f(x)|^p}{d(x, \pa\boz)^{(s+a)p}} dx \lesssim  \int_{\boz} \left[\int_{\boz} \frac{|f(x) - f(y)|^q}
{|x - y|^{s q + n}d(y,\pa\boz)^{aq}}\,dy \right]^{\frac{p}{q}}\,dx.
\end{equation}
\end{assumption}



Using Assumption \ref{1617}, the following result provides an affirmative answer to  Question {\bf I-1}
with  the intrinsic norm in \eqref{eqn-1.3} replaced by the restriction norm in \eqref{eqn-wrctls}.


\begin{theorem}\label{co-main} 
Let $a\in(0,\infty)$, $p,q \in [1, \infty)$ and $s\in(\max\{0,\frac{n}{p} - \frac{n}{q}\},1)$. 
Assume that \(\Omega\subset\mathbb{R}^n\) is a bounded uniform domain that  supports Hardy's inequality
$(HI)_{s,p,q,a}$. Then the following statements are true.
\begin{enumerate}
\item[\rm(i)]$F_{p,q}^{s}(\boz,d^{-aq})=\mathcal F_{p,q}^{s}(\boz,d^{-aq})$ with equivalent norms.
\item[\rm(ii)]
$\widetilde{F}_{p,q}^{s,1}
(\boz,d^{-aq})
=\widetilde{F}_{p,q}^{s,2}
(\boz,d^{-aq})
=\widetilde{\mathcal{F}}_{p,q}^{s,1}(\boz,d^{-aq}) =\widetilde{\mathcal{F}}_{p,q}^{s,2}(\boz,d^{-aq})$
with equivalent norms.
\item[\rm(iii)]
Assume further that $aq<1$ and $\Omega$ is a Lipschitz domain.
Then 
\begin{align*}
F_{p,q}^{s,1}(\boz,d^{-aq})
&= F_{p,q}^{s,2}(\boz,d^{-aq})
=\widehat{\mathcal{F}}_{p,q}^{s,1}(\boz,d^{-aq})
=\widehat{\mathcal{F}}_{p,q}^{s,2}(\boz,d^{-aq})\\
&=\mathcal{F}_{p,q}^{s,1}(\boz,d^{-aq})= \mathcal{F}_{p,q}^{s,2}(\boz,d^{-aq})
\end{align*}
with equivalent norms.
\end{enumerate}
\end{theorem}

The most crucial part of Theorem \ref{co-main} is (i). To prove it, we need the following boundedness of linear extension operator, which 
when $a=0$ was proven by Prats and Saksman in [\cite{ps2017}, Lemma 3.10]. The extension problem in Triebel--Lizorkin spaces has been extensively studied 
(see, e.g., [\cite{ayy2022}, \cite{hit2016}, \cite{zyy2024}]) .

\begin{proposition}\label{pro-main}
Let $a\in(0,\infty)$, $p,q \in [1, \infty)$ and $s\in(\max\{0,\frac{n}{p} - \frac{n}{q}\},1)$. Assume that $\boz$ is a bounded uniform domain
that  supports Hardy's inequality $(HI)_{s,p,q,a}$.
Then there exists a bounded linear extension operator $\Lambda_0\colon \huafwe \to F_{p,q}^{s}(\rn,d^{-aq})$.
\end{proposition}

The rest of this paper is organized as follows.
Section~\ref{s2} collects some preliminaries on the Whitney decomposition of a domain,
which will be used throughout the proofs.
Section~\ref{s3} is devoted to the density property of the weighted Triebel--Lizorkin space
in the case $s\in(0,1)$.
To be precise, we prove Theorem~\ref{thm-main} in this section.
In Section~\ref{s4} we show the density property for Triebel--Lizorkin spaces on $\Omega$
in the higher-order smoothness case $s\in(1,\infty)\setminus\mathbb{N}$
by proving Theorem~\ref{thm-secmain}.
Section~\ref{sN} deals with the integer-order smoothness case;
we establish Theorem~\ref{thm-third} there.
Finally, in Section~\ref{s5} we prove that the density property is stable
if the intrinsic norm is replaced by the restriction norm defined in \eqref{eqn-wrctls}.
This is achieved by first proving the bounded linear extension operator
in Proposition~\ref{pro-main} and then using it to obtain Theorem~\ref{co-main}.

We end this introduction by making some convention on the notation.

\noindent {\bf Notation}. Throughout the paper, let $\nn:=\left\{1,2,...\right\}$ and $\nn_0:=\nn\cup\{0\}$. 
We use $C$ to denote a positive constant that is independent of the main parameters involved; its value may differ from 
line to line, and we do not distinguish between different constants $C$ appearing in different steps.
The symbol $A\lesssim B$ means that there exists a universal constant $C$ such that $A\leq CB$, while $A\sim B$ 
means $A\lesssim B$ and $B\lesssim A$. For any cube $Q\subset\mathbb{R}^n$, let $\mathrm{diam\,}Q:=
\sup_{x,y\in Q}|x-y|$ and $l(Q)$ be its sidelength. By a domain, we mean a connected open set. 
For any set $E\subset \rn$, denote by $E^\complement $ the set $\rn\setminus E$, by $\overline{E}$ the 
closure of $E$, by $|E|$ the Lebesgue measure of $E$ and by $\#E$ the cardinality of $E$. Moreover, for any 
$x\in\rn$, define $d(x,E):=\inf_{y\in E}|x-y|$ and $d(E,F):=\inf_{y\in E}d(y,F)$.
We also use $\chi_E$ to denote the characteristic function of $E$. When $|E|\in(0,\infty)$, for any 
$f\in L_{\mathrm{loc}}^1(\rn)$ (the set of all locally integrable functions on $\rn$), we denote by 
$f_E:=\frac{1}{|E|}\int_E f(x)\,dx$ the integral mean of $f$ over $E$.


%%%%%%%%%%%%%%%%%%%%%%%%%%%%%%%%%%%%%%%%%%%%%%%%%%%%%%%%%%%%%%%%%%%%%

%%%%%%%%%%%%%%%%%%%%%%%  section 2  %%%%%%%%%%%%%%%%%%%%%%%%%%%%%%%%%

%%%%%%%%%%%%%%%%%%%%%%%%%%%%%%%%%%%%%%%%%%%%%%%%%%%%%%%%%%%%%%%%%%%%%
\section{Preliminaries}\label{s2}

\hspace{1em} In this section, we collect some basic properties of Whitney decomposition of domain $\boz$, with emphasis on those Whitney cubes that are close to the boundary. These properties play an important role in the proofs of our main results.

\begin{lemma}[\cite{g2008}, \cite{s1970}]\label{lem2.1}
Let \( \boz \subsetneq \rn \) be a domain. There exists a family of cubes \( \{Q_j\}_{j=1}^{\infty} \) with sides parallel to the coordinate axes and satisfying
\begin{enumerate}
\item[\rm(i)]\( Q_j^o \cap Q_k^o = \emptyset \) if \( j \neq k \), where \( Q_j^o \) denotes the interior of \( Q_j \);
\item[\rm(ii)]for any \( j \in \nn \), \( \mathrm{diam} \, Q_j \leq d(Q_j, \pa\boz) \leq 4 \, \mathrm{diam} \, Q_j \);
\item[\rm(iii)]
\( \boz = \bigcup_{j=1}^{\infty}Q_j \);
\item[\rm(iv)]
each \( x \in \boz \) is contained in at most \( 12^n \) cubes \( \frac{11}{10}Q_j \), where, for any \(j\in\nn\), \( \frac{11}{10}Q_j \) is the concentric cube of \( Q_j \) with sidelength \( \frac{11}{10} l_j \) and $l_j$ denotes the side length of $Q_j$;
\item[\rm(v)]if \( Q_i \) and \( Q_j \) touch, namely, \( \overline{Q_i} \cap \overline{Q_j} \neq \emptyset \) and \( Q_i^o \cap Q_j^o = \emptyset \), then
\[
\frac{1}{4} \, \mathrm{diam} \, Q_i \leq \textup{diam} \, Q_j \leq 4 \, \mathrm{diam} \, Q_i.
\]
\end{enumerate}
\end{lemma}
Let \( \{Q_j\}_{j=1}^{\infty} \) be the Whitney decomposition of \( \Omega \) as in Lemma \ref{lem2.1}. It is known that there exists a family of bump functions \( \{\psi_j\}_{j=1}^{\infty} \subset C_{\mathrm{c}}^\infty(\mathbb{R}^n) \) satisfying the following properties.
\begin{enumerate}
\item[\rm(i)]
For any \(j\in\nn\), \( \psi_j \equiv 1 \) on \( \frac{10}{11}Q_j \), \(0\leq\ \psi_j\leq1\) on \(\rn\) and \( \text{supp}\,\psi_j \subset\frac{11}{10}Q_j \).
\item[\rm(ii)]For any \( x \in \Omega \), it holds that
\begin{equation}\label{eqn-2.1}
\sum_{j=1}^{\infty} \psi_j(x) = 1.
\end{equation}
\item[\rm(iii)]There exists a positive constant \( C \) such that for any \( x \in \mathbb{R}^n \), \(j\in\nn\) and \( k \in \mathbb{N} \),
\begin{equation}\label{eqn-2.2}
|D^{k}\psi_j(x)| \leq \frac{C}{(\mathrm{diam}\,Q_j)^{k}}. 
\end{equation}
\end{enumerate}
Let \( \varepsilon > 0 \). The following classes of index sets represent three subgroups of \( \{Q_j\}_{j=1}^{\infty} \) that are closely related to the boundary layer in \( \boz \). Let
\begin{align}\label{eqn-2.3}
\Lambda_1 := \left\{ j \in \mathbb{N} : d(Q_j, \partial\Omega) \geq \varepsilon \right\}, \ \ \Lambda_2 := \left\{ j \in \mathbb{N} : d(Q_j, \partial\Omega) < \varepsilon \right\}
\end{align}
and
\begin{align}\label{eqn-2.4}
\Lambda_3 &:= \left\{ j \in \mathbb{N} : Q_j \cap (\Omega \setminus \Omega_{14\varepsilon}) \neq \emptyset \right\}
\end{align} 
with \( \boz_{14\varepsilon} \) as in \eqref{eqn-OBD}. Note that $\Lambda_1\cup\Lambda_2=\mathbb{N}$ and $\Lambda_1\cap\Lambda_2=\emptyset$.

\medskip
The following lemmas collect some basic properties  from [\cite{cjlz2022}] on the above three index sets.

\begin{lemma}[\cite{cjlz2022}]\label{lem2.2}
Let \( \varepsilon > 0 \) and \( \blz_1 \), \( \blz_2 \), \(\Lambda_3\) be the index sets as in \eqref{eqn-2.3} and \eqref{eqn-2.4}. Then the following three assertions hold.
\begin{enumerate}
\item[\rm(i)]There exists a positive constant $C$ such that
\begin{equation}\label{eqn-2.5}
\bigcup_{j \in \blz_1} \frac{11}{10}Q_j \subseteq \left\{ x \in \boz : d(x, \pa\boz) \geq C\varepsilon \right\}.
\end{equation}
\item[\rm(ii)]\begin{equation}\label{eqn-2.6}
\bigcup_{j \in \blz_2} \frac{6}{5}Q_j \subseteq \boz_{3\varepsilon} 
\end{equation}
with \( \boz_{3\varepsilon} \) as in \eqref{eqn-OBD}.
\item[\rm(iii)]For any \( k \in \blz_2 \), \( j \in \blz_3 \), \( x \in \Qkk \) and \( y \in Q_j \), it holds that
\begin{align*}
|x - y| \sim D(Q_j, Q_k),
\end{align*}
where 
$$D(Q_j, Q_k) := d(Q_j, Q_k) + l_j + l_k$$
and the positive equivalence constants are independent of \( \varepsilon \), \( j \), \( k \), \( x \) and \( y \).
\end{enumerate}
\end{lemma}

We recall the following useful estimate from \cite{pt2015}.

\begin{lemma}[\cite{pt2015}, Lemma 3.13]\label{lem2.3}
Let \( \boz \) be a bounded Lipschitz domain and \( \{Q_j\}_{j=1}^{\infty} \) its Whitney decomposition. Let $b>a>n-1$. Then there exists a positive constant \( C \) such that for any $j_0\in\nn$,
\[
\sum_{j=1}^{\infty} \frac{l(Q_j)^a}{D(Q_j, Q_{j_0})^{b}} \leq Cl(Q_{j_0})^{a-b}.
\]
\end{lemma}	

We end this section by giving properties of two subgroups of Whitney cubes from \( \blz_2 \) as in \eqref{eqn-2.3}. 
To this end, for any fixed \( i \in \blz_2 \), we make a subdivision of \( \blz_2 \) by setting
\begin{equation}
\label{eqn-2.8}
\begin{aligned}
\blz_{2,1}(i) &:= \left\{ k \in \blz_2 : \Qkk \cap \Qi \neq \emptyset \right\}, \\
\text{and}\ \ \blz_{2,2}(i) &:= \left\{ k \in \blz_2 : \Qkk \cap \Qi = \emptyset \right\}.
\end{aligned}
\end{equation}
For any \( i \in \blz_2 \) and \( k \in \blz_{2,1}(i) \), let
\begin{equation}\label{eqn-2.9}
\blz_{2,3}(i,k) := \left\{ j \in \blz_2 : \Qj \cap \Qi \neq \emptyset \text{ or } \Qj \cap \Qkk \neq \emptyset \right\}.
\end{equation}

\begin{lemma}[\cite{cjlz2022}]\label{lem2.4}
Let \( \boz \) be a bounded domain, \( \varepsilon > 0 \) and \( \blz_2 \) be as in \eqref{eqn-2.3}. Then the following two assertions hold.
\begin{enumerate}
\item[\rm(i)] For any \( i \in \blz_2 \), let \( \blz_{2,1}(i) \) be the index set as in \eqref{eqn-2.8}. Then there exists a constant $c$ such that for any \( x \in \Qi \),
\[
\bigcup_{k \in \blz_{2,1}(i)} \Qkk \subseteq B(x, cl_i).
\]
\item[\rm(ii)]For any \( i \in \blz_2 \) and \( k \in \blz_{2,1}(i) \), let \( \blz_{2,3}(i,k) \) be the index set as in \eqref{eqn-2.9}.
Then there exists \( N \in \nn \), independent of \( i \) and \( k \), such that
\begin{align*}
\#\blz_{2,3}(i,k) \leq N,
\end{align*}
moreover, for any \( j \in \blz_{2,3}(i,k) \),
\begin{align*}
l_i \sim l_j 
\end{align*}
with positive equivalence constants independent of \( i \) and \( j \).
\end{enumerate}
\end{lemma}	

%%%%%%%%%%%%%%%%%%%%%%%%%%%%%%%%%%%%%%%%%%%%%%%%%%%%%%%%%%%%%%%%%%%%%

%%%%%%%%%%%%%%%%%%%%%%%  section 3  %%%%%%%%%%%%%%%%%%%%%%%%%%%%%%%%%

%%%%%%%%%%%%%%%%%%%%%%%%%%%%%%%%%%%%%%%%%%%%%%%%%%%%%%%%%%%%%%%%%%%%%

\section{The case $s\in(0,1)$}\label{s3}

This section is devoted to the density property of the weighted 
Triebel--Lizorkin space in the case $s\in (0,1)$. To be precise, 
we prove Theorem \ref{thm-main} in this section.

\begin{proof}[Proof of Theorem \ref{thm-main}]
We first prove (i). 
It is obvious to see that $\widetilde{\mathcal{F}}_{p,q}^{s,1}(\boz,d^{-aq}) \subset\widetilde{\mathcal{F}}_{p,q}^{s,2}(\boz,d^{-aq})$
since $C_{\mathrm{c}}^\infty(\boz)\subset C^\infty(\boz)$. Thus, we only need to prove $\widetilde{\mathcal{F}}_{p,q}^{s,2}(\boz,d^{-aq})\subset\widetilde{\mathcal{F}}_{p,q}^{s,1}(\boz,d^{-aq})$. To this end, it suffices to show 
\begin{align}\label{eqn-thm-mainc}
C^\infty(\boz)\cap \widetilde{\mathcal F}_{p, q}^{s}(\boz,d^{-aq})
\subset\widetilde{\mathcal{F}}_{p,q}^{s,1}(\boz,d^{-aq}).
\end{align}
Let \( f \in C^\infty(\boz)\cap\widetilde{\mathcal F}_{p, q}^{s}(\boz,d^{-aq}) \), and let \( \{Q_j\}_{j=1}^{\infty} \) and \( \{\psi_j\}_{j=1}^{\infty} \subset C_{\mathrm{c}}^\infty(\mathbb{R}^n) \) be, respectively, the Whitney decomposition and its associated partition of unity fixed in Section~\ref{s2}.
For any \( \varepsilon > 0 \) and \( x \in \Omega \), using the index sets $\Lambda_1$ and 
$\Lambda_2$ in \eqref{eqn-2.3}, we divide $f$ by    
\begin{align}\label{eqn-4.1}
f(x)= \sum_{j \in \Lambda_1} \psi_j(x) f(x) + \sum_{j \in \Lambda_2} \psi_j(x) f(x) 
=: v_{\varepsilon}(x) + w_{\varepsilon}(x).
\end{align}
For \( v_{\varepsilon} \), let \( \psi := \sum_{j \in \Lambda_1} \psi_j \). From the fact that $\psi_j\in C_{\mathrm{c}}^\infty(\mathbb{R}^n) $ and \eqref{eqn-2.5}, it follows that \( \psi \in C_{\mathrm{c}}^{\fz}(\boz) \). Then, by the fact \( f \in C^\infty(\boz) \) and \eqref{eqn-4.1}, we find that 
$v_{\varepsilon}\in C_{\mathrm{c}}^{\fz}(\boz)$.  Thus, to prove \eqref{eqn-thm-mainc}, it suffices 
to show the following identity
\begin{equation}\label{eqn-4.2}
\lim_{\huae \to 0} \| \we \|_{\widetilde{\mathcal F}_{p, q}^{s}(\boz,d^{-aq}) }=0.
\end{equation}
Indeed, if \eqref{eqn-4.2} holds, then it is easy to see 
$v_{\varepsilon}\in C_{\mathrm{c}}^{\fz}(\boz)\cap\widetilde{\mathcal F}_{p, q}^{s}(\boz,d^{-aq}) $. This, together with \eqref{eqn-4.1} and \eqref{eqn-4.2}, further implies $f\in\widetilde{\mathcal{F}}_{p,q}^{s,1}(\boz,d^{-aq})$ and hence \eqref{eqn-thm-mainc} holds. 

We now prove \eqref{eqn-4.2}. To this end, we divide the proof into several 
steps.

{\bf Step 1}. By \eqref{eqn-4.1}, \eqref{eqn-2.6} and \eqref{eqn-2.1}, we obtain
\begin{align*}
\int_{\boz}| \we(x)|^p\,dx\leq\int_{{\boz}_{3\varepsilon}}\left|f(x)\sum_{j =1}^\infty \psi_j(x)\right|^p\,dx=\int_{{\boz}_{3\varepsilon}}| f(x)|^p\,dx,
\end{align*}
which tends to $0$ as $\varepsilon\to 0$ and hence 
$$
\lim_{\huae \to 0} \| \we \|_{L^p(\Omega)}=0.
$$
Similarly, we have
$$
\lim_{\huae \to 0}\left(\left\|\frac{\we}{d^{s+a}}\right\|_{L^p(\Omega)}
+\left\|\frac{\we}{d^{s+a+\frac{n}{q}}}\right\|_{L^q(\Omega)}\right)=0.
$$
Thus, to prove \eqref{eqn-4.2}, it
remains to show that
\begin{equation}\label{eqn-4.3}
\lim_{\huae \to 0} | \we |_{\huafwe}=0.
\end{equation}
By \eqref{eqn-4.1} and the fact that \( \text{supp}\,\psi_j \subset \frac{11}{10}Q_j \) for any $j\in\Lambda_2$, we write
\begin{align}\label{eqn-4.4}
&|\we|_{\huafwe}^p \nonumber\\
&\quad=\int_{\boz} \left[\int_{\boz} \frac{| f(x)\sum_{j \in \Lambda_2} \psi_j(x) -  f(y)\sum_{j \in \Lambda_2} \psi_j(y)|^q}{|x - y|^{s q + n}d(y,\pa\boz)^{aq}}\,dy \right]^{\frac{p}{q}}\,dx  \nonumber \\
&\quad=\left[ \int_{\bigcup_{i \in \blz_2} \Qi}+ \int_{\bozfour \setminus \bigcup_{i \in \blz_2} \Qi}+ \int_{\boz \setminus \bozfour} \right]\left(\int_{\boz}\dots\,dy\right)^{\frac{p}{q}}\,dx \nonumber \\
&\quad\sim\int_{\bigcup_{i \in \blz_2} \Qi} \left[ \int_{\bigcup_{k \in \blz_2} \Qkk} \frac{| f(x)\sum_{j \in \Lambda_2} \psi_j(x) -  f(y)\sum_{j \in \Lambda_2} \psi_j(y)|^q}{|x - y|^{s q + n} {d(y,\pa\boz)^{aq}}}\,dy \right]^{\frac{p}{q}}\,dx \nonumber  \\
&\qquad+  \int_{\bigcup_{i \in \blz_2} \Qi} \left[ \int_{\boz \setminus \bigcup_{k \in \blz_2} \Qkk} \frac{| f(x)\sum_{j \in \Lambda_2} \psi_j(x) |^q}{|x - y|^{s q + n}{d(y,\pa\boz)^{aq}}}\,dy \right]^{\frac{p}{q}}\,dx  \nonumber \\
&\qquad+ \int_{\bozfour \setminus \bigcup_{i \in \blz_2} \Qi} \left[ \int_{\bigcup_{k \in \blz_2} \Qk} \frac{| f(x)\sum_{j \in \Lambda_2} \psi_j(x) -  f(y)\sum_{j \in \Lambda_2} \psi_j(y)|^q}{|x - y|^{s q + n}{d(y,\pa\boz)^{aq}}}\,dy \right]^{\frac{p}{q}}\,dx  \nonumber \\
&\qquad+ \int_{\boz \setminus \bozfour} \left[ \int_{\bigcup_{k \in \blz_2} \Qk} \frac{| f(y)\sum_{j \in \Lambda_2} \psi_j(y) |^q}{|x - y|^{s q + n}d(y,\pa\boz)^{aq}}\,dy \right]^{\frac{p}{q}}\,dx  \nonumber \\
&\quad=:\mathrm{I}_1(\vez)+\mathrm{I}_2(\vez)+\mathrm{II}(\vez)+\mathrm{III}(\vez).
\end{align}

{\bf Step 2}. In this step, we estimate $\mathrm{I}_2(\vez)$.
By Lemma \ref{lem2.1}(ii), we find that,
for any $x\in\bigcup_{i \in \blz_2} \Qi$ and $y\in\boz \setminus\bigcup_{k \in \blz_2} \Qkk$, 
$$d(x,\pa\boz)<d(y,\pa\boz),$$ which further implies
\begin{align}\label{1353}
\mathrm{I}_2(\vez)\leq \sum_{i \in \Lambda_2} \int_{\Qi} \frac{| f(x) |^p}{d(x,\pa\boz)^{ap}} \left( \int_{\boz \setminus \bigcup_{k \in \blz_2} \Qkk} \frac{1}{|x - y|^{s q + n}}\,dy \right)^{\frac{p}{q}} \,dx. 
\end{align}
Note that for any given $i\in\blz_2$, since $\boz \setminus\bigcup_{k \in \blz_2}\Qkk\subset\boz \setminus \Qii$, it follows that, for any \(x \in \Qi\) and \( y \in \boz \setminus \bigcup_{k \in \blz_2}\Qkk\),
\begin{align}\label{eqn-7.5}
|x-y|\geq \frac{6}{5}\cdot\frac{l_i}{2}-\frac{11}{10}\cdot\frac{l_i}{2} =\frac{l_i}{20}.
\end{align}
From this, \eqref{1353} and polar coordinates, we deduce that
\begin{align*}
\mathrm{I}_2(\vez)
&\lesssim \sum_{i \in \Lambda_2} \int_{\Qi} \frac{| f(x) |^p}{d(x,\pa\boz)^{ap}} \left( \int_{\frac{l_i}{20}}^\infty \rho^{-s q - 1}\,d\rho \right)^{\frac{p}{q}}\,dx\\
&\lesssim \sum_{i \in \Lambda_2} \int_{\Qi} \frac{| f(x) |^p}{d(x,\pa\boz)^{ap}}   {l_i}^{-s p}\,dx .
\end{align*}
By this, Lemma \ref{lem2.1}, \eqref{eqn-2.6} and the fact that $f$ satisfies the Hardy condition $(H)_{s,p,q,a,p,(s+a)p}$, we obtain
\begin{align*}
\mathrm{I}_2(\vez)
&\lesssim \sum_{i \in \Lambda_2} \int_{\Qi} \frac{| f(x)|^p}{d(x, \partial\Omega)^{(s+a) p}}\,dx 
\sim \int_{\bigcup_{i \in \Lambda_2} \Qi} \frac{| f(x)|^p}{d(x, \partial\Omega)^{(s+a) p}}\,dx \nonumber \\
&\leq \int_{\bozthreee}\frac{ |f(x)|^p}{d(x, \partial\Omega)^{(s+a)p}}\,dx \to 0 
\end{align*}
as $\varepsilon \to 0$. Thus,
\begin{equation}\label{eqn-4.6}
\lim_{\varepsilon \to 0} \mathrm{I}_2(\vez) = 0. 
\end{equation}

{\bf Step 3}. In this step, we estimate $\mathrm{III}(\vez)$. 
From \eqref{eqn-2.6}, \eqref{eqn-2.3} and Lemma \ref{lem2.1}(ii),
we infer that, for any $k\in\blz_2$,  
$$\Qk\subset\Omega_{3\varepsilon}
\ \ \text{and}\ \ l_k\lesssim \varepsilon.
$$
This implies that for any $x\in\Omega \setminus \Omega_{14\varepsilon}$ and $y\in\Qk$ with $k\in\blz_2$,
\begin{align}\label{eqn-8.8}
|x-y|\geq 14\varepsilon-3\varepsilon\gtrsim l_k,
\end{align}
which, combined with Lemma \ref{lem2.1}, \eqref{eqn-2.6},
the fact that $|\Omega|<\infty$ since $\Omega$ is bounded, and the fact that $f$ satisfies the Hardy condition $(H)_{s,p,q,a,q,(s+a)q+n}$, further implies that
\begin{align*}
\mathrm{III}(\vez) &\leq \int_{\Omega \setminus \Omega_{14\varepsilon}} \left[ \sum_{k\in\Lambda_2}\int_{ \Qk} \frac{| f(y)|^q}{|x - y|^{s q + n}d(y,\pa\boz)^{aq}}\,dy \right]^{\frac{p}{q}}\,dx \\
&\lesssim \int_{\Omega \setminus \Omega_{14\varepsilon}} \left[ \sum_{k\in\Lambda_2}\int_{ \Qk} \frac{| f(y)|^q}{{l_k}^{s q + n}d(y,\pa\boz)^{aq}}\,dy \right]^{\frac{p}{q}}\,dx  \\
&\lesssim \int_{\Omega \setminus \Omega_{14\varepsilon}} \left[ \sum_{k\in\Lambda_2}\int_{ \Qk} \frac{| f(y)|^q}{d(y,\pa\boz)^{(s+a)q+n}}\,dy \right]^{\frac{p}{q}}\,dx  \\
&\sim\int_{\Omega \setminus \Omega_{14\varepsilon}} \left[ \int_{\bigcup_{k \in \Lambda_2} \Qk} \frac{| f(y)|^q}{d(y,\pa\boz)^{(s+a)q+n}}\,dy \right]^{\frac{p}{q}}\,dx  \\
&\leq \left|\boz\right| \left[ \int_{\bozthreee} \frac{| f(y)|^q}{d(y,\pa\boz)^{(s+a)q+n}}\,dy \right]^{\frac{p}{q}}\to 0 
\end{align*}
as $\varepsilon \to 0$. Thus,
\begin{align}\label{eqn-4.7}
\lim_{\varepsilon \to 0}\mathrm{III}(\vez)= 0.
\end{align}

{\bf Step 4}. In this step, we estimate \( \mathrm{I}_1(\vez) \). To this end, we first decompose \( \mathrm{I}_1(\vez) \) as
\begin{align*}
\mathrm{I}_1(\vez)
&\lesssim \int_{\bigcup_{i \in \Lambda_2} \Qi} \left[ \int_{\bigcup_{k \in \Lambda_2} \Qkk} \frac{| f(x) - f(y)|^q |\sum_{j \in \Lambda_2} \psi_j(y)|^q}{|x - y|^{s q + n}d(y,\pa\boz)^{aq}}\,dy \right]^{\frac{p}{q}}\,dx \\
&\quad+ \int_{\bigcup_{i \in \Lambda_2} \Qi} \left[ \int_{\bigcup_{k \in \Lambda_2} \Qkk} \frac{| f(x)|^q |\sum_{j \in \Lambda_2} \psi_j(x) - \sum_{j \in \Lambda_2} \psi_j(y)|^q}{|x - y|^{s q + n}d(y,\pa\boz)^{aq}}\,dy \right]^{\frac{p}{q}}\,dx \\
&=: \mathrm{I}_{1,1}(\vez) + \mathrm{I}_{1,2}(\vez).
\end{align*}
For \( \mathrm{I}_{1,1}(\vez) \), by \eqref{eqn-2.1} and \eqref{eqn-2.6}, we obtain
\begin{align}\label{eqn-4.8}
\mathrm{I}_{1,1}(\vez)
&\leq \int_{\bigcup_{i \in \Lambda_2} \Qi} \left[ \int_{\bigcup_{k \in \Lambda_2} \Qkk} \frac{| f(x) - f(y)|^q}{|x - y|^{s q + n}d(y,\pa\boz)^{aq}}\,dy \right]^{\frac{p}{q}}\,dx \nonumber \\
&\leq \int_{\Omega_{3\varepsilon}} \left[ \int_{\Omega_{3\varepsilon}} \frac{| f(x) -  f(y)|^q}{|x - y|^{s q + n}d(y,\pa\boz)^{aq}}\,dy \right]^{\frac{p}{q}}\,dx.
\end{align}
Since \( f \in \huafwe \), it follows from \eqref{eqn-4.8} that $\mathrm{I}_{1,1}(\vez)\to 0$ as \( \varepsilon \to 0 \).

Next, we estimate \( \mathrm{I}_{1,2}(\vez) \). By \eqref{eqn-2.8}, we write
\begin{align}\label{eqn-4.9}
\mathrm{I}_{1,2}(\vez)
&\leq \sum_{i \in \Lambda_2} \int_{\Qi} \left[ \int_{\bigcup_{k \in \Lambda_2} \Qkk} \frac{| f(x)|^q |\sum_{j \in \Lambda_2} \psi_j(x) - \sum_{j \in \Lambda_2} \psi_j(y)|^q}{|x - y|^{s q + n}d(y,\pa\boz)^{aq}}\,dy \right]^{\frac{p}{q}}\,dx \nonumber \\
&\sim \sum_{i \in \Lambda_2} \int_{\Qi} \left[ \int_{\bigcup_{k \in \Lambda_{2,1}(i)} \Qkk} \frac{| f(x)|^q |\sum_{j \in \Lambda_2} \psi_j(x) - \sum_{j \in \Lambda_2} \psi_j(y)|^q}{|x - y|^{s q + n}d(y,\pa\boz)^{aq}}\,dy \right]^{\frac{p}{q}}\,dx \nonumber \\
&\quad+ \sum_{i \in \Lambda_2} \int_{\Qi} \left[ \int_{\bigcup_{k \in \Lambda_{2,2}(i)} \Qkk} \frac{| f(x)|^q |\sum_{j \in \Lambda_2} \psi_j(x) - \sum_{j \in \Lambda_2} \psi_j(y)|^q}{|x - y|^{s q + n}d(y,\pa\boz)^{aq}}\,dy \right]^{\frac{p}{q}}\,dx \nonumber \\
&=: \mathrm{I}_{1,2}^1(\vez) + \mathrm{I}_{1,2}^2(\vez).
\end{align}
By $\psi_j \in C_{\mathrm{c}}^{\fz}(\boz)$, \eqref{eqn-2.9}, the mean value theorem, \eqref{eqn-2.2}, Lemmas \ref{lem2.4} and \ref{lem2.1}, \eqref{eqn-2.8}, polar coordinates, \eqref{eqn-2.6} and the fact that $f$ satisfies the Hardy condition $(H)_{s,p,q,a,p,(s+a)p}$, we conclude that
\begin{align}\label{eqn-4.10}
\mathrm{I}_{1,2}^1(\vez)
&\leq \sum_{i \in \Lambda_2} \int_{\Qi} | f(x)|^p \left[ \sum_{k \in \Lambda_{2,1}(i)} \int_{\Qkk} \frac{|\sum_{j \in \Lambda_{2,3}(i,k)} \psi_j(x) -\sum_{j \in \Lambda_{2,3}(i,k)} \psi_j(y)|^q}{|x - y|^{s q + n}d(y,\pa\boz)^{aq}}\,dy \right]^{\frac{p}{q}}\,dx  \nonumber \\
&\lesssim \sum_{i \in \Lambda_2} \int_{\Qi} | f(x)|^p \left[ \sum_{k \in \Lambda_{2,1}(i)} \int_{\Qkk} \sum_{j \in \Lambda_{2,3}(i,k)}{\| \nabla \psi_j \|_{L^\infty}^q |x - y|^{q(1 - s)-n}{l_k}^{-aq}}\,dy \right]^{\frac{p}{q}}\,dx \nonumber \\
&\lesssim \sum_{i \in \Lambda_2} \int_{\Qi} | f(x)|^p \left[ \sum_{k \in \Lambda_{2,1}(i)} \int_{\Qkk} {{l_i}^{-q} |x - y|^{q(1 - s)-n}{l_i}^{-aq}}\,dy \right]^{\frac{p}{q}}\,dx  \nonumber \\
&\sim \sum_{i \in \Lambda_2} \int_{\Qi} | f(x)|^p \left[ \sum_{k \in \Lambda_{2,1}(i)} \int_{\Qkk} {{d(x,\pa\boz)}^{-(1+a)q} |x - y|^{q(1 - s)-n}}\,dy \right]^{\frac{p}{q}}\,dx  \nonumber \\
&\leq \sum_{i \in \Lambda_2} \int_{\Qi} \frac{ |f(x)|^p}{{d(x,\pa\boz)}^{(1+a)p}} \left[ \int_0^{cl_i} \rho^{q(1 - s) - 1}\,d\rho \right]^{\frac{p}{q}}\,dx \nonumber \\
&\sim \sum_{i \in \Lambda_2} \int_{\Qi} \frac{ |f(x)|^p}{{d(x,\pa\boz)}^{(1+a)p}}  {l_i}^{p(1-s)}\,dx \sim \int_{\bigcup_{i \in \Lambda_2} \Qi} \frac{| f(x)|^p}{d(x, \partial\Omega)^{(s+a) p}}\,dx \nonumber \\
&\leq\int_{\bozthreee}\frac{ |f(x)|^p}{d(x, \partial\Omega)^{(s+a)p}}\, dx  \to 0
\end{align}
as $\varepsilon \to 0$. For \( \mathrm{I}_{1,2}^2(\vez) \), we write
\begin{align*}
\mathrm{I}_{1,2}^2(\vez)
&\sim \sum_{i \in \Lambda_2} \int_{\Qi} \left[ \sum_{\substack{k\in\Lambda_{2,2}(i) \\ l_k\geq l_i}}\int_{\Qkk} \frac{| f(x)|^q |\sum_{j \in \Lambda_2} \psi_j(x) - \sum_{j \in \Lambda_2} \psi_j(y)|^q}{|x - y|^{s q + n}d(y,\pa\boz)^{aq}}\,dy \right]^{\frac{p}{q}}\,dx \nonumber \\
&\quad+\sum_{i \in \Lambda_2} \int_{\Qi} \left[ \sum_{\substack{k\in\Lambda_{2,2}(i) \\ l_k< l_i}}\int_{\Qkk} \frac{| f(x)|^q |\sum_{j \in \Lambda_2} \psi_j(x) - \sum_{j \in \Lambda_2} \psi_j(y)|^q}{|x - y|^{s q + n}d(y,\pa\boz)^{aq}}\,dy \right]^{\frac{p}{q}}\,dx \nonumber \\
&=:\mathrm{I}_{1,2}^{2,1}(\vez)+\mathrm{I}_{1,2}^{2,2}(\vez).
\end{align*} 
To estimate $\mathrm{I}_{1,2}^{2,1}(\vez)$, by Lemma \ref{lem2.1}(v), we obtain, for any $x\in\Qi$ and $y\in\Qkk$ with $i\in\Lambda_2$ and $k\in\Lambda_{2,2}(i)$ satisfying $l_k\geq l_i$,
\begin{align}\label{eqn-8.13}
|x-y|\geq\frac{l_k}{4}-\left(\frac{l_i}{20}+\frac{l_k}{10}\right)\geq\frac{l_k}{4}-\frac{3}{20}l_k=\frac{l_k}{20}\geq\frac{l_i}{20},
\end{align}
which, combined with \eqref{eqn-2.1}, Lemma \ref{lem2.1}, polar coordinates, \eqref{eqn-2.6} and the fact that $f$ satisfies the Hardy condition $(H)_{s,p,q,a,p,(s+a)p}$, further implies that
\begin{align}\label{eqn-4.11}
\mathrm{I}_{1,2}^{2,1}(\vez)&\lesssim \sum_{i \in \Lambda_2} \int_{\Qi} \left[ \sum_{k \in \Lambda_{2,2}(i)} \int_{\Qkk} \frac{| f(x)|^q }{|x - y|^{s q + n}d(y,\pa\boz)^{aq}}\,dy \right]^{\frac{p}{q}}\,dx \nonumber \\
&\sim \sum_{i \in \Lambda_2} \int_{\Qi} \left[ \sum_{k \in \Lambda_{2,2}(i)} \int_{\Qkk} \frac{| f(x)|^q }{|x - y|^{s q + n}{l_i}^{aq}}\,dy \right]^{\frac{p}{q}}\,dx \nonumber \\
&\sim\sum_{i \in \Lambda_2} \int_{\Qi} \frac{ |f(x)|^p}{{d(x,\pa\boz)}^{ap}} \left( \int_{\frac{l_i}{20}}^{\infty} \rho^{- sq - 1}\,d\rho \right)^{\frac{p}{q}}\,dx \nonumber \\
&\sim\sum_{i \in \Lambda_2} \int_{\Qi} \frac{ |f(x)|^p}{{d(x,\pa\boz)}^{ap}}  {l_i}^{-sp}\,dx
\nonumber\\
&\sim \int_{\bigcup_{i \in \Lambda_2} \Qi} \frac{| f(x)|^p}{d(x, \partial\Omega)^{(s+a) p}}\,dx \nonumber \\
&\leq \int_{\bozthreee}\frac{ |f(x)|^p}{d(x, \partial\Omega)^{(s+a)p}}\,dx \to 0 
\end{align}
as $\varepsilon \to 0$. To estimate $\mathrm{I}_{1,2}^{2,2}(\vez)$, from the fact that $0\leq\sum_{j\in\Lambda_2}\psi_j(z)\leq\sum_{j\in\nn}\psi_j(z)=1$ for any $z\in\boz$, we infer that
\begin{align}\label{eqn-8.14}
\mathrm{I}_{1,2}^{2,2}(\vez)
&\leq\sum_{i \in \Lambda_2} \int_{\Qi} \left[ \sum_{\substack{k\in\Lambda_{2,2}(i) \\ l_k< l_i}}\int_{\Qkk} \frac{| f(x)|^q }{|x - y|^{s q + n}d(y,\pa\boz)^{aq}}\,dy \right]^{\frac{p}{q}} \,dx \nonumber\\
&\lesssim \sum_{i \in \Lambda_2} \int_{\Qi} \left[ \sum_{\substack{k\in\Lambda_{2,2}(i) \\ l_k< l_i}}\int_{\Qkk} \frac{| f(x)-f(y)|^q }{|x - y|^{s q + n}d(y,\pa\boz)^{aq}}\,dy \right]^{\frac{p}{q}}\,dx \nonumber\\
&\quad+ \sum_{i \in \Lambda_2} \int_{\Qi} \left[ \sum_{\substack{k\in\Lambda_{2,2}(i) \\ l_k< l_i}}\int_{\Qkk}\frac{| f(y)|^q }{|x - y|^{s q + n}d(y,\pa\boz)^{aq}}\,dy \right]^{\frac{p}{q}} \,dx. 
\end{align}
By \( f \in \huafwe \) and an argument similar to that used in the estimation of \eqref{eqn-4.8}, we find that the first term on the right-hand side of \eqref{eqn-8.14} tends to 0 as $\vez\to 0$, and by  an argument similar to that used in the estimation of  $\mathrm{III}(\vez)$, we find that the second term on the right-hand side of \eqref{eqn-8.14} tends to 0 as $\vez\to 0$. Hence,
\begin{align}\label{eqn-4.12}
\lim_{\varepsilon \to 0}\mathrm{I}_{1,2}^{2,2}(\vez)= 0. 
\end{align}
Combining \eqref{eqn-4.8}--\eqref{eqn-4.12}, we obtain
\begin{align}\label{eqn-4.13}
\lim_{\varepsilon \to 0}\mathrm{I}_{1}(\vez)= 0. 
\end{align}

{\bf Step 5}. Finally, we estimate \( \mathrm{II}(\vez) \). To this end, we write
\begin{align}\label{1713}
\mathrm{II}(\vez)
&\leq \int_{\bozfour} \left[ \int_{\bigcup_{k \in \blz_2} \Qk} \frac{| f(x)\sum_{j \in \Lambda_2} \psi_j(x) -  f(y)\sum_{j \in \Lambda_2} \psi_j(y)|^q}{|x - y|^{s q + n}{d(y,\pa\boz)^{aq}}}\,dy \right]^{\frac{p}{q}}\,dx \nonumber\\
&\lesssim \int_{\bozfour} \left[ \int_{\bigcup_{k \in \blz_2} \Qk} \frac{| f(x) - f(y)|^q |\sum_{j \in \Lambda_2} \psi_j(y)|^q}{|x - y|^{s q + n}d(y,\pa\boz)^{aq}}\,dy \right]^{\frac{p}{q}}\,dx \nonumber\\
&\quad+ \int_{\bozfour} \left[ \int_{\bigcup_{k \in \blz_2} \Qk} \frac{| f(x)|^q |\sum_{j \in \Lambda_2} \psi_j(x) - \sum_{j \in \Lambda_2} \psi_j(y)|^q}{|x - y|^{s q + n}d(y,\pa\boz)^{aq}}\,dy \right]^{\frac{p}{q}}\,dx.
\end{align}
By an argument similar to that used in the estimation of \( \mathrm{I}_{1,1}(\vez) \), we find that the first term on the right-hand side of \eqref{1713} tends to 0 as $\vez\to 0$, and by an argument similar to that used in the estimation of  \( \mathrm{I}_{1,2}(\vez) \), we find that the second term on the right-hand side of \eqref{1713} tends to 0 as $\vez\to 0$. Hence,
\begin{align}\label{eqn-4.14}
\lim_{\varepsilon \to 0} \mathrm{II}(\vez) = 0.
\end{align}
Combining \eqref{eqn-4.4}, \eqref{eqn-4.6}, \eqref{eqn-4.7}, \eqref{eqn-4.13} and \eqref{eqn-4.14}, we obtain 
$$\lim_{\huae \to 0} | \we |_{\huafwe}=0,$$ which proves \eqref{eqn-4.3} and hence 
completes the proof of (i).

Next, we prove (ii). Based on the proof of (i), we only need to show 
how to estimate the term $\mathrm{III}(\vez)$ in Step 3
without using the Hardy condition $\|\frac{f}{d^{s+a+n/q}}\|_{L^q(\Omega)}<\infty$.
Indeed, by the fact that $0\leq\sum_{j\in\Lambda_2}\psi_j(z)\leq\sum_{j\in\nn}\psi_j(z)=1$ 
for any $z\in\boz$, we write
\begin{align*}
\mathrm{III}(\vez)
&\leq \int_{\Omega \setminus \Omega_{14\varepsilon}} \left[ \int_{\bigcup_{k \in \Lambda_2} \Qk} \frac{| f(y)|^q}{|x - y|^{s q + n}d(y,\pa\boz)^{aq}}\,dy \right]^{\frac{p}{q}}\,dx \nonumber\\
&\lesssim \int_{\Omega \setminus \Omega_{14\varepsilon}} \left[ \int_{\bigcup_{k \in \Lambda_2} \Qk} \frac{| f(x) - f(y)|^q}{|x - y|^{s q + n}d(y,\pa\boz)^{aq}}\,dy \right]^{\frac{p}{q}}\,dx \nonumber\\
&\quad+ \int_{\Omega \setminus \Omega_{14\varepsilon}} \left[ \int_{\bigcup_{k \in \Lambda_2} \Qk} \frac{|f(x) |^q}{|x - y|^{s q + n}d(y,\pa\boz)^{aq}}\,dy \right]^{\frac{p}{q}}\,dx \nonumber\\
&=:\mathrm{III}_1(\vez) + \mathrm{III}_2(\vez).
\end{align*}
To estimate $\mathrm{III}_1(\vez)$, according to \( f \in \huafwe \) and \eqref{eqn-2.6}, we have
\begin{align*}
\lim_{\varepsilon \to 0}\mathrm{III}_1(\vez) \leq \lim_{\varepsilon \to 0}\int_\Omega \left[ \int_{\Omega_{3\varepsilon}} \frac{|f(x) -  f(y)|^q}{|x - y|^{s q + n}d(y,\pa\boz)^{aq}}\,dy \right]^{\frac{p}{q}} \,dx = 0. 
\end{align*}
To estimate \( \mathrm{III}_2(\vez) \), for any \( x \in \Omega \), define
$$
K_\varepsilon(x)
:=
\int_{\bigcup_{k\in\Lambda_2}\frac{11}{10}Q_k}
\frac{dy}
{|x-y|^{sq+n}d(y,\partial\Omega)^{aq}}.
$$
Then it is easy to see that
\begin{align*}
\mathrm{III}_2(\varepsilon)
\le
\int_{\Omega}
|f(x)|^p K_\varepsilon(x)^\frac{p}{q}dx.
\end{align*}
We claim that for a.e. $x\in \Omega$, it holds that

\begin{align}\label{eqn=3.x22}
K_\varepsilon(x)
\lesssim
d(x,\partial\Omega)^{-(s+a)q}
\end{align}
and 
\begin{align}\label{eqn=3.x22x}
\lim_{\vez\to0}K_\varepsilon(x)= 0 .
\end{align}

Assuming these two estimates for the moment, we obtain
\[
|f(x)|^pK_\varepsilon(x)^\frac{p}{q}
\lesssim
\frac{|f(x)|^p}{d(x,\partial\Omega)^{(s+a)p}},
\]
which, combined with the Lebesgue dominated convergence theorem (we used the fact that $f$ satisfies the Hardy condition $(H)_{s,p,q,a,p,(s+a)p}$ here), further implies that 
\[
\lim_{\varepsilon\to0}{\rm III}_2(\varepsilon)=0.
\]

It remains to verify \eqref{eqn=3.x22} and \eqref{eqn=3.x22x}.
Fix $x\in\Omega$, and choose the Whitney cube $Q_i$ such that $x\in Q_i$. 
If $x\in\Omega\setminus\Omega_{14\varepsilon}$, then $i\in\Lambda_3$. For $y\in\frac{11}{10}Q_k$ with $k\in\Lambda_2$, by 
Lemma \ref{lem2.2}(iii), with harmless changes of constants due to the dilation $\frac{11}{10}Q_k$, we have
\[
|x-y|\sim D(Q_i,Q_k).
\]
Moreover, for any $y\in \frac{11}{10}Q_k$, it holds that
$
d(y,\partial\Omega)\sim l_k.
$
Hence,
\[
K_\varepsilon(x)
\lesssim
\sum_{k\in\Lambda_2}
\frac{l_k^n}
{D(Q_i,Q_k)^{sq+n}l_k^{aq}}
=
\sum_{k\in\Lambda_2}
\frac{l_k^{n-aq}}
{D(Q_i,Q_k)^{sq+n}}.
\]
Since $\Lambda_2\subset\mathbb N_+$, we may enlarge the sum and obtain
\[
K_\varepsilon(x)
\lesssim
\sum_{k=1}^\infty
\frac{l_k^{n-aq}}
{D(Q_i,Q_k)^{sq+n}}.
\]
Now we apply Lemma~\ref{lem2.3} to obtain
\[
\sum_{k=1}^\infty
\frac{l_k^{n-aq}}
{D(Q_i,Q_k)^{sq+n}}
\lesssim
l_i^{n-aq-(sq+n)}
=
l_i^{-(s+a)q}.
\]
Since $x\in Q_i$, we have 
\[
K_\varepsilon(x)
\lesssim
d(x,\partial\Omega)^{-(s+a)q}.
\]
This proves \eqref{eqn=3.x22}.

Finally, we prove the pointwise convergence \eqref{eqn=3.x22x}. For the fixed cube $Q_i$ containing $x$, we have
\[
K_\varepsilon(x)
\lesssim
\sum_{k\in\Lambda_2}
\frac{l_k^{n-aq}}
{D(Q_i,Q_k)^{sq+n}}.
\]
The full series
\[
\sum_{k=1}^\infty
\frac{l_k^{n-aq}}
{D(Q_i,Q_k)^{sq+n}}
\]
is finite by Lemma~\ref{lem2.3}. Moreover, for each fixed $k$, the condition $k\in\Lambda_2$ means
$
d(Q_k,\partial\Omega)<\varepsilon.
$
Thus, for every fixed $k$, when $\varepsilon<d(Q_k,\partial\Omega)$, we have $k\notin\Lambda_2$. Hence,
$
\mathbf{1}_{\Lambda_2}(k)\to0
$
as $\vez\to0$ for every fixed $k$. By dominated convergence for series,
\[
\sum_{k\in\Lambda_2}
\frac{l_k^{n-aq}}
{D(Q_i,Q_k)^{sq+n}}
\to0
\]
as $\vez\to0$. Consequently, $K_\varepsilon(x)\to 0$ as $\vez\to0$
for every $x\in\Omega$.
This verifies \eqref{eqn=3.x22x}, which finally proves (ii)
and hence completes the proof of Theorem \ref{thm-main}.
\end{proof}

\section{The case $s\in(1,\infty)\setminus\mathbb{N}$}\label{s4}

In this section, we show the density property for Triebel--Lizorkin 
spaces on $\boz$ in higher-order smoothness 
case $s\in(1,\infty)\setminus\mathbb{N}$ by proving Theorem \ref{thm-secmain}.

\begin{proof}[Proof of Theorem \ref{thm-secmain}]
We first prove (i). Let 
\( f \in C^\infty(\boz)\cap\widetilde{\mathscr{F}}_{p,q}^{s}(\boz)\),  
\( \{Q_j\}_{j=1}^{\infty} \) and \( \{\psi_j\}_{j=1}^{\infty} \subset C_{\mathrm{c}}^\infty(\mathbb{R}^n) \) 
be the Whitney decomposition and its associated partition of unity fixed in Section~\ref{s2}. For any \( \varepsilon > 0 \) and \( x \in \Omega \), by \eqref{eqn-2.1} and the definitions of \( \Lambda_1\) and \( \Lambda_2 \) in \eqref{eqn-2.3}, we write
\begin{align}\label{eqn-4.19}
f(x)= \sum_{j \in \Lambda_1} \psi_j(x) f(x) + \sum_{j \in \Lambda_2} \psi_j(x) f(x)
=: v_{\varepsilon}(x) + w_{\varepsilon}(x).
\end{align}
Similar to the argument used in the proof of Theorem \ref{thm-main},  it suffices to prove
$$
\lim_{\huae \to 0} \| \we \|_{\widetilde{\mathscr{F}}_{p,q}^{s}(\boz)}=0,
$$
which, combined with \eqref{eqn-GWN}, is equivalent to verify the following four identities:
\begin{align}\label{eqn-4.20}
\lim_{\huae \to 0} \| \we \|_{W^{k,p}(\Omega)}=0,
\end{align}
\begin{align}\label{1048}
\lim_{\varepsilon\to0}
\sum_{\substack{\bz\in\nno^n \\ |\bz| \leq k}}
\left(\left\|\frac{D^\beta w_\varepsilon}{d^{s-|\beta|}}\right\|_{L^p(\Omega)}+\left\|\frac{D^\beta w_\varepsilon}{d^{s-|\beta|+\frac{n}{q}}}\right\|_{L^q(\Omega)}\right)=0,
\end{align}
\begin{align}\label{eqn-4.21}
\lim_{\huae \to 0} \adk \left| D^{\az} \we \right|_{\cf_{p,q}^{\sz}(\boz)}=0
\end{align}
and
\begin{align}\label{eqn-4.22}
\lim_{\huae \to 0}
\sum_{\substack{\bz\in\nno^n \\ | \bz| \leq k}} \left|D^{\bz}\we\right|_{\cf_{p,q}^{\sz}(\boz,d^{-(k-|\bz|)q})}=0.
\end{align}
We divide the proofs of \eqref{eqn-4.20}--\eqref{eqn-4.22} into several steps.

{\bf Step 1}. In this step, we show \eqref{eqn-4.20}
and \eqref{1048}. By \eqref{eqn-1.3}, we have
\[
\| \we \|_{W^{k,p}(\boz)}=\sum_{\substack{\az\in\nno^n \\ |\alpha| \leq k}}
\left\| D^{\alpha}\we \right\|_{L^p(\Omega)}.
\]
Let $\az\in\nno^n$ and $|\az|\leq k$. To estimate \(\| D^{\alpha}\we \|_{L^p(\Omega)}\), by \eqref{eqn-4.19}, Lemma \ref{lem2.1}(iv), Leibniz's rule and \( \text{supp}\,\psi_j \subset \frac{11}{10}Q_j \) for any $j\in\Lambda_2$, we find that
\begin{align*}
\left\|D^{\alpha}\we\right\|_{L^p(\Omega)}^p
&= \int_{\boz} \left| \sum_{j \in \Lambda_2} D^{\alpha} (f \psi_j)(x) \right|^p\,dx\\
&\lesssim \sum_{j \in \Lambda_2} \int_{\Qj} | D^{\alpha} (f \psi_j)(x) |^p\,dx \\
&=\sum_{j \in \Lambda_2} \int_{\Qj} \left| \sum_{\substack{\bz\in\nno^n \\ \bz \leq \az}} \binom{\az}{\bz} D^{\bz}f(x) D^{\az - \bz} \psi_j(x) \right|^p\,dx \\
&\lesssim \sum_{\substack{\bz\in\nno^n \\ \bz \leq \az}} \binom{\az}{\bz} \sum_{j \in \Lambda_2} \int_{\Qj} | D^{\bz}f(x) D^{\az - \bz} \psi_j(x) |^p\,dx.
\end{align*}
According to \eqref{eqn-2.2}, Lemma \ref{lem2.1}(ii) and \eqref{eqn-2.6}, we obtain for any $\bz\in\nno^n$ with $\bz\leq\az$,
\begin{align*}
&\sum_{j \in \Lambda_2}\int_{\Qj} \left| D^{\bz}f(x) D^{\az - \bz} \psi_j(x) \right|^p\, dx\\
&\quad\lesssim \sum_{j \in \Lambda_2}\int_{\Qj} \frac{|D^{\bz}f(x)|^p}{(\mathrm{diam\,} Q_j)^{|\az-\bz|p}}\,dx\\
&\quad\lesssim \int_{\bigcup_{j \in \Lambda_2} \Qj} \frac{|D^{\bz}f(x)|^p}{d(x,\pa\boz)^{|\az-\bz|p}}\,dx 
\leq \int_{\boz_{3\huae}} \frac{|D^{\bz}f(x)|^p}{d(x,\pa\boz)^{|\az-\bz|p}}\,dx,
\end{align*}
which tends to 0 as \( \huae \to 0 \) by the fact
that $\|\frac{D^\beta f}{d^{|\alpha-\beta|}}\|_{L^p(\Omega)}<\infty$. Thus,
\[
\left\| D^{\alpha}\we\right\|_{L^p(\Omega)}^p \lesssim  \sum_{\substack{\bz\in\nno^n \\ \bz \leq \az}} \binom{\az}{\bz} \int_{\boz_{3\huae}} \frac{|D^{\bz}f(x)|^p}{d(x,\pa\boz)^{|\az-\bz|p}} \,dx \to 0
\]
as $\huae \to 0$.
Hence,
\[
\lim_{\huae \to 0} \| \we \|_{W^{k,p}(\Omega)}
=\lim_{\huae \to 0}  \sum_{|\alpha| \leq k} \left\| D^{\alpha}\we \right\|_{L^p(\Omega)}
=\sum_{|\alpha| \leq k} \lim_{\huae \to 0} \left\| D^{\alpha}\we \right\|_{L^p(\Omega)}=0.
\]
This completes the proof of \eqref{eqn-4.20}.
The proof of \eqref{1048} is similar.

{\bf Step 2}. To show \eqref{eqn-4.21}, by \eqref{eqn-4.19}, Lemma \ref{lem2.1}(iv) and Leibniz's rule, we obtain for any $\az\in\nno^n$ with $|\az|=k$,
\begin{align}\label{eqn-8.5}
&\left| D^{\az} \we \right|_{\cf_{p,q}^{\sz}(\boz)}^p \nonumber\\
&\quad= \left| D^{\alpha} \left(f \sum_{j \in \Lambda_2} \psi_j\right) \right|_{\cf_{p,q}^{\sz}(\boz)}^p \nonumber\\
&\quad= \int_{\boz} \left[\int_{\boz} \frac{|D^{\alpha} (f \sum_{j \in \Lambda_2} \psi_j)(x) - D^{\alpha} (f \sum_{j \in \Lambda_2} \psi_j)(y)|^q}{|x - y|^{\sz q + n}}\,dy \right]^{\frac{p}{q}}\,dx \nonumber\\
&\quad\lesssim \sum_{\substack{\bz\in\nno^n \\ \bz \leq \az}} \binom{\az}{\bz} \int_{\boz} \left[ \int_{\boz} \frac{| \fpjzuhe -  \fpjzuhey|^q}{|x - y|^{\sz q + n}}\,dy \right]^{\frac{p}{q}}\,dx.
\end{align}
Now, we fix $j\in\Lambda_2$ and $\az,\bz\in\nno^n$ satisfying both $|\az|=k$ and $\bz\leq\az$. Since supp \(\psi_j \subseteq \Qj\), we write
\begin{align}\label{eqn-4.23}
&\int_{\boz} \left[\int_{\boz} \frac{| \fpjzuhe -  \fpjzuhey|^q}{|x - y|^{\sz q + n}}\,dy \right]^{\frac{p}{q}}\,dx \nonumber \\
&\quad= \left[ \int_{\bigcup_{i \in \blz_2} \Qi}
+ \int_{\bozfour \setminus \bigcup_{i \in \blz_2} \Qi}
+ \int_{\boz \setminus \bozfour} \right]
\left(\int_{\boz}\dots\,dy\right)^{\frac{p}{q}}\,dx \nonumber \\
&\quad\sim \int_{\bigcup_{i \in \blz_2} \Qi}\nonumber \\
&\quad\quad\times\left[ \int_{\bigcup_{k \in \blz_2} \Qkk}
\frac{| \fpjzuhe -  \fpjzuhey|^q}{|x - y|^{\sz q + n}}\,dy \right]^{\frac{p}{q}}\,dx \nonumber \\
&\qquad+ \int_{\bigcup_{i \in \blz_2} \Qi}
\left[ \int_{\boz \setminus \bigcup_{k \in \blz_2} \Qkk}
\frac{| \fpjzuhe |^q}{|x - y|^{\sz q + n}} \,dy \right]^{\frac{p}{q}}\,dx \nonumber \\
&\qquad+ \int_{\bozfour \setminus \bigcup_{i \in \blz_2} \Qi} \nonumber \\
&\qquad\times \left[ \int_{\bigcup_{k \in \blz_2} \Qk}
\frac{| \fpjzuhe -  \fpjzuhey|^q}{|x - y|^{\sz q + n}}\,dy \right]^{\frac{p}{q}}\,dx \nonumber \\
&\qquad+ \int_{\boz \setminus \bozfour}
\left[ \int_{\bigcup_{k \in \blz_2} \Qk}
\frac{| \fpjzuhey |^q}{|x - y|^{\sz q + n}}\,dy \right]^{\frac{p}{q}}\,dx \nonumber \\
&\quad=:\tilde{\mathrm{I}}_1(\vez)+\tilde{\mathrm{I}}_2(\vez)
+\tilde{\mathrm{II}}(\vez)+\tilde{\mathrm{III}}(\vez).
\end{align}

{\bf Step 3}. In this step, we estimate $\tilde{\mathrm{I}}_2(\vez)$. By \eqref{eqn-2.2}, \eqref{eqn-2.6}, and both (ii) and (iv) of Lemma \ref{lem2.1}, we find that for any $z\in\boz$,
\begin{align}\label{eqn-xiutwo}
\left| \sum_{j\in\Lambda_2}D^{\az-\bz}\psi_j(z) \right|\lesssim \chi_{\boz_{3\vez}}(z)\sum_{j\in\Lambda_2}\frac{1}{(\mathrm{diam\,}Q_j)^{|\az-\bz|}}\lesssim \frac{\chi_{\boz_{3\vez}}(z)}{d(z,\pa\boz)^{|\az-\bz|}}.
\end{align}
From this, \eqref{eqn-7.5} and polar coordinates, we deduce that
\begin{align*}
\tilde{\mathrm{I}}_2(\vez)
&\lesssim \sum_{i \in \Lambda_2} \int_{\Qi} \frac{| D^{\bz}f(x) |^p}{d(x,\pa\boz)^{|\az-\bz|p}} \left( \int_{\boz \setminus \bigcup_{k \in \blz_2} \Qkk} \frac{1}{|x - y|^{\sz q + n}}\,dy \right)^{\frac{p}{q}}\,dx\\
&\lesssim \sum_{i \in \Lambda_2} \int_{\Qi} \frac{| D^{\bz}f(x) |^p}{d(x,\pa\boz)^{|\az-\bz|p}} \left( \int_{\frac{l_i}{20}}^\infty \rho^{-\sz q - 1}\,d\rho \right)^{\frac{p}{q}}\,dx \\
&\sim \sum_{i \in \Lambda_2} \int_{\Qi} \frac{| D^{\bz}f(x) |^p}{d(x,\pa\boz)^{|\az-\bz|p}} {l_i}^{-\sz p}\,dx.
\end{align*}
By this, Lemma \ref{lem2.1}, \eqref{eqn-2.6} and 
the fact
that $\|\frac{D^\beta f}{d^{s-|\beta|}}\|_{L^p(\Omega)}<\infty$, we obtain
\begin{align*}
\tilde{\mathrm{I}}_2(\vez)
&\lesssim \sum_{i \in \Lambda_2} \int_{\Qi} \frac{|D^\beta f(x)|^p}{d(x, \partial\Omega)^{(\sz+|\az-\bz|) p}}\,dx\\
&\sim \int_{\bigcup_{i \in \Lambda_2} \Qi} \frac{|D^\beta f(x)|^p}{d(x, \partial\Omega)^{(\sz+|\az-\bz|) p}}\,dx \nonumber \\
&\leq \int_{\bozthreee}  \frac{|D^\beta f(x)|^p}{d(x, \partial\Omega)^{(s-|\bz|) p}}\,dx \to 0
\end{align*}
as $\varepsilon \to 0.$ Thus,
\begin{equation}\label{eqn-4.25}
\lim_{\varepsilon \to 0} \tilde{\mathrm{I}}_2(\vez) = 0.
\end{equation}

{\bf Step 4}. In this step, we estimate \(\tilde{\mathrm{III}}(\vez)\). By \eqref{eqn-xiutwo}, \eqref{eqn-8.8}, Lemma \ref{lem2.1}, \eqref{eqn-2.6},
the fact that $|\Omega|<\infty$ (since $\Omega\subset\mathbb{R}^n$ is bounded),
and the fact
that $\|\frac{D^\beta f}{d^{s-|\beta|+n/q}}\|_{L^q(\Omega)}<\infty$, we find that
\begin{align*}
\tilde{\mathrm{III}}(\vez) &\leq\int_{\Omega \setminus \Omega_{14\varepsilon}} \left[ \sum_{k\in\Lambda_2}\int_{ \Qk} \frac{| D^{\bz} f(y)|^q}{|x - y|^{\sz q + n}d(y,\pa\boz)^{|\az-\bz|q}}\,dy \right]^{\frac{p}{q}}\,dx \\
&\lesssim \int_{\Omega \setminus \Omega_{14\varepsilon}} \left[ \sum_{k\in\Lambda_2}\int_{ \Qk} \frac{|D^\bz f(y)|^q}{{l_k}^{\sz q + n}d(y,\pa\boz)^{|\az-\bz|q}}\,dy \right]^{\frac{p}{q}}\,dx \\
&\sim \int_{\Omega \setminus \Omega_{14\varepsilon}} \left[ \sum_{k\in\Lambda_2}\int_{ \Qk} \frac{|D^\bz f(y)|^q}{d(y,\pa\boz)^{(\sz+|\az-\bz|)q+n}}\,dy \right]^{\frac{p}{q}} \,dx \\
&\sim \int_{\Omega \setminus \Omega_{14\varepsilon}} \left[ \int_{\bigcup_{k \in \Lambda_2} \Qk} \frac{|D^\bz f(y)|^q}{d(y,\pa\boz)^{(\sz+|\az-\bz|)q+n}}\,dy \right]^{\frac{p}{q}} \,dx \\
&\leq |\boz|\left[ \int_{\bozthreee} \frac{|D^\bz f(y)|^q}{d(y,\pa\boz)^{(s-|\bz|)q+n}}\,dy \right]^{\frac{p}{q}} \to 0
\end{align*}
as $\varepsilon \to 0.$ Thus,
\begin{equation}\label{eqn-4.26}
\lim_{\varepsilon \to 0}\tilde{\mathrm{III}}(\vez)= 0. 
\end{equation}

{\bf Step 5}. In this step, we estimate \( \tilde{\mathrm{I}}_1(\vez) \). To this end, we first decompose \( \tilde{\mathrm{I}}_1(\vez) \) as
\begin{align*}
\tilde{\mathrm{I}}_1(\vez)
&\lesssim \int_{\bigcup_{i \in \Lambda_2} \Qi} \left[ \int_{\bigcup_{k \in \Lambda_2}\Qkk} \frac{|D^\beta f(x) - D^\beta f(y)|^q |\sum_{j\in\Lambda_2}D^{\alpha - \beta} \psi_j(y)|^q}{|x - y|^{\sz q + n}}\,dy \right]^{\frac{p}{q}}\,dx \\
&\quad+ \int_{\bigcup_{i \in \Lambda_2} \Qi} \left[ \int_{\bigcup_{k \in \Lambda_2} \Qkk} \frac{|D^\beta f(x)|^q |\sum_{j\in\Lambda_2}D^{\alpha - \beta} \psi_j(x) - \sum_{j\in\Lambda_2}D^{\alpha - \beta} \psi_j(y)|^q}{|x - y|^{\sz q + n}}\,dy \right]^{\frac{p}{q}}\,dx \\
&=: \tilde{\mathrm{I}}_{1,1}(\vez) + \tilde{\mathrm{I}}_{1,2}(\vez).
\end{align*}
To estimate \( \tilde{\mathrm{I}}_{1,1}(\vez) \), by \eqref{eqn-xiutwo} and \eqref{eqn-2.6}, we conclude that
\begin{align}\label{eqn-4.27}
\tilde{\mathrm{I}}_{1,1}(\vez)
&\lesssim\int_{\bigcup_{i \in \Lambda_2} \Qi} \left[ \int_{\bigcup_{k \in \Lambda_2} \Qkk} \frac{|D^\beta f(x) - D^\beta f(y)|^q}{|x - y|^{\sz q + n}d(y,\pa\boz)^{|\az-\bz|q}}\,dy \right]^{\frac{p}{q}}\,dx\nonumber\\
&\leq \int_{\Omega_{3\varepsilon}} \left[ \int_{\Omega_{3\varepsilon}} \frac{|D^\beta f(x) - D^\beta f(y)|^q}{|x - y|^{\sz q + n}d(y,\pa\boz)^{(k-|\bz|)q}}\,dy \right]^{\frac{p}{q}}\,dx. 
\end{align}
Since \( f \in \huafweab \), it follows from \eqref{eqn-4.27} that $\tilde{\mathrm{I}}_{1,1}(\vez)\to 0$ as \( \varepsilon \to 0 \).

Next, we estimate \( \tilde{\mathrm{I}}_{1,2}(\vez) \). By \eqref{eqn-2.8}, we write
\begin{align*}
\tilde{\mathrm{I}}_{1,2}(\vez)
&\leq \sum_{i \in \Lambda_2} \int_{\Qi} \left[ \int_{\bigcup_{k \in \Lambda_2} \Qkk} \frac{|D^\beta f(x)|^q |\sum_{j\in\Lambda_2}D^{\alpha - \beta} \psi_j(x) - \sum_{j\in\Lambda_2}D^{\alpha - \beta} \psi_j(y)|^q}{|x - y|^{\sz q + n}}\,dy \right]^{\frac{p}{q}}\,dx  \nonumber \\
&\sim \sum_{i \in \Lambda_2} \int_{\Qi} \left[ \int_{\bigcup_{k \in \Lambda_{2,1}(i)} \Qkk} \frac{|D^\beta f(x)|^q |\sum_{j\in\Lambda_2}D^{\alpha - \beta} \psi_j(x) - \sum_{j\in\Lambda_2}D^{\alpha - \beta} \psi_j(y)|^q}{|x - y|^{\sz q + n}}\,dy \right]^{\frac{p}{q}}\,dx \nonumber \\
&\quad+ \sum_{i \in \Lambda_2} \int_{\Qi} \left[ \int_{\bigcup_{k \in \Lambda_{2,2}(i)} \Qkk} \frac{|D^\beta f(x)|^q |\sum_{j\in\Lambda_2}D^{\alpha - \beta} \psi_j(x) - \sum_{j\in\Lambda_2}D^{\alpha - \beta} \psi_j(y)|^q}{|x - y|^{\sz q + n}}\,dy \right]^{\frac{p}{q}}\,dx \nonumber \\
&=:\tilde{\mathrm{I}}_{1,2}^1(\vez) + \tilde{\mathrm{I}}_{1,2}^2(\vez).
\end{align*}
By the mean value theorem, \eqref{eqn-2.2}, Lemmas \ref{lem2.4} and \ref{lem2.1}, \eqref{eqn-2.8}, polar coordinates, \eqref{eqn-2.6} and 
the fact
that $\|\frac{D^\beta f}{d^{s-|\beta|}}\|_{L^p(\Omega)}<\infty$, we conclude that
\begin{align}\label{eqn-4.29}
\tilde{\mathrm{I}}_{1,2}^1(\vez)
&\leq \sum_{i \in \Lambda_2} \int_{\Qi} |D^\bz f(x)|^p \nonumber\\
&\quad\times\left[ \sum_{k \in \Lambda_{2,1}(i)} \int_{\Qkk} \frac{|\sum_{j \in \Lambda_{2,3}(i,k)} D^{\alpha - \beta}\psi_j(x) -\sum_{j \in \Lambda_{2,3}(i,k)} D^{\alpha - \beta}\psi_j(y)|^q}{|x - y|^{\sz q + n}}\,dy \right]^{\frac{p}{q}}\,dx \nonumber \\
&\lesssim \sum_{i \in \Lambda_2} \int_{\Qi} |D^\bz f(x)|^p \left[ \sum_{k \in \Lambda_{2,1}(i)} \int_{\Qkk} \sum_{j \in \Lambda_{2,3}(i,k)}{\| D^{\az-\bz+1} \psi_j \|_{L^\infty}^q |x - y|^{q(1 - \sz)-n}}\,dy \right]^{\frac{p}{q}}\,dx  \nonumber \\
&\lesssim \sum_{i \in \Lambda_2} \int_{\Qi} |D^\bz f(x)|^p \left[ \sum_{k \in \Lambda_{2,1}(i)} \int_{\Qkk} {{l_i}^{-q(|\az-\bz|+1)} |x - y|^{q(1 - \sz)-n}}\,dy \right]^{\frac{p}{q}}\,dx \nonumber \\
&\sim \sum_{i \in \Lambda_2} \int_{\Qi} |D^\bz f(x)|^p \left[ \sum_{k \in \Lambda_{2,1}(i)} \int_{\Qkk} {{d(x,\pa\boz)}^{-q(|\az-\bz|+1)} |x - y|^{q(1 - \sz)-n}}\,dy \right]^{\frac{p}{q}}\,dx \nonumber \\
&\lesssim \sum_{i \in \Lambda_2} \int_{\Qi} \frac{ |D^\bz f(x)|^p}{{d(x,\pa\boz)}^{(1+|\az-\bz|)p}} \left[ \int_0^{cl_i} \rho^{q(1 - \sz) - 1}\,d\rho \right]^{\frac{p}{q}}\,dx \nonumber \\
&\sim \sum_{i \in \Lambda_2} \int_{\Qi} \frac{ |D^\bz f(x)|^p}{{d(x,\pa\boz)}^{(1+|\az-\bz|)p}}  {l_i}^{p(1-\sz)}\,dx
\sim \int_{\bigcup_{i \in \Lambda_2} \Qi} \frac{|D^\bz f(x)|^p}{d(x, \partial\Omega)^{(\sz+|\az-\bz|) p}}\,dx  \nonumber \\
&\leq \int_{\bozthreee}\frac{|D^\bz f(x)|^p}{d(x, \partial\Omega)^{(s-|\bz|) p}}\,dx \to 0
\end{align}
as $\varepsilon \to 0.$ For $\tilde{\mathrm{I}}_{1,2}^2(\vez)$, we write
\begin{align*}
\tilde{\mathrm{I}}_{1,2}^2(\vez)
&\lesssim \sum_{i \in \Lambda_2} \int_{\Qi} \left[ \sum_{\substack{k\in\Lambda_{2,2}(i) \\ l_k\geq l_i}}\int_{\Qkk} \frac{|D^\beta f(x)|^q |\sum_{j\in\Lambda_2}D^{\alpha - \beta} \psi_j(x) - \sum_{j\in\Lambda_2}D^{\alpha - \beta} \psi_j(y)|^q}{|x - y|^{\sz q + n}}\,dy \right]^{\frac{p}{q}}\,dx \nonumber \\
&\quad+\sum_{i \in \Lambda_2} \int_{\Qi} \left[ \sum_{\substack{k\in\Lambda_{2,2}(i) \\ l_k< l_i}}\int_{\Qkk} \frac{|D^\beta f(x)|^q |\sum_{j\in\Lambda_2}D^{\alpha - \beta} \psi_j(x) - \sum_{j\in\Lambda_2}D^{\alpha - \beta} \psi_j(y)|^q}{|x - y|^{\sz q + n}}\,dy \right]^{\frac{p}{q}}\,dx \nonumber \\
&=:\tilde{\mathrm{I}}_{1,2}^{2,1}(\vez)+\tilde{\mathrm{I}}_{1,2}^{2,2}(\vez).
\end{align*} 
To estimate $\tilde{\mathrm{I}}_{1,2}^{2,1}(\vez)$, by \eqref{eqn-2.2}, \eqref{eqn-xiutwo}, Lemma \ref{lem2.1}, \eqref{eqn-8.13}, polar coordinates, \eqref{eqn-2.6} and 
$\|\frac{D^\beta f}{d^{s-|\beta|}}\|_{L^p(\Omega)}<\infty$, we find that
\begin{align}\label{eqn-4.30}
\tilde{\mathrm{I}}_{1,2}^{2,1}(\vez)
&\lesssim  \sum_{i \in \Lambda_2} \int_{\Qi} \left[ \sum_{k \in \Lambda_{2,2}(i)} \int_{\Qkk} \frac{| D^\bz f(x)|^q }{|x - y|^{\sz q + n}{l_i}^{|\az-\bz|q}}\,dy \right]^{\frac{p}{q}}\,dx \nonumber \\
&\lesssim \sum_{i \in \Lambda_2} \int_{\Qi} \left[ \sum_{k \in \Lambda_{2,2}(i)} \int_{\Qkk} \frac{|D^\bz f(x)|^q }{|x - y|^{\sz q + n}{d(x,\pa\boz)}^{|\az-\bz|p}}\,dy \right]^{\frac{p}{q}}\,dx \nonumber \\
&\lesssim \sum_{i \in \Lambda_2} \int_{\Qi} \frac{ |D^\bz f(x)|^p}{{d(x,\pa\boz)}^{|\az-\bz|p}} \left( \int_{\frac{l_i}{20}}^{\infty} \rho^{- \sz q - 1}\,d\rho \right)^{\frac{p}{q}}\,dx \nonumber \\
&\lesssim \sum_{i \in \Lambda_2} \int_{\Qi} \frac{ |D^\bz f(x)|^p}{{d(x,\pa\boz)}^{|\az-\bz|p}}  {l_i}^{-\sz p}\,dx\nonumber\\
&\sim \int_{\bigcup_{i \in \Lambda_2} \Qi} \frac{|D^\bz f(x)|^p}{d(x, \partial\Omega)^{(\sz+|\az-\bz|) p}}\,dx \nonumber \\
&\leq \int_{\bozthreee} \frac{|D^\bz f(x)|^p}{d(x, \partial\Omega)^{(s-|\bz|) p}}\,dx \to 0
\end{align}
as $\varepsilon \to 0.$ To estimate $\tilde{\mathrm{I}}_{1,2}^{2,2}(\vez)$, by \eqref{eqn-2.2}, \eqref{eqn-xiutwo} and Lemma \ref{lem2.1}(ii), we find that
\begin{align}\label{eqn-9.14}
\tilde{\mathrm{I}}_{1,2}^{2,2}(\vez)
&\lesssim \sum_{i \in \Lambda_2} \int_{\Qi} \left[ \sum_{\substack{k\in\Lambda_{2,2}(i) \\ l_k< l_i}}\int_{\Qkk} \frac{|D^\bz f(x)|^q }{|x - y|^{\sz q + n}d(y,\pa\boz)^{|\az-\bz|q}}\,dy \right]^{\frac{p}{q}}\,dx \nonumber\\
&\lesssim \sum_{i \in \Lambda_2} \int_{\Qi} \left[ \sum_{\substack{k\in\Lambda_{2,2}(i) \\ l_k< l_i}}\int_{\Qkk} \frac{| D^\bz f(x)-D^\bz f(y)|^q }{|x - y|^{\sz q + n}d(y,\pa\boz)^{|\az-\bz|q}} \,dy \right]^{\frac{p}{q}} \nonumber\\
&\quad+ \sum_{i \in \Lambda_2} \int_{\Qi} \left[ \sum_{\substack{k\in\Lambda_{2,2}(i) \\ l_k< l_i}}\int_{\Qkk}\frac{|D^\bz f(y)|^q }{|x - y|^{\sz q + n}d(y,\pa\boz)^{|\az-\bz|q}}\,dy \right]^{\frac{p}{q}}\,dx. 
\end{align}
By \( f \in \huafweab \) and an argument similar to that used in the estimation of \eqref{eqn-4.27}, we find that the first term on the right-hand side of \eqref{eqn-9.14} tends to 0 as $\vez\to 0$, and by an argument similar to that used in the estimation of  $\tilde{\mathrm{III}}(\vez)$, we find that the second term on the right-hand side of \eqref{eqn-9.14} tends to 0 as $\vez\to 0$. Hence,
\begin{equation}\label{eqn-4.31}
\lim_{\varepsilon \to 0}\tilde{\mathrm{I}}_{1,2}^{2,2}(\vez)= 0. 
\end{equation}
Combining \eqref{eqn-4.27}--\eqref{eqn-4.31}, we obtain
\begin{equation}\label{eqn-4.32}
\lim_{\varepsilon \to 0}\tilde{\mathrm{I}}_{1}(\vez)= 0.
\end{equation}

{\bf Step 6}. In this step, we estimate \( \tilde{\mathrm{II}}(\vez) \). To this end, we write
\begin{align}\label{1717}
\tilde{\mathrm{II}}(\vez)
&\leq \int_{\Omega_{14\varepsilon}} \left[ \int_{\bigcup_{k \in \Lambda_2} \Qk} \frac{|D^\beta f(x) \sum_{j\in\Lambda_2}D^{\alpha - \beta} \psi_j(x) - D^\beta f(y) \sum_{j\in\Lambda_2}D^{\alpha - \beta} \psi_j(y)|^q}{|x - y|^{\sz q + n}}\,dy \right]^{\frac{p}{q}}\,dx  \nonumber\\
&\lesssim \int_{\Omega_{14\varepsilon}} \left[ \int_{\bigcup_{k \in \Lambda_2} \Qk} \frac{|D^\beta f(x)|^q |\sum_{j\in\Lambda_2}D^{\alpha - \beta} \psi_j(x) - \sum_{j\in\Lambda_2}D^{\alpha - \beta} \psi_j(y)|^q}{|x - y|^{\sz q + n}}\,dy \right]^{\frac{p}{q}}\,dx \nonumber\\
&\quad+ \int_{\Omega_{14\varepsilon}} \left[ \int_{\bigcup_{k \in \Lambda_2} \Qk} \frac{|D^\beta f(x) - D^\beta f(y)|^q |\sum_{j\in\Lambda_2}D^{\alpha - \beta} \psi_j(y)|^q}{|x - y|^{\sz q + n}}\,dy \right]^{\frac{p}{q}}\,dx.
\end{align}
By an argument similar to that used in the estimation of \( \tilde{\mathrm{I}}_{1,1}(\vez) \), we find that the first term on the right-hand side of \eqref{1717} tends to 0 as $\vez\to 0$, and by an argument similar to that used in the estimation of \( \tilde{\mathrm{I}}_{1,2}(\vez) \), we find that the second term on the right-hand side of \eqref{1717} tends to 0 as $\vez\to 0$. Hence,
\begin{equation}\label{eqn-4.33}
\lim_{\varepsilon \to 0} \tilde{\mathrm{II}}(\vez) = 0.
\end{equation}
Combining \eqref{eqn-8.5}, \eqref{eqn-4.23}, \eqref{eqn-4.25}, \eqref{eqn-4.26}, \eqref{eqn-4.32} and \eqref{eqn-4.33}, we obtain
\[
\lim_{\huae \to 0} \adk \left| D^{\az} \we \right|_{\cf_{p,q}^{\sz}(\boz)}=0,
\]
which proves \eqref{eqn-4.21}.

{\bf Step 7}. In this step, we prove \eqref{eqn-4.22}. From an argument similar to that used in the estimation in Step 3, we deduce that
\begin{align*}
& \sum_{\substack{\bz\in\nno^n \\ | \bz| \leq k}} \left|D^{\bz}\we\right|_{\cf_{p,q}^{\sz}(\boz,d^{-(k-|\bz|)q})}^p \\ &\quad\lesssim \sum_{\substack{\bz\in\nno^n \\|\bz| \leq k}}  \sum_{\substack{\gamma\in\nno^n \\ \gamma \leq \bz}} \binom{\bz}{\gamma} \int_{\boz} \left[\int_{\boz} \frac{| D^\gamma f(x) \sum_{j\in\Lambda_2}D^{\bz - \gamma} \psi_j(x) - D^\gamma f(y) \sum_{j\in\Lambda_2}D^{\bz - \gamma} \psi_j(y)|^q}{|x - y|^{\sz q + n}d(y,\pa\boz)^{(k-|\bz|)q}}\,dy \right]^{\frac{p}{q}}\,dx.
\end{align*}
Similar to the decomposition in \eqref{eqn-4.23}, we obtain
\begin{align*}
&\int_{\boz} \left[ \int_{\boz} \frac{| D^\gamma f(x) \sum_{j\in\Lambda_2}D^{\bz - \gamma} \psi_j(x) - D^\gamma f(y) \sum_{j\in\Lambda_2}D^{\bz - \gamma} \psi_j(y)|^q}{|x - y|^{\sz q + n}d(y,\pa\boz)^{(k-|\bz|)q}}\,dy \right]^{\frac{p}{q}}\,dx \\
&\quad\lesssim \int_{\bigcup_{i \in \blz_2} \Qi} \\
&\qquad\times\left[ \int_{\bigcup_{k \in \blz_2} \Qkk} \frac{| D^\gamma f(x) \sum_{j\in\Lambda_2}D^{\bz - \gamma} \psi_j(x) - D^\gamma f(y)\sum_{j\in\Lambda_2} D^{\bz - \gamma} \psi_j(y)|^q}{|x - y|^{\sz q + n}d(y,\pa\boz)^{ (k-|\bz|)q}}\,dy \right]^{\frac{p}{q}}\,dx \\
&\qquad+ \int_{\bigcup_{i \in \blz_2} \Qi} \left[ \int_{\boz \setminus \bigcup_{k \in \blz_2} \Qkk} \frac{| D^\gamma f(x) \sum_{j\in\Lambda_2}D^{\bz - \gamma} \psi_j(x)|^q}{|x - y|^{\sz q + n}d(y,\pa\boz)^{(k-|\bz|)q}}\,dy \right]^{\frac{p}{q}}\,dx \\
&\qquad+\int_{\bozfour \setminus \bigcup_{i \in \blz_2} \Qi} \\
&\qquad\times\left[ \int_{\bigcup_{k \in \blz_2} \Qk} \frac{| D^\gamma f(x) \sum_{j\in\Lambda_2}D^{\bz - \gamma} \psi_j(x) - D^\gamma f(y) \sum_{j\in\Lambda_2}D^{\bz - \gamma} \psi_j(y)|^q}{|x - y|^{\sz q + n}d(y,\pa\boz)^{(k-|\bz|)q}}\,dy \right]^{\frac{p}{q}}\,dx \\
&\qquad+ \int_{\boz \setminus \bozfour} \left[ \int_{\bigcup_{k \in \blz_2} \Qk} \frac{| D^\gamma f(y) \sum_{j\in\Lambda_2}D^{\bz - \gamma} \psi_j(y)|^q}{|x - y|^{\sz q + n}d(y,\pa\boz)^{(k-|\bz|)q}}\,dy \right]^{\frac{p}{q}}\,dx. 
\end{align*}
Using this and following the same proof method as for \eqref{eqn-4.21}, \eqref{eqn-4.22} can be obtained via straightforward calculations by simply replacing $\bz$ with $\gamma$ in the proof process; the details are omitted. This completes the proof of part (i) of Theorem \ref{thm-secmain}.

The proof of  (ii) is similar to the corresponding part 
of Theorem \ref{thm-main}, the details being omitted.
This completes the proof of Theorem \ref{thm-secmain}.
\end{proof}

\section{The case $s\in\mathbb{N}$}\label{sN}

In this section, we show the density property for Triebel--Lizorkin 
spaces on $\boz$ in integer-order smoothness 
case by proving Theorem \ref{thm-third}. To this end,
we need the following three technical lemmas.
Throughout this section, we always fix the smoothness index
$s=k\in\mathbb{N}$.

\begin{lemma}\label{lembls}
There exists a constant $C\in(0,\infty)$ such that,
for any $x\in\boz$ and $h\in\mathbb{R}^n$ satisfying
\begin{align*}
|h|<\frac{d(x,\pa\boz)}{2(k+1)}
\ \ \text{and}\ \ 
x+\nu h\in\Omega_{3\varepsilon}
\end{align*}
for some $\nu\in\mathbb{N}_0\cap[0,k+1]$,
it holds that $x\in\Omega_{C\varepsilon}$.
\end{lemma}

\begin{proof}
If $z:=x+\nu h\in\boz_{3\vez}$ for some $0\le\nu\le k+1$, then
the triangle inequality implies that
\begin{align*}
d(x,\pa\boz)
\leq d(z,\pa\boz)+|z-x|
\leq 3\vez+(k+1)|h|
< 3\vez+\frac{1}{2} d(x,\pa\boz).
\end{align*}
The last inequality implies $d(x,\pa\boz)\lesssim\vez$. Hence $x\in\boz_{C\vez}$.
\end{proof}

\begin{lemma}[\cite{g2014}, Proposition~1.4.5]\label{lemdwc}
Let $f$ be a ${C}^b$ function on $\boz$ for some $b\in\nn$. Then for any $x\in\boz$ and $h\in\rn$ satisfying $b|h|<d(x,\pa\boz)$, the following identity holds:
\begin{align*}
\Delta_h^{b}(f)(x)&=
\int_{0}^{1} \cdots \int_{0}^{1} \sum_{j_1=1}^{n} \cdots \sum_{j_b=1}^{n} h_{j_1} \cdots h_{j_b} \\
&\quad \times \left(\partial_{j_1} \cdots \partial_{j_b} f\right)\big(x+(s_1+\cdots+s_b)h\big)
\,ds_1\cdots ds_b.
\end{align*}
\end{lemma}

\begin{lemma}[\cite{j1965}, p.\,95, (6)]\label{lemdlf}
Let $f$ and $g$ be two functions defined on $\boz$. Then, for any $m\in\nn$, $x\in\boz$ and $h\in\rn$ satisfying $m|h|<d(x,\pa\boz)$,
\begin{align*}
\Delta_h^m(fg)(x)
=\sum_{b=0}^m \binom{m}{b}
\Delta_h^{m-b}f(x)
\Delta_h^b g\left(x+(m-b)h\right).
\end{align*}
\end{lemma}

Now, we are ready to prove Theorem \ref{thm-third}.

\begin{proof}[Proof of Theorem \ref{thm-third}]
The inclusion
${\mathfrak F}^{k,1}_{p,q}(\boz)
\subset{\mathfrak F}^{k,2}_{p,q}(\boz)$
is from the fact that $C_{\mathrm c}^\infty(\Omega)\subset C^\infty(\boz)$. 

Next, we prove the reverse inclusion.
Let
$f\in C^\infty(\boz)\cap\mathfrak F^{k}_{p,q}(\boz)$,
$\{Q_j\}_{j=1}^\infty$ and $\{\psi_j\}_{j=1}^\infty\subset C_{\mathrm c}^\infty(\Omega)$ be the Whitney decomposition and its associated partition of unity fixed in Section~\ref{s2}. For any given \( \varepsilon\in(0,\infty)\), by \eqref{eqn-2.1} and \eqref{eqn-2.3}, we divide $f$ by, for any $x\in\Omega$,
\begin{align*}
f(x)= \sum_{j \in \Lambda_1} \psi_j(x) f(x) + \sum_{j \in \Lambda_2} \psi_j(x) f(x)
=: v_{\varepsilon}(x) + w_{\varepsilon}(x).
\end{align*}
Similar to the argument used in the proof of Theorem \ref{thm-main}, it suffices to  prove that
\begin{align}\label{eqn-integerw}
\lim_{\vez\to0}\|w_\vez\|_{\mathfrak F^{k}_{p,q}(\boz)}=0,
\end{align}
which, combined with \eqref{diffnorm}, is equivalent to verifying the following three identities:
\begin{align}\label{eqn5.2}
\lim_{\vez\to0}\|w_\vez\|_{L^p(\boz)}=0,
\end{align}
\begin{align}\label{eqn5.3}
\lim_{\vez\to0}\left\|\frac{w_\vez}{d^k}\right\|_{L^p(\boz)}=0
\end{align}
and
\begin{align}\label{eqn5.4}
\lim_{\vez\to0}\sum_{r=0}^k [w_\vez]_{r,k}=0.
\end{align}
We divide the proofs of \eqref{eqn5.2}--\eqref{eqn5.4} into three steps.

{\bf Step 1}. In this step, we show \eqref{eqn5.2}.
By \eqref{eqn-2.6} and $|w_\vez|\le |f|$, we obtain
\begin{align*}
\lim_{\vez\to0}\|w_\vez\|_{L^p(\boz)}^p
\le
\lim_{\vez\to0}\int_{\boz_{3\vez}} |f(x)|^p\,dx
=0.
\end{align*}

{\bf Step 2}. In this step, we show \eqref{eqn5.3}. Let $\theta_\vez:=\sum_{j \in \Lambda_2} \psi_j$.
By \eqref{eqn-2.6}, $|\theta_\vez|\le1$ and $f/d^k\in L^p(\boz)$, we obtain
\begin{align*}
\left\|\frac{w_\vez}{d^k}\right\|_{L^p(\boz)}^p
=
\int_{\boz}\frac{|f(x)\theta_\vez(x)|^p}{d(x,\pa\boz)^{kp}}\,dx
\le
\int_{\boz_{3\vez}}\frac{|f(x)|^p}{d(x,\pa\boz)^{kp}}\,dx,
\end{align*}
which tends to $0$ as $\vez\to0$.

{\bf Step 3}. In this step, we show \eqref{eqn5.4}.
Fix $r\in\{0,1,\dots,k\}$ and let
$m:=k+1-r$.
By Lemma \ref{lemdlf}, we find that, for any $x\in\boz$ and $h\in\rn$ satisfying $|h|<\frac{d(x,\pa\boz)}{2(k+1)}$,
\begin{align*}
\Delta_h^m(f\theta_\vez)(x)
=
\sum_{b=0}^m \binom{m}{b}
\Delta_h^{m-b}f(x)
\Delta_h^b\theta_\vez(x+(m-b)h).
\end{align*}
Hence, we have
\begin{align*}
[w_\vez]_{r,k}
\lesssim
\sum_{b=0}^m T_{r,b}(\vez),
\end{align*}
where, for any $b\in\{0,\ldots,m\}$,
\begin{align*}
T_{r,b}(\vez)^p
:=
\int_\boz
\left[
\int_{|h|<\frac{d(x,\pa\boz)}{2(k+1)}}
\frac{
|\Delta_h^{m-b}f(x)|^q
|\Delta_h^b\theta_\vez(x+(m-b)h)|^q
}
{|h|^{n+(k-r)q}d(x,\pa\boz)^{rq}}
\,dh
\right]^\frac{p}{q}
\,dx.
\end{align*}

We first consider the terms for which $0\le b\le m-1$.  Since $m=k+1-r$, this means $b\le k-r$.  By Lemma \ref{lemdwc}, the support condition of $\psi_j$ for any $j\in\Lambda_2$,
Lemma~\ref{lem2.1} and \eqref{eqn-2.2}, we obtain, for any $x\in\boz$ and $h\in\rn$ satisfying $|h|<\frac{d(x,\pa\boz)}{2(k+1)}$,
\begin{align}\label{deltatheta}
&\left|\Delta_h^b\theta_\vez(x+(m-b)h)\right|\nonumber\\
&\quad\leq bn|h|^b\sup_{\delta\in(0,1)}\sup_{\beta\in\mathbb{N}_0^n,\,|\beta|=b}
\left|D^\beta \theta_\varepsilon(x+(m-b)h+\delta bh)\right|\nonumber\\
&\quad\lesssim|h|^b\sup_{z\in B(x,m|h|)}\sup_{\beta\in\mathbb{N}_0^n,\,|\beta|=b}
\left|D^\beta \theta_\varepsilon(z)\right|\nonumber\\
&\quad\leq|h|^b\sup_{z\in B(x,\frac{1}{2}d(x,\pa\boz))}\sup_{\beta\in\mathbb{N}_0^n,\,|\beta|=b}
\left|D^\beta \theta_\varepsilon(z)\right|\nonumber\\
&\quad\leq|h|^b\sup_{z\in B(x,\frac{1}{2}d(x,\pa\boz))}\sup_{\beta\in\mathbb{N}_0^n,\,|\beta|=b}\sum_{j\in\Lambda_2,\,\mathrm{supp\,}\psi_j\cap B(x,\frac{1}{2}d(x,\pa\boz))\neq\emptyset}
\left|D^\beta\psi_j(z)\right|\nonumber\\
&\quad\lesssim|h|^b\sup_{z\in B(x,\frac{1}{2}d(x,\pa\boz))}\sup_{\beta\in\mathbb{N}_0^n,\,|\beta|=b}\sup_{j\in\Lambda_2,\,\mathrm{supp\,}\psi_j\cap B(x,\frac{1}{2}d(x,\pa\boz))\neq\emptyset}
\left|D^\beta\psi_j(z)\right|\nonumber\\
&\quad\lesssim|h|^b\sup_{z\in B(x,\frac{1}{2}d(x,\pa\boz))}\sup_{\beta\in\mathbb{N}_0^n,\,|\beta|=b}\sup_{j\in\Lambda_2,\,\mathrm{supp\,}\psi_j\cap B(x,\frac{1}{2}d(x,\pa\boz))\neq\emptyset}
\frac{1}{(\mathrm{diam\,}Q_j)^b}\nonumber\\
&\quad\sim\frac{|h|^b}{d(x,\pa\boz)^b}.
\end{align}
%On the other hand, since $0\leq\theta_\varepsilon\leq1$,
%it follows from \eqref{deltabh} that %$\left|\Delta_h^b\theta_\vez(x+(m-b)h)\right|\lesssim1$.
%This, together with \eqref{deltatheta},
%implies that
%\begin{align}\label{deltatheta1}
%\left|\Delta_h^b\theta_\vez(x+(m-b)h)\right|\lesssim\min\left\{1,\,\frac{|h|^%b}{d(x)^b}\right\}.
%\end{align}
Now observe that \eqref{eqn-2.6} implies $\supp\theta_\vez\subset\boz_{3\varepsilon}$, hence by Lemma \ref{lembls}, we find that $x\in\boz_{C\vez}$. Combining this and \eqref{deltatheta}, further implies that
\begin{align*}
T_{r,b}(\vez)^p
\lesssim
\int_{\boz_{C\vez}}
\left[
\int_{|h|<\frac{d(x,\pa\boz)}{2(k+1)}}
\frac{|\Delta_h^{k+1-(r+b)}f(x)|^q}
{|h|^{n+(k-(r+b))q}d(x,\pa\boz)^{(r+b)q}}
\,d h
\right]^\frac{p}{q}
\,d x.
\end{align*}
Because $[f]_{r+b,k}<\infty$, the last integral tends to $0$ as $\vez\to0$.
Thus, for any $0\le b\le m-1$,
\begin{align*}
\lim_{\vez\to0}T_{r,b}(\vez)	=0.
\end{align*}

It remains to treat the pure cutoff term $b=m$.  Since $m=k+1-r$, by an argument similar to that used in the estimation of \eqref{deltatheta}
with $b$ and $x+(m-b)h$ replaced, respectively, by $m$ and $x$, we obtain, for any $x\in\boz$ and $h\in\rn$ satisfying $|h|<\frac{d(x,\pa\boz)}{2(k+1)}$,
\begin{align*}
\frac{|f(x)|^q|\Delta_h^m\theta_\vez(x)|^q}
{|h|^{n+(k-r)q}d(x,\pa\boz)^{rq}}
\lesssim
\frac{|f(x)|^q |h|^{mq}}
{|h|^{n+(k-r)q}d(x,\pa\boz)^{(r+m)q}}
=\frac{|f(x)|^q |h|^q}{|h|^n d(x,\pa\boz)^{(k+1)q}}.
\end{align*}
Combining this and polar coordinates, further implies that, for any $x\in\boz$,
\begin{align*}
\int_{|h|<\frac{d(x,\pa\boz)}{2(k+1)}}
\frac{|f(x)|^q |h|^q}{|h|^n d(x,\pa\boz)^{(k+1)q}}\,dh
\lesssim
\frac{|f(x)|^q}{d(x,\pa\boz)^{kq}}.
\end{align*}
Again this term is supported in a boundary layer.  Consequently,
\begin{align*}
T_{r,m}(\vez)^p
\lesssim
\int_{\boz_{C\vez}}\frac{|f(x)|^p}{d(x,\pa\boz)^{kp}}\,dx
\to0
\end{align*}
as $\vez\to0$ because $f/d^k\in L^p(\boz)$.

Combining the estimates for $b\in\{0,1,\dots,m\}$, further implies that
\begin{align*}
\lim_{\vez\to0}[w_\vez]_{r,k}=0.
\end{align*}
Since $r\in\{0,1,\dots,k\}$ was arbitrary and there are only finitely many $r$, \eqref{eqn5.4} follows.

Combining \eqref{eqn5.2}, \eqref{eqn5.3} and \eqref{eqn5.4},
we then conclude that \eqref{eqn-integerw} holds. This completes the proof of Theorem \ref{thm-third}.
\end{proof}

\section{Stability under the restriction norm}\label{s5}

In this section, we show that the density property is stable if the intrinsic norm is replaced by 
the restriction norm defined in \eqref{eqn-wrctls}. To this end, let us recall the 
definition of linear extension operators.
Let $\boz\subset\mathbb{R}^n$ be a bounded uniform domain, $\mathcal{W}_1$ the Whitney
decomposition of $\Omega$, $\mathcal{W}_2$ the Whitney
decomposition of $\overline{\Omega}^\complement$, and 
$\mathcal{W}_3$ the collection of cubes in $\mathcal{W}_2$ whose side lengths are sufficiently small
(depending only on \(\Omega\)). By [\cite{j1981}, Lemma~2.4] (see also [\cite{ps2017}, Section~3]), for every $Q\in\mathcal W_3$ one may choose a cube $Q^*\in\mathcal W_1$ such that
\begin{align}\label{eqn-7.1}
D(Q,Q^*) \lesssim l(Q) \quad\text{and}\quad l(Q)\sim l(Q^*).
\end{align}
We call any such chosen $Q^*$ a \emph{symmetric cube} of $Q$; the choice need not be unique, and we fix one choice for each $Q$ before invoking the following lemma.

\begin{lemma}[\cite{ps2017}, Lemma 3.9]\label{lem5.1}
Let $Q_1$, $Q_2 \in \mathcal{W}_3$ and $S \in \mathcal{W}_1$. Then
\begin{enumerate}
\item[\rm(i)]there exists a positive constant $C$ such that $\#\{Q \in \mathcal{W}_3: Q^* = S\} \leq C$;
\item[\rm(ii)]
$D(Q_1^{*}, Q_2^{*}) \sim D(Q_1, Q_2)$ and $D(Q_1^{*}, S)\sim D(Q_1, S)$; 
\item[\rm(iii)]if $\overline{Q_1} \cap \overline{Q_2} \neq \varnothing$, then
$
D(Q_1^{*}, Q_2^{*}) \sim l(Q_1).
$
\end{enumerate}
\end{lemma}

With the family of bump functions $\{\psi_Q\}_{Q\in\mathcal{W}_2}$ defined as immediately after Lemma \ref{lem2.1}, simply by replacing $Q_j$ with $Q$, we define the extension operator $\Lambda_0$ by setting, for any $f \in L^1_{\text{loc}}(\Omega)$ and $x\in\rn$,
\begin{equation}\label{eqn-3.1}
\Lambda_0 f(x) := f(x)\chi_\Omega(x) + \sum_{Q\in\mathcal{W}_3} \psi_Q(x) f_{Q^{*}}.
\end{equation}
In particular, it follows from $\psi_Q\subset\frac{11}{10}Q$, Lemma \ref{lem2.1} and \eqref{eqn-3.1} that,
for any $x\in\pa\boz$, $\Lambda_0 f(x)=0$.

\begin{lemma}[\cite{ps2017}, Lemma 3.10]\label{L6.2}
Let $p,q \in [1, \infty)$ and $s\in(\max\{0,\frac{n}{p} - \frac{n}{q}\},1)$. Assume that $\boz$ is a uniform domain. Then there exists a bounded linear extension operator $\Lambda_0\colon \huaf \to F_{p,q}^{s}(\rn)$.
\end{lemma}

\begin{remark}
The result [\cite{ps2017}, Lemma 3.10] is stated only for $p,q>1$. The cases $p=1$ or $q=1$,
which are not covered there, also satisfy the conclusion of Lemma~\ref{L6.2}.
Indeed, the proof is exactly the same as that of [\cite{ps2017}, Lemma 3.10] and can be repeated verbatim.
\end{remark}
Now, we are ready to show Proposition \ref{pro-main}.
\begin{proof}[Proof of Proposition \ref{pro-main}]
Let $f\in\cf_{p,q}^{s}(\boz,d^{-aq})$ and $\Lambda_0 f$ be as in \eqref{eqn-3.1}. We only need to show
\begin{align}\label{eqn-6.2}
\left\| \Lambda_0 f \right\|_{F_{p,q}^{s}(\rn,d^{-aq})}
&\lesssim \| f \|_{\cf_{p,q}^{s}(\boz,d^{-aq})}.
\end{align}
To this end, from the proof of [\cite{ps2017}, Lemma~3.10]
(we need the assumption that $\Omega$ is a uniform domain here),
we deduce that
\begin{align*}
\| \Lambda_0 f \|_{F_{p,q}^{s}(\rn,d^{-aq})}
&= \| \Lambda_0 f \|_{\cf_{p,q}^{s}(\rn)} +\left\{\int_{\rn} \left[ \int_{\rn} \frac{|\Lambda_0 f(x) - \Lambda_0 f(y)|^q}{|x - y|^{s q + n}d(y,\pa\boz)^{aq}}\,dy \right]^{\frac{p}{q}}\,dx \right\}^{\frac{1}{p}} \\
&\lesssim\|f\|_{\cf_{p,q}^{s}(\Omega)} +\left\{\int_{\rn} \left[ \int_{\rn} \frac{|\Lambda_0 f(x) - \Lambda_0 f(y)|^q}{|x - y|^{s q + n}d(y,\pa\boz)^{aq}}\,dy \right]^{\frac{p}{q}}\,dx \right\}^{\frac{1}{p}},
\end{align*}
which, combined with the fact that
\begin{align}\label{1544}
&\int_{\boz} \left[ \int_{\boz} \frac{|f(x)-f(y)|^q}{|x - y|^{s q + n}}\,dy \right]^{\frac{p}{q}}\,dx\nonumber\\
&\quad\leq
(\mathrm{diam\,}\Omega)^{ap}\int_{\boz} \left[ \int_{\boz} \frac{|f(x) - f(y)|^q}{|x - y|^{s q + n}d(y,\pa\boz)^{aq}}\,dy \right]^{\frac{p}{q}}\,dx
\end{align}
(since $\Omega$ is bounded),
further implies that
\begin{align*}
\| \Lambda_0 f \|_{F_{p,q}^{s}(\rn,d^{-aq})}
&\lesssim\|f\|_{\cf_{p,q}^{s}(\boz,d^{-aq})}\\
&\quad+\left\{\int_{\rn} \left[ \int_{\rn} \frac{|\Lambda_0 f(x) - \Lambda_0 f(y)|^q}{|x - y|^{s q + n}d(y,\pa\boz)^{aq}}\,dy \right]^{\frac{p}{q}}\,dx \right\}^{\frac{1}{p}}.
\end{align*}
Therefore, to prove \eqref{eqn-6.2}, we only need to prove
\begin{align}\label{eqn-aim}
\int_{\rn} \left[ \int_{\rn} \frac{|\Lambda_0 f(x) - \Lambda_0 f(y)|^q}{|x - y|^{s q + n}d(y,\pa\boz)^{aq}}\,dy \right]^{\frac{p}{q}} \,dx \lesssim \| f \|_{\cf_{p,q}^{s}(\boz,d^{-aq})}^p.
\end{align}
More precisely, we will prove that
\[
X+Y+Z \lesssim \| f \|_{\cf_{p,q}^{s}(\boz,d^{-aq})}^{p},
\]
where 
\begin{align*}
&X:= \int_{\boz} \left[ \int_{{\boz}^\complement} \frac{|\Lambda_0 f(x) - \Lambda_0 f(y)|^q}{|x - y|^{s q + n}d(y,\pa\boz)^{aq}}\,dy \right]^{\frac{p}{q}}\,dx,  \\
&Y:= \int_{{\boz}^\complement} \left[ \int_{\boz} \frac{|\Lambda_0 f(x) - \Lambda_0 f(y)|^q}{|x - y|^{s q + n}d(y,\pa\boz)^{aq}}\,dy \right]^{\frac{p}{q}}\,dx,
\end{align*}
and
\begin{align*}
Z:= \int_{{\boz}^\complement} \left[ \int_{{\boz}^\complement} \frac{|\Lambda_0 f(x) - \Lambda_0 f(y)|^q}{|x - y|^{s q + n}d(y,\pa\boz)^{aq}}\,dy \right]^{\frac{p}{q}}\,dx.
\end{align*}

We first estimate $X$.
By \eqref{eqn-3.1}, we write
\begin{align}\label{eqn-5.3}
X&= \int_{\boz} \left[ \int_{{\boz}^\complement} \frac{| f(x) - \sum_{S\in\mathcal{W}_3} \psi_S(y) f_{S^{*}}|^q}{|x - y|^{s q + n}d(y,\pa\boz)^{aq}}\,dy \right]^{\frac{p}{q}}\,dx  \nonumber \\
&\lesssim \sum_{Q\in\mathcal{W}_1}  \int_{Q} \left[ \int_{{\boz}^\complement} \frac{| f(x) |^q |1- \sum_{S\in\mathcal{W}_3} \psi_S(y) |^q}{|x - y|^{s q + n}d(y,\pa\boz)^{aq}}\,dy \right]^{\frac{p}{q}}\,dx \nonumber\\
&\quad+ \sum_{Q\in\cw_1} \int_{Q} \left[ \int_{{\boz}^\complement} \frac{| \sum_{S\in\mathcal{W}_3} \psi_S(y)|^q |f(x)- f_{S^{*}}|^q}{|x - y|^{s q + n}d(y,\pa\boz)^{aq}}\,dy \right]^{\frac{p}{q}}\,dx \nonumber \\
&=:X_1+X_2.
\end{align}
Let 
$$
\cw_4:=\left\{S\in\cw_3:\text {all the neighbors of S are in } \cw_3 \right\}.
$$
Note that, for any $y\in P\in\cw_4$, we have that $1-\sum_{S\in\mathcal{W}_3}\psi_{S}(y) \equiv 0$  and, otherwise 
\begin{align}\label{1719}
0\leq 1-\sum_{S\in\mathcal{W}_3}\psi_{S}(y)\leq 1.
\end{align}
By this and Lemma \ref{lem2.1}, we obtain
\begin{align}\label{eqn-6.6}
X_1&\leq\sum_{Q\in\mathcal{W}_1}  \int_{Q} \left[ \sum_{P\in\cw_2\setminus\cw_4} \int_{P}\frac{| f(x) |^q}{|x - y|^{s q + n}d(y,\pa\boz)^{aq}}\,dy \right]^{\frac{p}{q}}\,dx \nonumber\\
&\lesssim \sum_{Q\in\mathcal{W}_1}  \int_{Q} \left[\sum_{P\in\cw_2\setminus\cw_4} \int_{P} \frac{| f(x) |^q}{D(Q,P)^{s q + n}l(P)^{aq}}\,dy \right]^{\frac{p}{q}}\,dx \nonumber\\
&=\sum_{Q\in\mathcal{W}_1}  \int_{Q} \left[\sum_{P\in\cw_2\setminus\cw_4} \frac{l(P)^{n-aq}}{D(Q,P)^{s q + n}} \right]^{\frac{p}{q}}|f(x)|^p\,dx \nonumber\\
&\leq \left[\sum_{P\in\cw_2\setminus\cw_4} \frac{l(P)^{n-aq}}{d(P,\pa\boz)^{s q + n}} \right]^{\frac{p}{q}}\|f\|_{L^{p}(\boz)}^p\nonumber\\
&\sim \left[\sum_{P\in\cw_2\setminus\cw_4} \frac{1}{l(P)^{(s+a)q}} \right]^{\frac{p}{q}}\|f\|_{L^{p}(\boz)}^p.
\end{align}
Note that for any $P\in\cw_2\setminus\cw_4$, we have $l(P)>l_0$,
where $l_0$ is a positive constant depending only on $\boz$. Let \(m\in\mathbb N_0\) and
\begin{align*}
N_m:=\#\left\{P\in\cw_2\setminus\cw_4:
2^ml_0<l(P)\leq2^{m+1}l_0\right\}.
\end{align*}
From this and the assumption that \(\Omega\) is bounded, it follows that the
numbers \(N_m\) are uniformly bounded by a constant \(N\) depending only
on \(\Omega\) and the dimension \(n\). Hence,
\begin{align}\label{eqn-3.3}
\left[\sum_{P\in\cw_2\setminus\cw_4} \frac{1}{l(P)^{(s+a)q}} \right]^{\frac{p}{q}}
\lesssim \left[N\sum_{m=0}^{\infty}(2^m l_0)^{-(s+a)q}\right]^{\frac pq}
\lesssim l_0^{-(s+a)p}.
\end{align}
By this and \eqref{eqn-6.6}, we conclude that
\begin{equation}\label{eqn-3.4}
X_1\lesssim l_0^{-(s+a)p}\|f\|_{L^{p}(\boz)}^p\lesssim\|f\|_{L^{p}(\boz)}^p.
\end{equation}
To estimate $X_2$, by $\text{supp}\,\psi_S \subset \frac{11}{10}S$ for any $S\in\mathcal{W}_3$ and Lemma \ref{lem2.1}, we obtain
\begin{align}\label{1518}
X_2&\lesssim \sum_{Q\in\cw_1} \int_{Q} \left[ \sum_{S\in\mathcal{W}_3}\int_{\frac{11}{10}S} \frac{ |\psi_S(y)|^q |f(x)- f_{S^{*}}|^q}{|x - y|^{s q + n}d(y,\pa\boz)^{aq}}\,dy \right]^{\frac{p}{q}}\,dx\nonumber\\
&\lesssim \sum_{Q\in\cw_1} \int_{Q} \left[ \sum_{S\in\mathcal{W}_3} \frac{|f(x)- f_{S^{*}}|^q}{D(Q,S)^{s q + n}l(S)^{aq}} \int_{\frac{11}{10}S} |\psi_S(y)|^q\,dy \right]^{\frac{p}{q}}\,dx.
\end{align}
By the fact that $0\leq\psi_S\leq1$,
\eqref{eqn-7.1} and H\"older's inequality, we find that for any $S\in\mathcal{W}_3$,
$$
|f(x)- f_{S^{*}}|^q\int_{\frac{11}{10}S} |\psi_S(y)|^q\,dy\lesssim
|f(x)- f_{S^{*}}|^q l(S^{*})^n
\leq\int_{S^{*}} |f(x)-f(\xi)|^q\,d\xi,
$$
which, together with
\eqref{1518} and Lemmas \ref{lem2.1} and \ref{lem5.1}, further implies that
\begin{align}\label{eqn-3.5}
X_2&\lesssim \sum_{Q\in\cw_1} \int_{Q} \left[ \sum_{S\in\mathcal{W}_3} \frac{\int_{S^{*}} |f(x)-f(\xi)|^q\,d\xi}{D(Q,\sxing)^{s q + n}l(\sxing)^{aq}} \right]^{\frac{p}{q}}\,dx \nonumber \\
&\lesssim \sum_{Q\in\cw_1} \int_{Q} \left[ \sum_{S\in\mathcal{W}_3} \int_{S^{*}} \frac{|f(x)-f(\xi)|^q}{|x - \xi|^{s q + n}d(\xi,\pa\boz)^{aq}}\,d\xi \right]^{\frac{p}{q}}\,dx \nonumber  \\
&\leq \int_{\boz} \left[ \int_{\boz} \frac{| f(x) - f(\xi)|^q}{|x - \xi|^{s q + n}d(\xi,\pa\boz)^{aq}}\,d\xi \right]^{\frac{p}{q}}\,dx
\leq \| f \|_{\cf_{p,q}^{s}(\boz,d^{-aq})}^{p}.
\end{align}
Combining \eqref{eqn-5.3}, \eqref{eqn-3.4} and \eqref{eqn-3.5}, we have
\begin{equation}\label{eqn-3.6}
X\lesssim \| f \|_{\cf_{p,q}^{s}(\boz,d^{-aq})}^{p}.
\end{equation}

Next, we estimate Y. Write
\begin{align}\label{eqn-5.8}
Y&= \int_{{\boz}^\complement} \left[ \int_{\boz} \frac{| \sum_{Q\in\mathcal{W}_3} \psi_Q(x) f_{Q^{*}}-f(y)|^q}{|x - y|^{s q + n}d(y,\pa\boz)^{aq}}\,dy \right]^{\frac{p}{q}}\,dx  \nonumber \\
&\lesssim \int_{{\boz}^\complement} \left[ \int_{\boz} \frac{| \sum_{Q\in\mathcal{W}_3} \psi_Q(x)|^q | f_{Q^{*}}-f(y)|^q}{|x - y|^{s q + n}d(y,\pa\boz)^{aq}}\,dy \right]^{\frac{p}{q}}\,dx \nonumber \\
&\quad+ \int_{{\boz}^\complement} \left[ \int_{\boz} \frac{ |1- \sum_{Q\in\mathcal{W}_3} \psi_Q(x) |^q | f(y) |^q}{|x - y|^{s q + n}d(y,\pa\boz)^{aq}}\,dy \right]^{\frac{p}{q}}\,dx  \nonumber \\
&=:Y_1+Y_2.
\end{align}
For $Y_1$,
by $\text{supp}\,\psi_Q \subset \frac{11}{10}Q$ for any $Q\in\mathcal{W}_3$ and Lemma \ref{lem2.1}, we obtain
\begin{align*}
Y_1
&\lesssim \sum_{Q\in\mathcal{W}_3} \int_{\frac{11}{10}Q} | \psi_Q(x)|^p \left[ \sum_{S\in\mathcal{W}_1}\int_{S} \frac{ | f_{Q^{*}}-f(y)|^q}{|x - y|^{s q + n}d(y,\pa\boz)^{aq}}\,dy \right]^{\frac{p}{q}}\,dx \\
&\lesssim \sum_{Q\in\mathcal{W}_3} l(Q)^n \left\{ \sum_{S\in\mathcal{W}_1}\int_{S} \frac{\left[ \frac{1}{l(Q)^n}\int_{Q^{*}} |f(\xi)-f(y)| d\xi\right]^q}{D(Q,S)^{s q + n}l(S)^{aq}}\,dy \right\}^{\frac{p}{q}}.
\end{align*}
By this, Minkowski's inequality and H\"older's inequality, we obtain
\begin{align}\label{eqn-3.7}
Y_1
&\lesssim \sum_{Q\in\mathcal{W}_3} \frac{l(Q)^n}{l(Q)^{np}} \left\{ \int_{\qxing} \left[\sum_{S\in\mathcal{W}_1}\int_{S}\frac{ |f(\xi)-f(y)|^q }{|\xi - y|^{s q + n}d(y,\pa\boz)^{aq}} \,dy\right]^{\frac{1}{q}}\,d\xi \right\}^{p}  \nonumber \\
&\lesssim \sum_{Q\in\mathcal{W}_3} \frac{1}{l(Q)^{n(p-1)}} \left\{ \int_{\qxing} \left[\int_{\boz}\frac{ |f(\xi)-f(y)|^q }{|\xi - y|^{s q + n}d(y,\pa\boz)^{aq}}\,dy\right]^{\frac{p}{q}}\,d\xi \right\}l(Q)^{n(p-1)} \nonumber \\
&\leq \int_{\boz} \left[\int_{\boz}\frac{ |f(\xi)-f(y)|^q }{|\xi - y|^{s q + n}d(y,\pa\boz)^{aq}}\,dy\right]^{\frac{p}{q}}\,d\xi
\leq \| f \|_{\cf_{p,q}^{s}(\boz,d^{-aq})}^{p}.
\end{align}
For $Y_2$, by \eqref{1719} and Lemma \ref{lem2.1}, we obtain
\begin{align}\label{1538}
Y_2
&= \sum_{P\in\cw_2\setminus\cw_4} \int_{P}\left[ \int_{\boz} \frac{ |1- \sum_{Q\in\mathcal{W}_3} \psi_Q(x) |^q | f(y) |^q}{|x - y|^{s q + n}d(y,\pa\boz)^{aq}}\,dy \right]^{\frac{p}{q}}\,dx \nonumber \\
&\leq\sum_{P\in\cw_2\setminus\cw_4} \int_{P}\left[ \int_{\boz} \frac{ | f(y) |^q}{|x - y|^{s q + n}d(y,\pa\boz)^{aq}}\,dy \right]^{\frac{p}{q}} \,dx \nonumber\\
&\sim \sum_{P\in\cw_2\setminus\cw_4} \int_{P} \left[ \sum_{S\in\mathcal{W}_1} \int_{S} \frac{ | f(y) |^q}{D(P,S)^{s q + n}d(y,\pa\boz)^{aq}}\,dy \right]^{\frac{p}{q}}\,dx \nonumber\\
&\leq \sum_{P\in\cw_2\setminus\cw_4} \int_{P} \left[ \sum_{S\in\mathcal{W}_1} \int_{S} \frac{ | f(y) |^q}{d(P,\pa\boz)^{s q + n}d(y,\pa\boz)^{aq}}\,dy \right]^{\frac{p}{q}}\,dx \nonumber\\
&\sim\sum_{P\in\cw_2\setminus\cw_4}l(P)^{n-sp-\frac{np}{q}}\left[ \sum_{S\in\mathcal{W}_1} \int_{S} \frac{ | f(y) |^q}{d(y,\pa\boz)^{aq}} \,dy \right]^{\frac{p}{q}}.
\end{align}
Since $s>\frac{n}{p}-\frac{n}{q}$, we have that $s p + \frac{np}{q}-n>0$,
which, combined with \eqref{1538} and an argument similar to that used in the estimation of \eqref{eqn-3.3}, further implies that
\begin{align}\label{1539}
Y_2 \lesssim l_0^{n-sp - \frac{np}{q}}\left[ \sum_{S\in\mathcal{W}_1} \int_{S} \frac{ | f(y) |^q}{d(y,\pa\boz)^{aq}}\,dy \right]^{\frac{p}{q}}. 
\end{align}
To further estimate the right-hand side of \eqref{1539},
we consider the following two cases for $p$ and $q$.

{\bf Case (1)} $q\leq p$. In this case, by Lemma \ref{lem2.1}(i), H\"older's inequality, the assumption that $\boz$ is bounded and Hardy's inequality \eqref{eqn-Hardy} with $a+s$ replaced by $a$, we find that
\begin{align}
Y_2 &\lesssim\left[ \int_{\boz} \frac{ | f(y) |^q}{d(y,\pa\boz)^{aq}} \,dy \right]^{\frac{p}{q}}
\lesssim \int_{\boz}\frac{ | f(x) |^p}{d(x,\pa\boz)^{ap}}\,dx 
\lesssim \int_{\boz}\frac{ | f(x) |^p}{d(x,\pa\boz)^{(a+s)p}}\,dx\nonumber \\
&\lesssim \int_{\boz} \left[ \int_{\boz} \frac{|f(x) - f(y)|^q}{|x - y|^{s q + n}d(y,\pa\boz)^{aq}}\,dy \right]^{\frac{p}{q}}\,dx
\leq \| f \|_{\cf_{p,q}^{s}(\boz,d^{-aq})}^{p}.
\end{align}

{\bf Case (2)} $p<q$. In this case, we first write
\begin{align*}
Y_2&\lesssim \left[ \sum_{S\in\mathcal{W}_1} \int_{S} \frac{ | f(y) |^q}{d(y,\pa\boz)^{aq}}\,dy \right]^{\frac{p}{q}} \\
&\leq \sum_{S\in\mathcal{W}_1} \frac{1}{l(S)^n} \int_{S }\left[\int_{S} \frac{ | f(y) |^q}{d(y,\pa\boz)^{aq}}\,dy \right]^{\frac{p}{q}}\,dx \\
&\lesssim \sum_{S\in\mathcal{W}_1} \frac{1}{l(S)^n} \int_{S }\left[\int_{S} \frac{ | f(x)-f(y) |^q}{d(y,\pa\boz)^{aq}}\,dy \right]^{\frac{p}{q}}\,dx \\
&\quad+\sum_{S\in\mathcal{W}_1} \frac{1}{l(S)^n} \int_{S }\left[\int_{S} \frac{ | f(x) |^q}{d(y,\pa\boz)^{aq}}\,dy \right]^{\frac{p}{q}}\,dx \\
&=:Y_{2,1}+Y_{2,2}.
\end{align*}
To estimate $Y_{2,1}$, since $\boz$ is bounded,
it follows that for any $S\in\mathcal{W}_1$ and $x,y\in S$,
$$
|x-y|\leq \mathrm{diam\,}\boz\lesssim1
\ \ \text{and}\ \ 
|x-y|\leq\mathrm{diam\,}S\sim l(S)$$
and hence
\begin{align}\label{eqn-3.8}
Y_{2,1}&\lesssim \sum_{S\in\mathcal{W}_1}\int_{S }\left[\int_{S} \frac{ | f(x)-f(y) |^q}{|x-y|^{sq+n}d(y,\pa\boz)^{aq}}\,dy \right]^{\frac{p}{q}}\,dx \nonumber \\
&\leq \int_{\boz}\left[\int_{\boz} \frac{ | f(x)-f(y) |^q}{|x-y|^{sq+n}d(y,\pa\boz)^{aq}}\,dy \right]^{\frac{p}{q}}\,dx 
\leq \| f \|_{\cf_{p,q}^{s}(\boz,d^{-aq})}^{p}.
\end{align}
To estimate $Y_{2,2}$,
by Lemma \ref{lem2.1}(ii),
the fact that $d(x,\partial\Omega)\leq\mathrm{diam\,}\Omega\lesssim1$
for any $x\in\Omega$,
$s>\frac{n}{p}-\frac{n}{q}$ and Hardy's inequality \eqref{eqn-Hardy}, we obtain
\begin{align}\label{eqn-3.9}
Y_{2,2}&\sim \sum_{S\in\mathcal{W}_1} \frac{1}{l(S)^{p(\frac{n}{p}-\frac{n}{q})}} \int_{S }\frac{ | f(x) |^p}{d(x,\pa\boz)^{ap}}\,dx \nonumber \\
&\sim \sum_{S\in\mathcal{W}_1}\int_{S}\frac{ | f(x) |^p}{d(x,\pa\boz)^{(a+\frac{n}{p}-\frac{n}{q})p}}\,dx \nonumber \\
&\lesssim\int_{\boz}\frac{ | f(x) |^p}{d(x,\pa\boz)^{(a+s)p}}\,dx
\lesssim \| f \|_{\cf_{p,q}^{s}(\boz,d^{-aq})}^{p}.
\end{align}
By \eqref{eqn-5.8}--\eqref{eqn-3.9}, we conclude that
\begin{align}\label{eqn-3.10}
Y\lesssim \| f \|_{\cf_{p,q}^{s}(\boz,d^{-aq})}^{p}.
\end{align}

Finally, we estimate Z. Write
\begin{align}\label{eqn-6.18}
Z
&= \int_{{\boz}^\complement} \left[ \int_{{\boz}^\complement} \frac{| \sum_{P\in\mathcal{W}_3} \psi_P(x) f_{P^{*}} - \sum_{S\in\mathcal{W}_3} \psi_S(y) f_{S^{*}}|^q}{|x - y|^{s q + n}d(y,\pa\boz)^{aq}}\,dy \right]^{\frac{p}{q}}\,dx \nonumber\\
& \lesssim \sum_{Q\in\cw_2} \int_{Q} \left[ \int_{{\boz}^\complement} \frac{| \sum_{P\in\mathcal{W}_3} \psi_P(x) f_{P^{*}} - \sum_{S\in\mathcal{W}_3} \psi_S(y) f_{S^{*}}|^q}{|x - y|^{s q + n}d(y,\pa\boz)^{aq}}\,dy \right]^{\frac{p}{q}}\,dx.
\end{align}
For any $ x \in Q\in \cw_2$ and $y\in{\boz}^\complement$, we have
\begin{align}\label{eqn-3.11}
&\sum_{P\in\mathcal{W}_3} \psi_P(x) f_{P^{*}} - \sum_{S\in\mathcal{W}_3} \psi_S(y) f_{S^{*}} \nonumber \\
&\quad=\sum_{P\in\mathcal{W}_3} \psi_P(x) f_{P^{*}} \left[\sum_{S\in\mathcal{W}_3} \psi_S(y)+1-\sum_{S\in\mathcal{W}_3} \psi_S(y)\right] \nonumber \\
&\qquad- \sum_{S\in\mathcal{W}_3} \psi_S(y) f_{S^{*}} \left[\sum_{P\in\mathcal{W}_3} \psi_P(x)+1-\sum_{P\in\mathcal{W}_3} \psi_P(x)\right] \nonumber \\
&\quad=\sum_{\substack{P\in\mathcal{W}_3\\ \frac{11}{10}\overline{P}\cap\overline{Q} \neq\varnothing}} \sum_{S\in\mathcal{W}_3} \psi_P(x) \psi_S(y)(f_{\pxing}-f_{\sxing})\nonumber \\
&\qquad+ \sum_{\substack{P\in\mathcal{W}_3\\ \frac{11}{10}\overline{P}\cap\overline{Q} \neq\varnothing}}  \psi_P(x) f_{\pxing} \left[1-\sum_{S\in\mathcal{W}_3} \psi_S(y)\right] \nonumber \\
&\qquad-\sum_{S\in\mathcal{W}_3} \psi_S(y)f_{S^{*}}\left[1-\sum_{P\in\mathcal{W}_3} \psi_P(x)\right].
\end{align}
From this and an argument similar to that used in the proof of [\cite{cyz2025}, Lemma 3.10(ii)], we infer that for any $x \in Q\in \cw_2$ and $y\in B(x,\frac{l(Q)}{10})$,
\begin{align*}
\sum_{P\in\mathcal{W}_3} \psi_P(x) f_{P^{*}} - \sum_{S\in\mathcal{W}_3} \psi_S(y) f_{S^{*}}=\sum_{\substack{S\in\mathcal{W}_3\\ \frac{11}{10}\overline{S}\cap\overline{Q} \neq\varnothing}} (\psi_S(x)-\psi_S(y))f_{S^{*}}.
\end{align*}
By this, \eqref{eqn-6.18} and \eqref{eqn-3.11}, we find that
\begin{align}\label{eqn-6.19}
Z&\lesssim \sum_{Q\in\cw_4} \int_{Q} \left[ \int_{\boz^\complement\setminus B(x,\frac{l(Q)}{10})} \frac{|\sum_{P\in\mathcal{W}_3,\frac{11}{10}\overline{P}\cap\overline{Q} \neq\varnothing} \sum_{S\in\mathcal{W}_3} \psi_P(x) \psi_S(y)(f_{\pxing}-f_{\sxing})|^q
}{|x - y|^{s q + n}d(y,\pa\boz)^{aq}}\,dy \right]^{\frac{p}{q}}\,dx \nonumber\\
&\quad+\sum_{Q\in\cw_4} \int_{Q} \left[ \int_{\boz^\complement\setminus B(x,\frac{l(Q)}{10})} \frac{|\sum_{P\in\mathcal{W}_3, \frac{11}{10}\overline{P}\cap\overline{Q} \neq\varnothing}  \psi_P(x) f_{\pxing} [1-\sum_{S\in\mathcal{W}_3} \psi_S(y)]|^q}{|x - y|^{s q + n}d(y,\pa\boz)^{aq}}\,dy \right]^{\frac{p}{q}}\,dx \nonumber\\
&\quad+ \sum_{Q\in\cw_4} \int_{Q} \left[ \int_{ B(x,\frac{l(Q)}{10})} \frac{| \sum_{S\in\mathcal{W}_3, \frac{11}{10}\overline{S}\cap\overline{Q} \neq\varnothing} (\psi_S(x)-\psi_S(y))f_{S^{*}}|^q}{|x - y|^{s q + n}d(y,\pa\boz)^{aq}}\,dy \right]^{\frac{p}{q}}\,dx \nonumber\\
&\quad+ \sum_{Q\in\cw_2\setminus\cw_4} \int_{Q} \left[ \int_{B(x,\frac{l(Q)}{10})} \frac{|\sum_{S\in\mathcal{W}_3, \frac{11}{10}\overline{S}\cap\overline{Q} \neq\varnothing} (\psi_S(x)-\psi_S(y))f_{S^{*}}|^q}{|x - y|^{s q + n}d(y,\pa\boz)^{aq}}\,dy \right]^{\frac{p}{q}}\,dx \nonumber\\
&\quad+ \sum_{Q\in\cw_2\setminus\cw_4} \int_{Q} \left[ \int_{\boz^\complement\setminus B(x,\frac{l(Q)}{10})} \frac{| \sum_{P\in\mathcal{W}_3} \psi_P(x) f_{P^{*}} - \sum_{S\in\mathcal{W}_3} \psi_S(y) f_{S^{*}}|^q}{|x - y|^{s q + n}d(y,\pa\boz)^{aq}}\,dy \right]^{\frac{p}{q}}\,dx \nonumber\\
&=:Z_{1,1}+Z_{1,2}+Z_2+Z_3+Z_4.
\end{align}
To estimate $Z_{1,1}$, by the support condition and the size condition of bump functions,
Lemma \ref{lem2.1}, \eqref{eqn-7.1}, the fact that $|x-y|\gtrsim D(P,S)$ for any $x\in Q\in \mathcal{W}_4$ and $y\in S\in\mathcal{W}_3$ with $\overline{Q}\cap\frac{11}{10}\overline{P}\neq\varnothing$ and Lemma \ref{lem5.1}, we obtain
\begin{align*}
Z_{1,1}&\lesssim \sum_{Q\in\cw_4} \int_{Q} \left[ \sum_{\substack{P\in\mathcal{W}_3\\ \frac{11}{10}\overline{P}\cap\overline{Q} \neq\varnothing}} \sum_{S\in\mathcal{W}_3}\int_{\frac{11}{10}S}\frac{| \psi_P(x) \psi_S(y)(f_{\pxing}-f_{\sxing})|^q
}{|x - y|^{s q + n}d(y,\pa\boz)^{aq}}\,dy \right]^{\frac{p}{q}}\,dx \\
&\lesssim \sum_{Q\in\cw_4} \int_{Q} \left[  \sum_{\substack{P\in\mathcal{W}_3\\ \frac{11}{10}\overline{P}\cap\overline{Q} \neq\varnothing}} \sum_{S\in\mathcal{W}_3}\int_{\frac{11}{10}S}\frac{| f_{\pxing}-f_{\sxing}|^q
}{D(P,S)^{s q + n}l(S)^{aq}}\,dy \right]^{\frac{p}{q}}\,dx \\
&\lesssim \sum_{Q\in\cw_4} l(Q)^n \left[  \sum_{\substack{P\in\mathcal{W}_3\\ \frac{11}{10}\overline{P}\cap\overline{Q} \neq\varnothing}} \sum_{S\in\mathcal{W}_3} l(\sxing)^{n} \frac{| f_{\pxing}-f_{\sxing}|^q
}{D(\pxing,\sxing)^{s q + n}l(\sxing)^{aq}}  \right]^{\frac{p}{q}}  \\
&\leq \sum_{Q\in\cw_4} l(Q)^n \left\{\sum_{\substack{P\in\mathcal{W}_3\\ \frac{11}{10}\overline{P}\cap\overline{Q} \neq\varnothing}} \sum_{S\in\mathcal{W}_3} l(\sxing)^{n} \frac{\left[ \frac{1}{l(\sxing)^n l(\pxing)^n} \int_{\sxing} \int_{\pxing} |f(x)-f(y)|\,dx\,dy \right]^q
}{D(\pxing,\sxing)^{s q + n}l(\sxing)^{aq}}  \right\}^{\frac{p}{q}}. 
\end{align*}
By this, \eqref{eqn-7.1}, Lemma \ref{lem2.1}(v), H$\ddot{\text{o}}$lder's inequality, Minkowski's inequality, Tonelli's theorem, the fact that $D(P^*,S^*)\gtrsim |x-y|$ for any $x\in P^*$ and $y\in S^*$, and Lemma \ref{lem2.1}(ii), we conclude that
\begin{align}\label{eqn-3.13}
Z_{1,1}&\lesssim \sum_{Q\in\cw_4} l(Q)^n \left\{ \sum_{\substack{P\in\mathcal{W}_3\\ \frac{11}{10}\overline{P}\cap\overline{Q} \neq\varnothing}} \sum_{S\in\mathcal{W}_3} \frac{1}{ l(\pxing)^{nq}} \frac{  \int_{\sxing} [\int_{\pxing} |f(x)-f(y)|\,dx]^q\,dy 
}{D(\pxing,\sxing)^{s q + n}l(\sxing)^{aq}}  \right\}^{\frac{p}{q}} \nonumber  \\
&\sim \sum_{Q\in\cw_4} \frac{l(Q)^n}{l(Q)^{np}}\sum_{\substack{P\in\mathcal{W}_3\\ \frac{11}{10}\overline{P}\cap\overline{Q} \neq\varnothing}} \left\{ \sum_{S\in\mathcal{W}_3} \int_{\sxing} \frac{   [\int_{\pxing} |f(x)-f(y)|\,dx]^q 
}{D(\pxing,\sxing)^{s q + n}l(\sxing)^{aq}}\,dy \right\}^{\frac{p}{q}} \nonumber \\
&\leq \sum_{Q\in\cw_4} \frac{1}{l(Q)^{n(p-1)}}\sum_{\substack{P\in\mathcal{W}_3\\ \frac{11}{10}\overline{P}\cap\overline{Q} \neq\varnothing}} \left\{ \int_{\pxing} \left[\sum_{S\in\mathcal{W}_3}\int_{\sxing} \frac{ |f(x)-f(y)|^q 
}{D(\pxing,\sxing)^{s q + n}l(\sxing)^{aq}}\,dy \right]^{\frac{1}{q}} \,dx \right\}^p \nonumber \\
&\lesssim \sum_{Q\in\cw_4} \frac{1}{l(Q)^{n(p-1)}}   \sum_{\substack{P\in\mathcal{W}_3\\ \frac{11}{10}\overline{P}\cap\overline{Q} \neq\varnothing}} l(\pxing)^{\frac{np}{p'}} \int_{\pxing} \left[\int_{\boz} \frac{    |f(x)-f(y)|^q 
}{|x-y|^{s q + n}d(y,\pa\boz)^{aq}}\,dy \right]^{\frac{p}{q}}\,dx
\nonumber \\
&\sim \sum_{\substack{P\in\mathcal{W}_3}} \sum_{\substack{Q\in\mathcal{W}_4\\ \frac{11}{10}\overline{P}\cap\overline{Q} \neq\varnothing}} \int_{\pxing} \left[\int_{\boz} \frac{    |f(x)-f(y)|^q 
}{|x-y|^{s q + n}d(y,\pa\boz)^{aq}}\,dy \right]^{\frac{p}{q}}\,dx \nonumber \\
&\lesssim \sum_{\substack{P\in\mathcal{W}_3}}  \int_{\pxing} \left[\int_{\boz} \frac{    |f(x)-f(y)|^q 
}{|x-y|^{s q + n}d(y,\pa\boz)^{aq}}\,dy \right]^{\frac{p}{q}}\,dx
\leq \| f \|_{\cf_{p,q}^{s}(\boz,d^{-aq})}^{p}.
\end{align}
To estimate $Z_{1,2}$, by \eqref{1719}, Lemma~\ref{lem2.1}(iv), the facts that $0\leq\psi_P\leq1$ for any $P\in\mathcal{W}_3$ and $d(y,\pa\boz)>l_0$ for any $y\in R\in\mathcal{W}_2\setminus\mathcal{W}_4$, and H\"older's inequality, we obtain
\begin{align}\label{1222}
Z_{1,2}
&\lesssim\sum_{Q\in\cw_4} \int_{Q} \left[\sum_{\substack{P\in\mathcal{W}_3\\ \frac{11}{10}\overline{P}\cap\overline{Q} \neq\varnothing}} \int_{(\bigcup_{R\in\cw_2\setminus\cw_4}R)\setminus B(x,\frac{l(Q)}{10})
} \frac{|  \psi_P(x) f_{\pxing} |^q}{|x - y|^{s q + n}d(y,\pa\boz)^{aq}}\,dy \right]^{\frac{p}{q}}\,dx\nonumber\\
&\lesssim\sum_{Q\in\cw_4} \int_{Q} \sum_{\substack{P\in\mathcal{W}_3\\ \frac{11}{10}\overline{P}\cap\overline{Q} \neq\varnothing}}|f_{\pxing}|^p \left[ \int_{(\bigcup_{R\in\cw_2\setminus\cw_4}R)\setminus B(x,\frac{l(Q)}{10})} \frac{1}{|x - y|^{s q + n}d(y,\pa\boz)^{aq}}\,dy \right]^{\frac{p}{q}}\,dx\nonumber\\
&\leq \sum_{Q\in\cw_4} \int_{Q} \sum_{\substack{P\in\mathcal{W}_3\\ \frac{11}{10}\overline{P}\cap\overline{Q} \neq\varnothing}}\left[ \frac{1}{l(\pxing)^n}\int_{P^{*}} |f(\xi)|\,d\xi\right]^p \nonumber \\
&\quad\times\left[ \int_{(\bigcup_{R\in\cw_2\setminus\cw_4}R)\setminus B(x,\frac{l(Q)}{10})} \frac{1}{|x - y|^{s q + n}{l_0}^{aq}}\,dy \right]^{\frac{p}{q}}\,dx\nonumber\\
&\lesssim \sum_{Q\in\cw_4} \int_{Q} \sum_{\substack{P\in\mathcal{W}_3\\ \frac{11}{10}\overline{P}\cap\overline{Q} \neq\varnothing}} \frac{1}{l(\pxing)^n}\int_{P^{*}} |f(\xi)|^p\,d\xi \nonumber \\
&\quad\times\left[ \int_{(\bigcup_{R\in\cw_2\setminus\cw_4}R)\setminus B(x,\frac{l(Q)}{10})} \frac{1}{|x - y|^{s q + n}}\,dy \right]^{\frac{p}{q}}\,dx\nonumber\\
&\lesssim \sum_{Q\in\cw_4} \int_{Q} \sum_{\substack{P\in\mathcal{W}_3\\ \frac{11}{10}\overline{P}\cap\overline{Q} \neq\varnothing}} \frac{1}{l(\pxing)^n}\int_{P^{*}} |f(\xi)|^p\,d\xi\nonumber\\
&\quad\times \sum_{\substack{R\in\mathcal{W}_2\setminus\cw_4 \\ \overline{R}\cap\overline{Q} =\varnothing}}\left( \int_{R} \frac{1}{|x - y|^{s q + n}}\,dy \right)^{\frac{p}{q}}\,dx.
\end{align}
Note that, by Lemma~\ref{lem2.1}, we find that,
for any $x\in Q\in\cw_4$ and $y\in R\in\cw_2\setminus\cw_4$ with $\overline{R}\cap\overline{Q}=\varnothing$, there exists a positive constant $C$ such that $|x-y|\geq Cl(R)$, which,
combined with \eqref{1222} and polar coordinates, further implies that
\begin{align*}
Z_{1,2}&\lesssim\sum_{Q\in\cw_4} \int_{Q} \sum_{\substack{P\in\mathcal{W}_3\\ \frac{11}{10}\overline{P}\cap\overline{Q} \neq\varnothing}} \frac{1}{l(\pxing)^n}\int_{P^{*}} |f(\xi)|^p\,d\xi \sum_{\substack{R\in\mathcal{W}_2\setminus\cw_4 \\ \overline{R}\cap\overline{Q} =\varnothing}}\left[ \int_{Cl(R)}^{\infty} \rho^{-s q-1}\,d\rho  \right]^{\frac{p}{q}}\,dx \\
&\sim \sum_{Q\in\cw_4} l(Q)^n \sum_{\substack{P\in\mathcal{W}_3\\ \frac{11}{10}\overline{P}\cap\overline{Q} \neq\varnothing}} \frac{1}{l(\pxing)^n}\int_{P^{*}} |f(\xi)|^p\,d\xi \sum_{\substack{R\in\mathcal{W}_2\setminus\cw_4 \\ \overline{R}\cap\overline{Q} =\varnothing}} \frac{1}{l(R)^{sp}}.
\end{align*}
From this, an argument similar to that used in the estimation of \eqref{eqn-3.3}, and Lemma \ref{lem5.1}, it follows that
\begin{align}\label{eqn-3.14}
Z_{1,2}&\lesssim \sum_{Q\in\cw_4} l(Q)^n \sum_{\substack{P\in\mathcal{W}_3\\ \frac{11}{10}\overline{P}\cap\overline{Q} \neq\varnothing}} \frac{1}{l(\pxing)^n}\int_{P^{*}} |f(\xi)|^p\,d\xi \nonumber \\
&\lesssim \sum_{Q\in\cw_4}  \sum_{\substack{P\in\mathcal{W}_3\\ \frac{11}{10}\overline{P}\cap\overline{Q} \neq\varnothing}} \int_{P^{*}} |f(\xi)|^p \,d\xi
\lesssim \int_{\boz} |f(\xi)|^p\,d\xi
\leq \| f \|_{\cf_{p,q}^{s}(\boz,d^{-bq})}^{p}.
\end{align}
For $Z_2$, the trick is to use that $\left\{\psi_Q \right\}$ is a partition of the unity with $\psi_Q$ supported in $\frac{11}{10}Q$, that is, for any $x\in Q\in\cw_4$ and $y\in B(x,\frac{l(Q)}{10})$, $\sum_{S\in\mathcal{W}_3}\psi_S(x)=\sum_{S\in\mathcal{W}_3}\psi_S(y)=1$. Thus,
\[
Z_2=\sum_{Q\in\cw_4} \int_{Q} \left[ \int_{ B(x,\frac{l(Q)}{10})} \frac{| \sum_{S\in\mathcal{W}_3, \frac{11}{10}\overline{S}\cap\overline{Q} \neq\varnothing} (\psi_S(x)-\psi_S(y))(f_{S^{*}}-f_{\qxing})|^q}{|x - y|^{s q + n}d(y,\pa\boz)^{aq}}\,dy \right]^{\frac{p}{q}}\,dx.
\]
Using this, the mean value theorem and \eqref{eqn-2.2}, we conclude that
\begin{align*}
Z_2&\lesssim \sum_{Q\in\cw_4} \int_{Q} \left[ \sum_{\substack{S\in\mathcal{W}_3\\ \frac{11}{10}\overline{S}\cap\overline{Q} \neq\varnothing}} |f_{S^{*}}-f_{\qxing}|^q \int_{ B(x,\frac{l(Q)}{10})} \frac{1}{l(Q)^q} \frac{1}{|x - y|^{s q + n-q}d(y,\pa\boz)^{aq}}\,dy \right]^{\frac{p}{q}}\,dx \\
&\lesssim \sum_{Q\in\cw_4} \int_{Q} \left[ \sum_{\substack{S\in\mathcal{W}_3\\ \frac{11}{10}\overline{S}\cap\overline{Q} \neq\varnothing}} |f_{S^{*}}-f_{\qxing}|^q \frac{1}{l(Q)^q l(Q)^{aq}}\int_{ 0}^{\frac{l(Q)}{10}} \rho^{(1-s)q-1}\,d\rho \right]^{\frac{p}{q}}\,dx \\
&\lesssim \sum_{Q\in\cw_4} l(Q)^n \left[ \sum_{\substack{S\in\mathcal{W}_3\\ \frac{11}{10}\overline{S}\cap\overline{Q} \neq\varnothing}}\frac{ |f_{S^{*}}-f_{\qxing}|^q }{l(Q)^{sq} l(Q)^{aq}} \right]^{\frac{p}{q}}.
\end{align*}
From this and Lemma \ref{lem5.1}, we deduce that
\[
Z_2 \lesssim \sum_{Q\in\cw_4} l(Q)^n \left[ \sum_{\substack{S\in\mathcal{W}_3\\ \frac{11}{10}\overline{S}\cap\overline{Q} \neq\varnothing}}\frac{ |f_{S^{*}}-f_{\qxing}|^q }{D(\qxing,\sxing)^{sq} l(Q)^{aq}} \right]^{\frac{p}{q}}.
\]
By this and an argument similar to that used in the estimation of $Z_{1,1}$, it is easy to show that
\begin{equation}\label{eqn-3.15}
Z_2\lesssim \| f \|_{\cf_{p,q}^{s}(\boz,d^{-aq})}^{p}.
\end{equation}
For $Z_3$, by the mean value theorem, \eqref{eqn-2.2}, Lemma \ref{lem2.1}(ii), polar coordinates and H$\ddot{\text{o}}$lder's inequality, we conclude that
\begin{align*}
Z_3&\lesssim \sum_{Q\in\cw_2\setminus\cw_4} \int_{Q} \left[ \sum_{\substack{S\in\mathcal{W}_3\\ \frac{11}{10}\overline{S}\cap\overline{Q} \neq\varnothing}}\int_{B(x,\frac{l(Q)}{10})} \frac{| (\psi_S(x)-\psi_S(y))f_{S^{*}}|^q}{|x - y|^{s q + n}d(y,\pa\boz)^{aq}} \,dy \right]^{\frac{p}{q}}\,dx \\
&\lesssim \sum_{Q\in\cw_2\setminus\cw_4} \int_{Q}  \sum_{\substack{S\in\mathcal{W}_3\\ \frac{11}{10}\overline{S}\cap\overline{Q} \neq\varnothing}} |f_{S^{*}}|^p \left[\int_{B(x,\frac{l(Q)}{10})} \frac{1}{l(Q)^q}\frac{1}{|x - y|^{(s-1) q + n}l(Q)^{aq}}\,dy \right]^{\frac{p}{q}}\,dx \\
&=\sum_{Q\in\cw_2\setminus\cw_4} \int_{Q}  \sum_{\substack{S\in\mathcal{W}_3\\ \frac{11}{10}\overline{S}\cap\overline{Q} \neq\varnothing}} \frac{|f_{S^{*}}|^p}{l(Q)^p l(Q)^{ap}} \left[\int_{0}^{\frac{l(Q)}{10} }\rho^{(1-s) q -1 }\,d\rho \right]^{\frac{p}{q}}\,dx \\
&\sim \sum_{Q\in\cw_2\setminus\cw_4} \frac{l(Q)^n}{l(Q)^{ap} l(Q)^{sp}} \sum_{\substack{S\in\mathcal{W}_3\\ \frac{11}{10}\overline{S}\cap\overline{Q} \neq\varnothing}} |f_{S^{*}}|^p \\
&\lesssim \sum_{Q\in\cw_2\setminus\cw_4} \frac{1}{l(Q)^{(a+s)p}} \sum_{\substack{S\in\mathcal{W}_3\\ \frac{11}{10}\overline{S}\cap\overline{Q} \neq\varnothing}} \int_{\sxing} |f(x)|^p\,dx,
\end{align*}
which, together with an argument similar to that used in the estimation of \eqref{eqn-3.3}, implies that
\begin{align}\label{eqn-3.16}
Z_3&\lesssim {l_0}^{-(a+s)p} \sum_{S\in\mathcal{W}_3} \int_{\sxing} |f(x)|^p\,dx
\lesssim \int_{\boz} |f(x)|^p\,dx
\lesssim \| f \|_{\cf_{p,q}^{s}(\boz,d^{-aq})}^{p}.
\end{align}
For $Z_4$, by \eqref{eqn-3.11}, we write
\begin{align*}
Z_4
&\lesssim \sum_{Q\in \cw_2\setminus\cw_4} \int_{Q} \left[ \int_{\boz^\complement\setminus B(x,\frac{l(Q)}{10})} \frac{|\sum_{P\in\mathcal{W}_3, \frac{11}{10}\overline{P}\cap\overline{Q} \neq\varnothing} \sum_{S\in\mathcal{W}_3} \psi_P(x) \psi_S(y)(f_{\pxing}-f_{\sxing})|^q
}{|x - y|^{s q + n}d(y,\pa\boz)^{aq}}\,dy \right]^{\frac{p}{q}}\,dx \\
&\quad+\sum_{Q\in \cw_2\setminus\cw_4} \int_{Q} \left[ \int_{\boz^\complement\setminus B(x,\frac{l(Q)}{10})} \frac{|\sum_{P\in\mathcal{W}_3, \frac{11}{10}\overline{P}\cap\overline{Q} \neq\varnothing}  \psi_P(x) f_{\pxing} [1-\sum_{S\in\mathcal{W}_3} \psi_S(y)]|^q}{|x - y|^{s q + n}d(y,\pa\boz)^{aq}}\,dy \right]^{\frac{p}{q}}\,dx \\
&\quad+\sum_{Q\in \cw_2\setminus\cw_4} \int_{Q} \left[ \int_{\boz^\complement\setminus B(x,\frac{l(Q)}{10})} \frac{|\sum_{S\in\mathcal{W}_3}  \psi_S(y) f_{\sxing} [1-\sum_{P\in\mathcal{W}_3} \psi_P(x)]|^q}{|x - y|^{s q + n}d(y,\pa\boz)^{aq}}\,dy \right]^{\frac{p}{q}}\,dx \\
&=:Z_{4,1}+Z_{4,2}+Z_{4,3}.
\end{align*}
By an argument similar to that used in the estimations of $Z_{1,1}$ and $Z_{1,2}$, it is easy to see that
\begin{equation}\label{eqn-3.17}
Z_{4,1}\lesssim \| f \|_{\cf_{p,q}^{s}(\boz,d^{-aq})}^{p}
\end{equation}
and
\begin{equation}\label{eqn-3.18}
Z_{4,2}\lesssim \| f \|_{\cf_{p,q}^{s}(\boz,d^{-aq})}^{p}.
\end{equation}
For $Z_{4,3}$, we write
\begin{align*}
Z_{4,3}&\lesssim \sum_{Q\in \cw_2\setminus\cw_4} \int_{Q} \left[ \int_{(\bigcup_{R\in\cw_2\setminus\cw_4}R)\setminus B(x,\frac{l(Q)}{10})}\sum_{S\in\mathcal{W}_3} |\psi_S(y) f_{\sxing}|^q \frac{|  1-\sum_{P\in\mathcal{W}_3} \psi_P(x)|^q}{|x - y|^{s q + n}d(y,\pa\boz)^{aq}}\,dy \right]^{\frac{p}{q}}\,dx \\
&\quad+\sum_{Q\in \cw_2\setminus\cw_4} \int_{Q} \left[ \int_{(\bigcup_{R\in\cw_4}R)\setminus B(x,\frac{l(Q)}{10})} \sum_{S\in\mathcal{W}_3}|\psi_S(y) f_{\sxing}|^q \frac{|   1-\sum_{P\in\mathcal{W}_3} \psi_P(x)|^q}{|x - y|^{s q + n}d(y,\pa\boz)^{aq}}\,dy \right]^{\frac{p}{q}}\,dx \\
&=:Z_{4,3}^1+Z_{4,3}^2.
\end{align*}
For $Z_{4,3}^1$, by the properties of bump functions, the fact that $d(y,\pa\boz)>l_0$ for any $y\in\frac{11}{10}R$ with $R\in \mathcal{W}_2\setminus\mathcal{W}_4$, Lemma \ref{lem5.1} and H$\ddot{\text{o}}$lder's inequality, we obtain
\begin{align*}
Z_{4,3}^1&\leq \sum_{Q\in \cw_2\setminus\cw_4} \int_{Q} \left[ \sum_{S\in\mathcal{W}_3}\int_{\frac{11}{10}S} |\psi_S(y) f_{\sxing}|^q \frac{1}{|x - y|^{s q + n}d(y,\pa\boz)^{aq}}\,dy \right]^{\frac{p}{q}} \,dx \\
&\leq \sum_{Q\in \cw_2\setminus\cw_4} \int_{Q} \left[\sum_{S\in\mathcal{W}_3} |f_{\sxing}|^q \int_{\frac{11}{10}S} |\psi_S(y) |^q \frac{1}{|x - y|^{s q + n}{l_0}^{aq}}\,dy \right]^{\frac{p}{q}}\,dx \\
&\lesssim \sum_{Q\in \cw_2\setminus\cw_4} \int_{Q} \left[\sum_{S\in\mathcal{W}_3} |f_{\sxing}|^q \int_{\frac{11}{10}S} \frac{1}{D(Q,S)^{s q + n}}\,dy \right]^{\frac{p}{q}}\,dx \\
&\leq \sum_{Q\in \cw_2\setminus\cw_4} l(Q)^n \left\{\sum_{S\in\mathcal{W}_3} \frac{1}{l(\sxing)^n} \left[\int_{\sxing} |f(\xi)|^q\,d\xi\right] l(\sxing)^n \frac{1}{D(Q,S)^{s q + n}} \right\}^{\frac{p}{q}} \\
&\sim \sum_{Q\in \cw_2\setminus\cw_4} l(Q)^n \left[\sum_{S\in\mathcal{W}_3} \frac{1}{D(Q,\sxing)^{sq+n}} \int_{\sxing}|f(\xi)|^q\,d\xi\right]^{\frac{p}{q}} \\
&\lesssim \sum_{Q\in \cw_2\setminus\cw_4} \frac{1}{l(Q)^{sp+\frac{np}{q}-n}} \left[ \int_{\boz}|f(\xi)|^q \,d\xi\right]^{\frac{p}{q}}
\lesssim \|f\|_{L^q(\boz)}^{p}.
\end{align*}
If $q\leq p$, then $\|f\|_{L^q(\boz)}\lesssim \|f\|_{L^p(\boz)}$ by H\"older's inequality and the assumption that $\boz$ is bounded. If $p<q$, by an argument similar to that used in the estimation of $Y_2$, [\cite{ps2017}, Lemma 3.10] and \eqref{1544}, we find that 
\[
\|f\|_{L^q(\boz)} \lesssim \| f \|_{\cf_{p,q}^{s}(\boz)} \lesssim \| f \|_{\cf_{p,q}^{s}(\boz,d^{-aq})}.
\]
Therefore,
\begin{equation}\label{eqn-3.19}
Z_{4,3}^1 \lesssim  \| f \|_{\cf_{p,q}^{s}(\boz,d^{-aq})}^{p}.
\end{equation}
For $Z_{4,3}^2$, by the properties of bump functions, Lemma \ref{lem2.1}(ii), \eqref{eqn-7.1}, Lemma \ref{lem5.1} and H$\ddot{\text{o}}$lder's inequality, we obtain
\begin{align*}
Z_{4,3}^2 &\leq \sum_{Q\in \cw_2\setminus\cw_4} \int_{Q} \left[ \int_{(\bigcup_{R\in\cw_4}R)\setminus B(x,\frac{l(Q)}{10})} \sum_{S\in\mathcal{W}_3}|f_{\sxing}|^q \frac{|   \psi_S(y)|^q}{|x - y|^{s q + n}d(y,\pa\boz)^{aq}}\,dy \right]^{\frac{p}{q}}\,dx \\
&\leq \sum_{Q\in \cw_2\setminus\cw_4} \int_{Q} \left[\sum_{S\in\mathcal{W}_3} \int_{\frac{11}{10}S} |f_{\sxing}|^q \frac{1}{|x - y|^{s q + n}d(y,\pa\boz)^{aq}}\,dy \right]^{\frac{p}{q}} \,dx \\
&\lesssim \sum_{Q\in \cw_2\setminus\cw_4} \int_{Q} \left[\sum_{S\in\mathcal{W}_3} |f_{\sxing}|^q \int_{\frac{11}{10}S} \frac{1}{D(Q,S)^{s q + n}l(\sxing)^{aq}}\,dy \right]^{\frac{p}{q}}\,dx \\
&\lesssim \sum_{Q\in \cw_2\setminus\cw_4} l(Q)^n \left[\sum_{S\in\mathcal{W}_3} \frac{\int_{\sxing}|f(\xi)|^q \,d\xi}{l(\sxing)^{aq} l(\sxing)^n}  \frac{l(\sxing)^n}{D(Q,\sxing)^{sq+n}}\right]^{\frac{p}{q}} \\
&\lesssim \sum_{Q\in \cw_2\setminus\cw_4} \frac{1}{l(Q)^{sp+\frac{np}{q}-n}} \left[\sum_{S\in\mathcal{W}_3} \int_{\sxing} \frac{|f(\xi)|^q}{d(\xi,\pa\boz)^{aq}}\,d\xi \right]^{\frac{p}{q}} \\
&\lesssim \left[\int_{\boz} \frac{|f(\xi)|^q}{d(\xi,\pa\boz)^{aq}}\,d\xi \right]^{\frac{p}{q}}.
\end{align*}
By this, \eqref{eqn-Hardy} and an argument similar to that used in the estimation of $Y_2$, we conclude that
\begin{equation}\label{eqn-3.20}
Z_{4,3}^2\lesssim \| f \|_{\cf_{p,q}^{s}(\boz,d^{-aq})}^p.
\end{equation}
Combining \eqref{eqn-6.19}--\eqref{eqn-3.20}, we conclude that
\[
Z\lesssim \| f \|_{\cf_{p,q}^{s}(\boz,d^{-aq})}^p.
\] 
From this, \eqref{eqn-3.6} and \eqref{eqn-3.10}, we deduce that 
\[
X+Y+Z \lesssim \| f \|_{\cf_{p,q}^{s}(\boz,d^{-aq})}^{p},
\]
which completes the proof of \eqref{eqn-aim} and hence Proposition \ref{pro-main}.
\end{proof}

\begin{remark}
Let all the symbols be as in Proposition \ref{pro-main}.
Since, for any $g\in F_{p,q}^{s}(\rn,d^{-aq})$,
\begin{align*}
\|g\|_{\mathcal{F}_{p,q}^{s}(\rn,d^{-aq})}
&\leq\|g\|_{\cf_{p,q}^{s}(\rn)} +\left\{\int_{\rn} \left[ \int_{\rn} \frac{|g(x) -g(y)|^q}{|x - y|^{s q + n}d(y,\pa\boz)^{aq}}\,dy \right]^{\frac{p}{q}}\,dx \right\}^{\frac{1}{p}}\\
&=\|g \|_{F_{p,q}^{s}(\rn,d^{-aq})},
\end{align*}
it follows from Proposition \ref{pro-main}
that $\Lambda_0$ defined in \eqref{eqn-3.1}
is also a bounded linear extension operator
from $\huafwe$
to $\mathcal{F}_{p,q}^{s}(\rn,d^{-aq})$.
\end{remark}

Using Proposition \ref{pro-main}, we are now ready to prove Theorem \ref{co-main}.

\begin{proof}[Proof of Theorem \ref{co-main}]
First, we prove (i). On the one hand, by Proposition \ref{pro-main}, we find that for any $f\in \huafwe$,
\begin{align}\label{eqn-5.31}
\|f\|_{F_{p, q}^{s}(\boz,d^{-aq})}
&=\inf_{{g\in F_{p,q}^s(\rn,d^{-aq})} , g|_{\boz}=f} \|g\|_{F_{p,q}^s(\rn,d^{-aq})}\nonumber\\
&\leq \|\Lambda_0 f\|_{{F_{p,q}^s(\rn,d^{-aq})}} \lesssim \|f\|_{\huafwe}.
\end{align}
On the other hand, for any given $f\in F_{p,q}^s(\boz,d^{-aq})$ and for any $g\in F_{p,q}^s(\rn,d^{-aq})$ such that $g|_{\boz}=f$,
\begin{align*}
\| f \|_{\cf_{p,q}^{s}(\boz,d^{-aq})}
&=\| f \|_{L^p(\boz)} + \left[ \int_{\boz} \left( \int_{\boz} \frac{|f(x) - f(y)|^q}{|x - y|^{s q + n}d(y,\pa\boz)^{aq}}\,dy \right)^{\frac{p}{q}}\,dx \right]^{\frac{1}{p}} \\ 
&\leq \| f \|_{\cf_{p,q}^{s}(\boz)} + \left[ \int_{\boz} \left( \int_{\boz} \frac{|f(x) - f(y)|^q}{|x - y|^{s q + n}d(y,\pa\boz)^{aq}} \,dy \right)^{\frac{p}{q}}\,dx \right]^{\frac{1}{p}} \\ 
&\leq \| g \|_{\cf_{p,q}^{s}(\rn)} + \left[ \int_{\rn} \left( \int_{\rn} \frac{|g(x) - g(y)|^q}{|x - y|^{s q + n}d(y,\pa\boz)^{aq}} \,dy \right)^{\frac{p}{q}}\,dx \right]^{\frac{1}{p}}.
\end{align*}
Taking the infimum over all $g$ as above, we then conclude that 
\[
\| f \|_{\cf_{p,q}^{s}(\boz,d^{-aq})} \lesssim \inf_{{g\in F_{p,q}^s(\rn,d^{-aq})} , g|_{\boz}=f} \|g\|_{F_{p,q}^s(\rn,d^{-aq})} = \|f\|_{F_{p, q}^{s}(\boz,d^{-aq})},
\]
which, together with \eqref{eqn-5.31}, completes the proof of (i).

(ii) and (iii) follow from (i), the definitions
of these function spaces, Theorem \ref{thm-main}
and Remark~\ref{remark}(ii).
This completes the proof of Theorem \ref{co-main}.
\end{proof}

\noindent\textbf{Acknowledgements}\quad
The authors would like to thank both referees for their careful reading
and valuable comments,
which have definitely improved the quality of this article.

\medskip

\noindent\textbf{Funding}\quad
This work was partially supported by
the National Natural Science Foundation of China
(Grant Nos. 12571108 and 12501129).

\medskip

\noindent\textbf{Data Availability}\quad
Data sharing is not applicable to this article
as no datasets were generated or analyzed during the current study.

\section*{Declarations}

\noindent\textbf{Conflict of interest}\quad
The authors declare
that they have no conflict of interest.


\begin{thebibliography}{30}

\vspace{-0.3cm}
\bibitem{a2015}
M. S. Agranovich, Sobolev Spaces, Their Generalizations and Elliptic Problems in Smooth and Lipschitz Domains, Monographs in Mathematics, Springer, 2015.

\vspace{-0.3cm}
\bibitem{ayy2022}
R. Alvarado, D. Yang and W. Yuan, A measure characterization of embedding and extension domains for Sobolev, Triebel--Lizorkin, and Besov spaces on spaces of homogeneous type, J. Funct. Anal. 283 (2022), Paper No. 109687, 71 pp.

\vspace{-0.3cm}
\bibitem{c2000}
A. M. Caetano, Approximation by functions of compact support in Besov--Triebel--Lizorkin spaces on irregular domains, Studia. Math. 142 (2000), 47--63.

\vspace{-0.3cm}
\bibitem{cjlz2022}
J. Cao, Y. Jin, Y. Li and Q. Zhang, Hardy inequalities and interrelations of fractional Triebel--Lizorkin spaces in a bounded uniform domain, Mathematics 10 (2022), Paper No. 637, 25 pp.

\vspace{-0.3cm}
\bibitem{cyz2025}
J. Cao, D. Yang and Q. Zhang, Equivalent characterizations of fractional Sobolev spaces with partially vanishing traces on $(\vez, \delta, D)$-domains supporting $
D$-adapted fractional Hardy inequalities, Acta Math. Sin. Engl. Ser. (to appear).

%\vspace{-0.3cm}
%\bibitem{ci2017}
%A. Cianchi and N. Ioku, Canceling effects in higher-order %Hardy--Sobolev inequalities, Calc. Var. Partial Differential Equations %56 (2017), Paper No. 31, 18 pp.

\vspace{-0.3cm}
\bibitem{dilt2019}
B. Dyda, L. Ihnatsyeva, J. Lehrb\"{a}ck, H. Tuominen and A. V. V\"{a}h\"{a}kangas, Muckenhoupt $A_p$-properties of distance functions and applications to Hardy--Sobolev-type inequalities, Potential Anal. 50 (2019), 83--105.

\vspace{-0.3cm}
\bibitem{dk2021}
B. Dyda and M. Kijaczko, On density of smooth functions in weighted fractional Sobolev spaces, Nonlinear Anal. 205 (2021), Paper No. 112231, 10 pp.

\vspace{-0.3cm}
\bibitem{dk2022}
B. Dyda and M. Kijaczko, On density of compactly supported smooth functions in fractional Sobolev spaces, Ann. Mat. Pura Appl. (4) 201 (2022), 1855--1867.

\vspace{-0.3cm}
\bibitem{dk2024}
B. Dyda and M. Kijaczko, Sharp weighted fractional Hardy inequalities, Studia Math. 274 (2024), 153--171.

\vspace{-0.3cm}
\bibitem{go1979}
F. W. Gehring and B. G. Osgood, Uniform domains and the quasihyperbolic metric, J. Anal. Math. 36 (1979), 50--74 (1980).

\vspace{-0.3cm}
\bibitem{g2008}
L. Grafakos, Classical Fourier Analysis, Third edition, Graduate Texts in Math. 249, Springer, New York, 2014.

\vspace{-0.3cm}
\bibitem{g2014}
L. Grafakos, Modern Fourier Analysis, Third edition, Graduate Texts in Math. 250, Springer, New York, 2014.

\vspace{-0.3cm}
\bibitem{hit2016}
T. Heikkinen, L. Ihnatsyeva and H. Tuominen, Measure density and extension of Besov and Triebel--Lizorkin functions, J. Fourier Anal. Appl. 22 (2016), 334--382.

\vspace{-0.3cm}
\bibitem{iv2013}
L. Ihnatsyeva and A. V. V\"{a}h\"{a}kangas, Hardy inequalities in Triebel--Lizorkin spaces, Indiana Univ. Math. J. 62 (2013), 1785--1807.

\vspace{-0.3cm}
\bibitem{iv2015}
L. Ihnatsyeva and A. V. V\"{a}h\"{a}kangas, Hardy inequalities in Triebel--Lizorkin spaces II. Aikawa dimension, Ann. Mat. Pura Appl. (4) 194 (2015), 479--493.

\vspace{-0.3cm}
\bibitem{j1981}
P. W. Jones, Quasiconformal mappings and extendability of functions in Sobolev spaces, Acta Math. 147 (1981), 71--88.

\vspace{-0.3cm}
\bibitem{j1965}
K. Jord\'an, Calculus of Finite Differences, American Mathematical Soc. 1965.

\vspace{-0.3cm}
\bibitem{k2023}
M. Kijaczko, Fractional Sobolev spaces with power weights, Ann. Sc. Norm. Super. Pisa Cl. Sci. (5) 24 (2023), 1741--1764.

\vspace{-0.3cm}
\bibitem{klv2013}
A. K\"{a}enm\"{a}ki, J. Lehrb\"{a}ck and M. Vuorinen, Dimensions, Whitney covers, and tubular neighborhoods, India Univ. Math. J. 62 (2013), 1861--1889.

%\vspace{-0.3cm}
%\bibitem{o1999}
%M. P. Owen, The Hardy--Rellich inequality for polyharmonic operators, %Proc. Roy. Soc. Edinburgh Sect. A 129 (1999), 825--839.

\vspace{-0.3cm}
\bibitem{p2019}
M. Prats, Measuring Triebel--Lizorkin fractional smoothness on domains in terms of first-order differences, J. Lond. Math. Soc. (2) 100 (2019), 692--716.

\vspace{-0.3cm}
\bibitem{ps2017}
M. Prats and E. Saksman, A T(1) theorem for fractional Sobolev spaces on domains, J. Geom. Anal. 27 (2017), 2490--2538.

\vspace{-0.3cm}
\bibitem{pt2015}
M. Prats and X. Tolsa, A T(P) theorem for Sobolev spaces on domains, J. Funct. Anal. 268 (2015), 2946--2989.

\vspace{-0.3cm}
\bibitem{s1970}
E. M. Stein, Singular Integrals and Differentiability Properties of Functions, Princeton Mathematical Series, Princeton University Press, Princeton, 1970.


\vspace{-0.3cm}
\bibitem{t2001}
H. Triebel, The Structure of Functions, Monographs in Mathematics 97, Birkh\"{a}user Verlag, Basel, 2001.

\vspace{-0.3cm}
\bibitem{t2008}
H. Triebel, Function Spaces and Wavelets on Domains, EMS Tracts in Mathematics 7, European Mathematical Society, Z\"{u}rich, 2008.

\vspace{-0.3cm}
\bibitem{t2010}
H. Triebel, Theory of Function Spaces, Modern Birkh\"{a}user Classics, Birkh\"{a}user/Springer Basel AG, Basel, 2010.

\vspace{-0.3cm}
\bibitem{v1988}
J. V\"{a}is\"{a}l\"{a}, Uniform domains, Tohoku Math. J. (2) 40 (1988), 101--118.

%\vspace{-0.3cm}
%\bibitem{ysy10}
%W. Yuan, W. Sickel and D. Yang,
%Morrey and Campanato Meet Besov, Lizorkin
%and Triebel,
%Lecture Notes in Math. 2005, Springer-Verlag, Berlin, 2010.

\vspace{-0.3cm}
\bibitem{zyy2024}
Z. Zeng, D. Yang and W. Yuan, Extension and embedding of Triebel--Lizorkin-type spaces built on ball quasi-Banach spaces, J. Geom. Anal. 34 (2024), Paper No. 335, 71 pp.

\end{thebibliography}

\bigskip

\noindent Jun Cao, Haonan Lu, Guangzheng Wang and Chenfeng Zhu (Corresponding author)

\medskip

\noindent  	School of Mathematical Sciences,
Zhejiang University of Technology, Hangzhou 310023,
The People's Republic of China

\smallskip

\noindent {\it E-mails}: \texttt{caojun1860@zjut.edu.cn} (J. Cao)

\noindent\phantom{{\it E-mails:}} \texttt{haonan@zjut.edu.cn} (H. Lu)

\noindent\phantom{{\it E-mails:}} \texttt{211124090004@zjut.edu.cn} (G. Wang)

\noindent\phantom{{\it E-mails:}} \texttt{chenfengzhu@zjut.edu.cn} (C. Zhu)
\end{document} 
